\documentclass[11pt,oneside,openany]{book}

\usepackage[T1]{fontenc}
\usepackage[utf8]{inputenc}
\usepackage{amsmath,mathtools,amsthm}
\usepackage{newtxtext,newtxmath}
\usepackage[final]{microtype}
\usepackage{geometry}
\geometry{margin=1in,headheight=15pt}
\usepackage{bm}
\usepackage{booktabs,longtable,tabularx,array,multirow}
\usepackage{enumitem}
\usepackage{xcolor}
\usepackage{fancyhdr}
\usepackage{imakeidx}
\usepackage{hyperref}
\usepackage[nameinlink,capitalise,noabbrev]{cleveref}
\usepackage{bookmark}
\usepackage{url}

\definecolor{linkblue}{RGB}{18,74,130}
\definecolor{citegreen}{RGB}{24,104,74}
\hypersetup{
  colorlinks=true,
  linkcolor=linkblue,
  citecolor=citegreen,
  urlcolor=linkblue,
  pdftitle={q-Pochhammer Symbols and q-Binomial Coefficients},
  pdfsubject={A proof-oriented monograph on q-shifted factorials, Gaussian coefficients, q-series, combinatorics, and arithmetic},
  pdfkeywords={q-Pochhammer, q-binomial, Gaussian coefficient, basic hypergeometric series, q-series}
}

\makeindex[intoc]
\allowdisplaybreaks[3]
\setcounter{tocdepth}{2}
\setcounter{secnumdepth}{3}
\setlist{itemsep=2pt,topsep=5pt}

\pagestyle{fancy}
\fancyhf{}
\fancyhead[L]{\nouppercase{\leftmark}}
\fancyhead[R]{\thepage}
\renewcommand{\headrulewidth}{0.4pt}

\newtheorem{theorem}{Theorem}[chapter]
\newtheorem{lemma}[theorem]{Lemma}
\newtheorem{proposition}[theorem]{Proposition}
\newtheorem{corollary}[theorem]{Corollary}
\newtheorem{identity}[theorem]{Identity}
\newtheorem{principle}[theorem]{Principle}
\theoremstyle{definition}
\newtheorem{definition}[theorem]{Definition}
\newtheorem{example}[theorem]{Example}
\newtheorem{algorithm}[theorem]{Algorithm}
\theoremstyle{remark}
\newtheorem{remark}[theorem]{Remark}
\newtheorem{historicalnote}[theorem]{Historical note}

\newcommand{\qbinom}[2]{\genfrac{[}{]}{0pt}{}{#1}{#2}_{q}}
\newcommand{\Qbinom}[3]{\genfrac{[}{]}{0pt}{}{#1}{#2}_{#3}}
\newcommand{\qmultinom}[2]{\genfrac{[}{]}{0pt}{}{#1}{#2}_{q}}
\newcommand{\qint}[1]{[#1]_q}
\newcommand{\qfact}[1]{[#1]_q!}
\newcommand{\poch}[3]{(#1;#2)_{#3}}
\newcommand{\ph}{\varphi}
\newcommand{\eps}{\varepsilon}
\newcommand{\C}{\mathbb C}
\newcommand{\R}{\mathbb R}
\newcommand{\Q}{\mathbb Q}
\newcommand{\Z}{\mathbb Z}
\newcommand{\N}{\mathbb N}
\newcommand{\F}{\mathbb F}
\newcommand{\OO}{\mathcal O}
\newcommand{\Li}{\operatorname{Li}}
\newcommand{\Res}{\operatorname{Res}}
\newcommand{\ord}{\operatorname{ord}}
\newcommand{\inv}{\operatorname{inv}}
\newcommand{\area}{\operatorname{area}}
\newcommand{\Span}{\operatorname{span}}
\newcommand{\ch}{\operatorname{ch}}
\newcommand{\diag}{\operatorname{diag}}
\newcommand{\lcm}{\operatorname{lcm}}
\newcommand{\sgn}{\operatorname{sgn}}
\newcommand{\e}{\mathrm e}
\newcommand{\dd}{\,\mathrm d}
\newcommand{\one}{\mathbf 1}
\newcommand{\abs}[1]{\lvert #1\rvert}
\newcommand{\norm}[1]{\lVert #1\rVert}
\newcommand{\floor}[1]{\left\lfloor #1\right\rfloor}
\newcommand{\ceil}[1]{\left\lceil #1\right\rceil}
\newcommand{\set}[1]{\left\{#1\right\}}
\newcommand{\angles}[1]{\left\langle#1\right\rangle}
\newcommand{\defeq}{\mathrel{:=}}
\newcommand{\qbase}{\;q}

\DeclarePairedDelimiter{\paren}{(}{)}
\DeclarePairedDelimiter{\bracket}{[}{]}
\DeclarePairedDelimiter{\bracep}{\lbrace}{\rbrace}

\crefname{identity}{identity}{identities}
\Crefname{identity}{Identity}{Identities}
\crefname{historicalnote}{historical note}{historical notes}
\Crefname{historicalnote}{Historical note}{Historical notes}

\title{\Huge\bfseries $q$-Pochhammer Symbols and $q$-Binomial Coefficients\\[8pt]
\Large Algebra, Combinatorics, Analysis, Arithmetic, and Summation Theory}
\author{A self-contained, proof-oriented monograph}
\date{August 2026}

\begin{document}
\frontmatter
\hypersetup{pageanchor=false}
\maketitle
\hypersetup{pageanchor=true}

\chapter*{Abstract}
\addcontentsline{toc}{chapter}{Abstract}
The finite and infinite $q$-Pochhammer symbols and the Gaussian, or $q$-binomial, coefficients form a small set of notational primitives from which a remarkable amount of mathematics can be built.  They simultaneously encode weighted subsets, partitions inside rectangles, lattice paths, subspaces over finite fields, Schubert cells in Grassmannians, noncommutative binomial expansions, $q$-difference operators, basic hypergeometric summations, theta products, partition identities, and congruences at roots of unity.

This monograph develops those themes from first principles.  It proves the algebraic laws of shifted factorials; the polynomial, reciprocal, and cyclotomic structure of Gaussian coefficients; finite and infinite $q$-binomial theorems; weighted generalizations through elementary and complete symmetric functions; combinatorial and finite-geometric models; unimodality by $\mathfrak{sl}_2$ representation theory; the $q$-Gauss, $q$-Chu--Vandermonde, and $q$-Pfaff--Saalsch\"utz summations; Jackson calculus and the $q$-gamma and $q$-beta functions; Jacobi's triple product, Euler's pentagonal theorem, and the Rogers--Ramanujan identities through Bailey pairs; cyclotomic factorization, $q$-Lucas and $q$-Babbage congruences; negative upper indices and $q$-Newton interpolation; and a recent family of central $q$-binomial/Lambert-series evaluations.  Every displayed theorem used in the main development is accompanied by a proof, and the final appendices collect formulas, limiting transitions, and computational recipes.

\chapter*{Preface and source map}
\addcontentsline{toc}{chapter}{Preface and source map}
The aim is not merely to list identities, but to expose the mechanisms that repeatedly generate them.  Three such mechanisms dominate:
\begin{enumerate}
\item finite products split into residue classes, blocks, or initial and tail factors;
\item coefficient extraction turns product factorizations into convolution identities;
\item functional equations determine analytic $q$-series from one initial value.
\end{enumerate}
Once these mechanisms are visible, many formulas that initially appear unrelated become instances of the same principle.

The supplied archive was used throughout.  In particular, the weighted subset framework of Konvalina motivates \Cref{chap:weighted}; Kupershmidt's treatment of the $q$-Newton binomial informs \Cref{chap:qcalculus,chap:qnewton}; Straub's seminar on $q$-binomial congruences is reflected in \Cref{chap:cyclotomic}; the formalized $q$-series project of Lau, Lee, and Ono motivates the algebraic-versus-analytic distinctions in \Cref{chap:formalization} and supplies a modern Bailey-pair route to the Rogers--Ramanujan identities; and the central $q$-binomial paper of Chern, Dilcher, and Jiu is developed, with a complete proof, in \Cref{chap:central}.  The archive's reference chapters and encyclopedia snapshots were checked against the standard sources listed in the bibliography.  Additional material was drawn from the classical books and articles of Andrews, Askey, Bailey, Gasper, Rahman, Rota, Goldman, Reiner, Stanton, White, and others, as well as the NIST Digital Library of Mathematical Functions.

\section*{Two regimes for $q$}
Throughout the finite algebraic chapters, $q$ may be an indeterminate and identities are understood in a polynomial or rational-function ring.  In analytic chapters, unless stated otherwise,
\[
  q\in\C,\qquad 0<\abs q<1.
\]
For the Jackson integral and positivity statements we usually assume $0<q<1$.  A capital $Q$ is used for the cardinality of a finite field, so $Q$ is a prime power and normally $Q>1$.

\section*{Conventions}
Empty products equal $1$, empty sums equal $0$, and
\[
 \qbinom nk=0\qquad\text{when }k<0\text{ or }k>n
\]
for nonnegative integral $n$.  Multiple parameters are compressed as
\[
 (a_1,\ldots,a_r;q)_n=\prod_{j=1}^r(a_j;q)_n.
\]
The notation $(a;q)_\infty$ is analytic when $\abs q<1$; when formal power series are intended, that will be said explicitly.

\tableofcontents

\mainmatter

\part{The finite algebra}

\chapter{The algebra of \texorpdfstring{$q$}{q}-shifted factorials}\label{chap:pochhammer}
\index{$q$-Pochhammer symbol|see{$q$-shifted factorial}}
\index{$q$-shifted factorial}

\section{Definitions and first examples}

\begin{definition}
For $n\in\N_0$ the \emph{$q$-shifted factorial}, or \emph{$q$-Pochhammer symbol}, is
\[
 (a;q)_n=\prod_{j=0}^{n-1}(1-aq^j),\qquad (a;q)_0=1.
\]
For parameters $a_1,\dots,a_r$ we write
\[
 (a_1,\dots,a_r;q)_n=\prod_{i=1}^r(a_i;q)_n.
\]
When $0<\abs q<1$, the infinite symbol is
\[
 (a;q)_\infty=\prod_{j=0}^{\infty}(1-aq^j).
\]
\end{definition}

The word ``shifted'' refers to the geometric progression $a,aq,aq^2,\ldots$.  The first few values are
\[
 (a;q)_1=1-a,
 \qquad
 (a;q)_2=(1-a)(1-aq),
 \qquad
 (a;q)_3=(1-a)(1-aq)(1-aq^2).
\]
Important specializations include
\[
 (q;q)_n=\prod_{j=1}^n(1-q^j),
 \qquad
 (-q;q)_n=\prod_{j=1}^n(1+q^j),
 \qquad
 (a;q^r)_n=\prod_{j=0}^{n-1}(1-aq^{rj}).
\]

\section{Concatenation, dissection, and reversal}

\begin{proposition}[Concatenation law]\label{prop:concat}
For $m,n\in\N_0$,
\[
 (a;q)_{m+n}=(a;q)_m(aq^m;q)_n.
\]
Consequently, whenever the denominator is nonzero,
\[
 \frac{(a;q)_{m+n}}{(a;q)_m}=(aq^m;q)_n.
\]
\end{proposition}

\begin{proof}
Split the defining product at $j=m$:
\[
 \prod_{j=0}^{m+n-1}(1-aq^j)
 =\left(\prod_{j=0}^{m-1}(1-aq^j)\right)
  \left(\prod_{j=m}^{m+n-1}(1-aq^j)\right).
\]
In the second product put $j=m+r$, where $0\le r<n$.
\end{proof}

\begin{proposition}[Dissection into residue classes]\label{prop:dissection}
For integers $r,n\ge1$,
\[
 (a;q)_{rn}=\prod_{s=0}^{r-1}(aq^s;q^r)_n.
\]
If $\abs q<1$, then also
\[
 (a;q)_\infty=\prod_{s=0}^{r-1}(aq^s;q^r)_\infty.
\]
\end{proposition}

\begin{proof}
Every integer $j$ with $0\le j<rn$ has a unique representation $j=s+rt$ with $0\le s<r$ and $0\le t<n$.  Reordering the finite product gives the first identity.  For the infinite identity, the factors are absolutely summably close to $1$ because
\(
 \sum_{j\ge0}\abs{a q^j}<\infty,
\)
so the product may be regrouped into the $r$ residue classes.
\end{proof}

\begin{corollary}[Dissection with a remainder]\label{cor:dissection-remainder}
\index{dissection!with a remainder}
Let $r\ge1$ and $n\ge0$ be integers, let $0\le u<r$, and set $N=rn+u$.  Then
\[
 (a;q)_N
 =\prod_{s=0}^{u-1}(aq^s;q^r)_{n+1}
  \prod_{s=u}^{r-1}(aq^s;q^r)_n .
\]
For $u=0$ this is \Cref{prop:dissection}.
\end{corollary}

\begin{proof}
Division with remainder by $r$ is a bijection between the nonnegative integers $j$ and the pairs $(s,\ell)$ with $0\le s<r$ and $\ell\ge0$, given by $j=s+r\ell$.  Fix a residue $s$ and count the $\ell$ for which $j=s+r\ell$ lies in the exponent set $\{0,\ldots,N-1\}$ of the defining product.  The condition $s+r\ell<rn+u$ reads
\[
 r\ell<rn+(u-s).
\]

If $s<u$, then $0<u-s\le u<r$, so
\[
 rn<rn+(u-s)\le rn+r=r(n+1).
\]
As $r\ell$ runs through the multiples of $r$, the inequality $r\ell<rn+(u-s)$ therefore holds exactly for $\ell=0,1,\ldots,n$, which is $n+1$ values.

If $s\ge u$, then $u-s\le0$, so the inequality implies $r\ell<rn$, that is $\ell\le n-1$; conversely every such $\ell$ satisfies $s+r\ell\le s+r(n-1)=rn+(s-r)<rn\le N$.  Hence $\ell$ ranges over $0,\ldots,n-1$, which is $n$ values.

The defining product is finite and its factors commute, so it may be reordered according to the residue of $j$ modulo $r$:
\[
 (a;q)_N
 =\prod_{s=0}^{u-1}\ \prod_{\ell=0}^{n}\bigl(1-aq^{s}(q^{r})^{\ell}\bigr)\;
  \prod_{s=u}^{r-1}\ \prod_{\ell=0}^{n-1}\bigl(1-aq^{s}(q^{r})^{\ell}\bigr).
\]
Each inner product is by definition a $q^r$-shifted factorial with argument $aq^s$, namely $(aq^s;q^r)_{n+1}$ in the first family and $(aq^s;q^r)_n$ in the second.  This is the asserted identity.  When $u=0$ the first family is empty and the second has all $r$ residues, recovering \Cref{prop:dissection}.
\end{proof}

\begin{remark}\label{rem:dissection-remainder-checks}
Two consistency checks are worth recording.  The total number of linear factors on the right is
\[
 u(n+1)+(r-u)n=rn+u=N,
\]
as it must be, since $(a;q)_N$ is a product of $N$ linear factors in $a$.  At the other extreme $n=0$ the corollary reads $(a;q)_u=\prod_{s=0}^{u-1}(aq^s;q^r)_1=\prod_{s=0}^{u-1}(1-aq^s)$, again the definition.  Combining the corollary with \Cref{prop:concat} in the form $(a;q)_{rn+u}=(a;q)_{rn}(aq^{rn};q)_u$ and dividing by \Cref{prop:dissection} recovers the identity $\prod_{s=0}^{u-1}(1-aq^{s+rn})=(aq^{rn};q)_u$, so the two decompositions agree.
\end{remark}

\begin{proposition}[Base reversal]\label{prop:base-reversal}
For $n\ge0$ and $a\ne0$,
\[
 (a;q)_n=(-a)^n q^{\binom n2}(a^{-1};q^{-1})_n.
\]
Equivalently,
\[
 (a;q^{-1})_n=(-a)^nq^{-\binom n2}(a^{-1};q)_n.
\]
\end{proposition}

\begin{proof}
Factor $-aq^j$ from each term:
\[
 1-aq^j=-aq^j\bigl(1-a^{-1}q^{-j}\bigr).
\]
Multiplying for $0\le j<n$ gives the factor
\(
 (-a)^nq^{0+1+\cdots+(n-1)}=(-a)^nq^{\binom n2}.
\)
The remaining product is $(a^{-1};q^{-1})_n$.
\end{proof}

\section{Negative indices}

The concatenation law suggests the only extension to negative integers that preserves
\((a;q)_{m+n}=(a;q)_m(aq^m;q)_n\).

\begin{definition}
For $n\ge1$, define
\[
 (a;q)_{-n}=\frac{1}{(aq^{-n};q)_n},
\]
provided no factor in the denominator vanishes.
\end{definition}

\begin{proposition}[Closed form for a negative index]\label{prop:negative-index}
For $n\ge1$ and $a\ne0$,
\[
 (a;q)_{-n}
 =\frac{(-1)^n a^{-n}q^{\binom{n+1}{2}}}{(q/a;q)_n}.
\]
Moreover, for all integers $m,n$ for which both sides are defined,
\[
 (a;q)_{m+n}=(a;q)_m(aq^m;q)_n.
\]
\end{proposition}

\begin{proof}
By definition,
\[
 (aq^{-n};q)_n=\prod_{j=0}^{n-1}(1-aq^{-n+j}).
\]
Factor $-a q^{-n+j}$ from each factor.  Their product is
\[
 (-a)^nq^{-n^2+\binom n2}=(-a)^nq^{-\binom{n+1}{2}},
\]
and the remaining product is
\[
 \prod_{j=0}^{n-1}(1-a^{-1}q^{n-j})=(q/a;q)_n.
\]
Taking the reciprocal proves the formula.  The extended concatenation identity follows by cancelling the overlapping finite factors; the cases in which one of $m,n,m+n$ is negative reduce directly to the definition.
\end{proof}

\begin{corollary}[A terminating factor]\label{cor:qminusN}
If $0\le k\le N$, then
\[
 (q^{-N};q)_k=(-1)^kq^{-Nk+\binom k2}
 \frac{(q;q)_N}{(q;q)_{N-k}}.
\]
In particular, $(q^{-N};q)_k=0$ for $k>N$.
\end{corollary}

\begin{proof}
Apply base reversal to $(q^{-N};q)_k$:
\[
 (q^{-N};q)_k=(-1)^kq^{-Nk+\binom k2}(q^{N-k+1};q)_k.
\]
The last factor is $(q;q)_N/(q;q)_{N-k}$.  For $k>N$, the defining product contains the factor $1-q^0$.
\end{proof}

\section{Partial fractions and logarithmic derivatives}

\begin{theorem}[Finite partial-fraction decomposition]\label{thm:partial-fraction}
Assume $q^i\ne q^j$ for $0\le i<j\le n$.  Then
\[
 \frac{1}{(z;q)_{n+1}}
 =\sum_{j=0}^{n}
 \frac{(-1)^j q^{\binom{j+1}{2}}}
 {(q;q)_j(q;q)_{n-j}}\,
 \frac{1}{1-zq^j}.
\]
\end{theorem}

\begin{proof}
The left side has simple poles at $z=q^{-j}$, $0\le j\le n$.  The coefficient of $(1-zq^j)^{-1}$ is
\[
 A_j=\left.\frac{1-zq^j}{(z;q)_{n+1}}\right|_{z=q^{-j}}
 =\frac{1}{\prod_{r\ne j}(1-q^{r-j})}.
\]
For $r<j$, put $s=j-r$ and obtain
\[
 \prod_{r=0}^{j-1}(1-q^{r-j})
 =\prod_{s=1}^{j}(1-q^{-s})
 =(-1)^jq^{-\binom{j+1}{2}}(q;q)_j.
\]
The factors with $r>j$ multiply to $(q;q)_{n-j}$.  Thus $A_j$ is the displayed coefficient.  Subtracting the proposed right side from the left gives a rational function with no poles and tending to $0$ at infinity, hence it is identically zero.
\end{proof}

\begin{proposition}[Logarithmic derivatives]\label{prop:logder-finite}
Away from its zeros,
\[
 \frac{\partial}{\partial a}\log(a;q)_n
 =-\sum_{j=0}^{n-1}\frac{q^j}{1-aq^j}.
\]
If $a=a(x)$ is differentiable, then
\[
 \frac{\dd}{\dd x}\log(a(x);q)_n
 =-a'(x)\sum_{j=0}^{n-1}\frac{q^j}{1-a(x)q^j}.
\]
\end{proposition}

\begin{proof}
Differentiate the logarithm of the finite product term by term.
\end{proof}

\section{Infinite products and Lambert series}

\begin{theorem}[Convergence and zero set]\label{thm:poch-entire}
Let $0<\abs q<1$.  The product $(a;q)_\infty$ converges locally uniformly as a function of $a\in\C$ and therefore defines an entire function.  Its zeros are precisely
\[
 a=q^{-j},\qquad j=0,1,2,\ldots,
\]
and all are simple.
\end{theorem}

\begin{proof}
On a compact set $K\subset\C$ let $M=\max_{a\in K}\abs a$.  Then
\[
 \sum_{j=0}^{\infty}\sup_{a\in K}\abs{a q^j}
 \le \frac{M}{1-\abs q}<\infty.
\]
The standard convergence criterion for infinite products shows that
\(
 \prod_j(1-aq^j)
\)
converges uniformly on $K$; equivalently, one may apply uniform convergence to the tails after choosing an analytic branch of the logarithm near $1$.  Uniform limits of holomorphic functions are holomorphic, so the limit is entire.

If $a=q^{-j}$, one factor vanishes.  Conversely, if no factor vanishes, then the tail product is nonzero because the sum of the absolute deviations from $1$ converges, while the finite initial product is nonzero.  At $a=q^{-j}$ exactly one factor vanishes and its derivative is $-q^j\ne0$, whereas all other factors are nonzero; the zero is simple.
\end{proof}

\subsection{Complex orders}

Finite and negative integral indices extend naturally to a complex order once an infinite product is available.  Fix $0<\abs q<1$ and a branch of the logarithm at $q$, denoted by $\log q$, and write $q^\alpha=\exp(\alpha\log q)$.

\begin{definition}\label{def:complex-qshift}
For $\alpha\in\C$, define
\[
 (a;q)_\alpha
 :=\frac{(a;q)_\infty}{(aq^\alpha;q)_\infty},
 \tag{1.1}\label{eq:complex-qshift}
\]
whenever the denominator is nonzero.
\end{definition}

\begin{proposition}[Complex concatenation]\label{prop:complex-concat}
Whenever the displayed quantities are defined,
\[
 (a;q)_{\alpha+\beta}
 =(a;q)_\alpha(aq^\alpha;q)_\beta.
 \tag{1.2}\label{eq:complex-concat}
\]
In particular,
\[
 (a;q)_{\alpha+1}=(1-aq^\alpha)(a;q)_\alpha.
 \tag{1.3}\label{eq:complex-shift}
\]
For integral $\alpha$, \eqref{eq:complex-qshift} agrees with the finite and negative-index definitions already given.
\end{proposition}

\begin{proof}
By definition,
\[
 (a;q)_\alpha(aq^\alpha;q)_\beta
 =\frac{(a;q)_\infty}{(aq^\alpha;q)_\infty}
  \frac{(aq^\alpha;q)_\infty}{(aq^{\alpha+\beta};q)_\infty}
 =\frac{(a;q)_\infty}{(aq^{\alpha+\beta};q)_\infty}.
\]
This is \eqref{eq:complex-concat}.  Taking $\beta=1$ gives \eqref{eq:complex-shift}.  If $\alpha=n\ge0$, the tail factorization
\(
 (a;q)_\infty=(a;q)_n(aq^n;q)_\infty
\)
recovers the finite product.  If $\alpha=-n<0$, the same relation, applied after shifting by $-n$, gives
\(
 (a;q)_{-n}=1/(aq^{-n};q)_n.
\)
\end{proof}

\begin{proposition}[Derivative with respect to the order]\label{prop:order-derivative}
On every domain avoiding the poles of \eqref{eq:complex-qshift},
\[
 \frac{\partial}{\partial\alpha}\log(a;q)_\alpha
 =(\log q)\sum_{j=0}^{\infty}
 \frac{aq^{\alpha+j}}{1-aq^{\alpha+j}}.
 \tag{1.4}\label{eq:order-derivative}
\]
The series is locally uniformly convergent.
\end{proposition}

\begin{proof}
Only the denominator in \eqref{eq:complex-qshift} depends on $\alpha$.  Differentiate its locally uniformly convergent logarithm term by term:
\[
 -\frac{\partial}{\partial\alpha}
 \sum_{j\ge0}\log(1-aq^{\alpha+j})
 =(\log q)\sum_{j\ge0}
 \frac{aq^{\alpha+j}}{1-aq^{\alpha+j}}.
\]
On a compact set avoiding the poles, the denominators of the tail are bounded away from zero and the numerators form a geometric majorant, proving local uniform convergence.
\end{proof}

\begin{theorem}[Lambert-series expansion]\label{thm:lambert-log}
If $\abs q<1$ and $\abs a<1$, then
\[
 \log(a;q)_\infty
 =-\sum_{m=1}^{\infty}\frac{a^m}{m(1-q^m)}.
\]
Consequently,
\[
 \frac{\partial}{\partial a}\log(a;q)_\infty
 =-\sum_{m=1}^{\infty}\frac{a^{m-1}}{1-q^m}
 =-\sum_{j=0}^{\infty}\frac{q^j}{1-aq^j}.
\]
All series converge locally uniformly in $\abs a<1$.
\end{theorem}

\begin{proof}
For $\abs a<1$, use $\log(1-u)=-\sum_{m\ge1}u^m/m$ and absolute convergence:
\[
 \sum_{j\ge0}\sum_{m\ge1}\frac{\abs a^m\abs q^{jm}}{m}
 =\sum_{m\ge1}\frac{\abs a^m}{m(1-\abs q^m)}<\infty.
\]
Thus the order of summation may be exchanged:
\[
 \sum_{j\ge0}\log(1-aq^j)
 =-\sum_{m\ge1}\frac{a^m}{m}\sum_{j\ge0}q^{jm}.
\]
The geometric sum gives the first formula.  Local uniform convergence permits termwise differentiation.  Expanding $(1-aq^j)^{-1}$ geometrically gives the second form of the derivative.
\end{proof}

\begin{corollary}[Divisor sums]\label{cor:divisor-sum}
As a formal power series,
\[
 -q\frac{\dd}{\dd q}\log(q;q)_\infty
 =\sum_{n=1}^{\infty}\sigma_1(n)q^n,
\]
where $\sigma_1(n)=\sum_{d\mid n}d$.
\end{corollary}

\begin{proof}
Logarithmic differentiation gives
\[
 -q\frac{\dd}{\dd q}\log(q;q)_\infty
 =\sum_{j\ge1}\frac{j q^j}{1-q^j}
 =\sum_{j\ge1}\sum_{r\ge1}j q^{jr}.
\]
The coefficient of $q^n$ is the sum of all $j$ dividing $n$.
\end{proof}

\section{The dilogarithmic limit}

The singular regime $q\to1^-$ turns an infinite product into an exponential of order $(1-q)^{-1}$.  The dilogarithm is the natural leading term.

\begin{theorem}[Asymptotic expansion near $q=1$]\label{thm:poch-dilog}
Fix $0<\rho<1$ and an integer $N\ge1$.  Uniformly for $\abs a\le\rho$, as $t\to0^+$,
\[
\begin{aligned}
 \log(a;\e^{-t})_\infty
 ={}&-\frac{\Li_2(a)}{t}+\frac12\log(1-a)\\
 &-\sum_{r=1}^{N-1}\frac{B_{2r}}{(2r)!}
 t^{2r-1}\Li_{2-2r}(a)+O_{\rho,N}(t^{2N-1}),
\end{aligned}
\]
where $B_{2r}$ are Bernoulli numbers and
\(
 \Li_s(a)=\sum_{m\ge1}a^m/m^s
\)
for $\abs a<1$.
\end{theorem}

\begin{proof}
By \Cref{thm:lambert-log},
\[
 \log(a;\e^{-t})_\infty
 =-\sum_{m\ge1}\frac{a^m}{m(1-\e^{-mt})}.
\]
The Bernoulli expansion at the origin is
\[
 \frac1{1-\e^{-x}}
 =\frac1x+\frac12+\sum_{r=1}^{N-1}
 \frac{B_{2r}}{(2r)!}x^{2r-1}+R_N(x),
\]
with $R_N(x)=O_N(x^{2N-1})$ for $0\le x\le1$.  Substitution of the displayed polynomial terms gives respectively
\[
 -\frac1t\sum_{m\ge1}\frac{a^m}{m^2},
 \qquad
 -\frac12\sum_{m\ge1}\frac{a^m}{m}=\frac12\log(1-a),
\]
and the polylogarithmic terms in the statement.

It remains to bound the total remainder.  Split the $m$-sum at $m\le t^{-1}$.  In this range,
\[
 \sum_{m\le t^{-1}}\frac{\abs a^m}{m}\abs{R_N(mt)}
 \le C_Nt^{2N-1}\sum_{m\ge1}\rho^m m^{2N-2}
 =O_{\rho,N}(t^{2N-1}).
\]
For $m>t^{-1}$, subtracting finitely many asymptotic terms leaves a quantity bounded by a polynomial in $mt$ plus $(1-\e^{-mt})^{-1}$; the factor $\rho^m$ makes the tail exponentially small in $1/t$, hence smaller than every power of $t$.  This proves the uniform remainder estimate.
\end{proof}

\section{Classical shifted factorials as a limit}

\begin{proposition}[Finite classical limit]\label{prop:classical-poch-limit}
For fixed $x\in\C$ and $n\in\N_0$,
\[
 \lim_{q\to1}\frac{(q^x;q)_n}{(1-q)^n}
 =x(x+1)\cdots(x+n-1)=(x)_n,
\]
where the limit is taken through values for which a consistent branch of $q^x$ is chosen.
\end{proposition}

\begin{proof}
Write
\[
 \frac{(q^x;q)_n}{(1-q)^n}
 =\prod_{j=0}^{n-1}\frac{1-q^{x+j}}{1-q}.
\]
Each factor tends to $x+j$, since $q^u=1+u(q-1)+o(q-1)$ for fixed $u$.
\end{proof}

\chapter{\texorpdfstring{$q$}{q}-integers and Gaussian coefficients}\label{chap:gaussian}
\index{Gaussian coefficient}
\index{$q$-binomial coefficient|see{Gaussian coefficient}}

\section{Definitions}

\begin{definition}
For $n\in\N_0$ define the $q$-integer and $q$-factorial by
\[
 [n]_q=\frac{1-q^n}{1-q}=1+q+\cdots+q^{n-1},
 \qquad
 [n]_q!=\prod_{j=1}^{n}[j]_q.
\]
For $0\le k\le n$, the \emph{Gaussian coefficient} is
\[
 \qbinom nk
 =\frac{(q;q)_n}{(q;q)_k(q;q)_{n-k}}
 =\frac{[n]_q!}{[k]_q![n-k]_q!}.
\]
Outside $0\le k\le n$ it is defined to be $0$.
\end{definition}

The notation is initially a rational expression, but the recurrences below show that it is a polynomial in $q$ with nonnegative integral coefficients.

\begin{proposition}[Product forms]\label{prop:qbinom-products}
For $0\le k\le n$,
\[
 \qbinom nk
 =\frac{(q^{n-k+1};q)_k}{(q;q)_k}
 =\prod_{j=1}^{k}\frac{1-q^{n-k+j}}{1-q^j}.
\]
\end{proposition}

\begin{proof}
Cancel $(q;q)_{n-k}$ from $(q;q)_n$ using \Cref{prop:concat}.
\end{proof}

\section{The two Pascal recurrences}

\begin{theorem}[$q$-Pascal recurrences]\label{thm:qpascal}
For $1\le k\le n-1$,
\begin{align}
 \qbinom nk&=\qbinom{n-1}{k-1}+q^k\qbinom{n-1}{k},\label{eq:qpascal1}\\
 \qbinom nk&=q^{n-k}\qbinom{n-1}{k-1}+\qbinom{n-1}{k}.\label{eq:qpascal2}
\end{align}
Both identities remain valid at the boundary under the zero convention.
\end{theorem}

\begin{proof}
Put the two terms on the right of \eqref{eq:qpascal1} over the common denominator $(q;q)_k(q;q)_{n-k}$.  Their numerator is
\[
 (q;q)_{n-1}\bigl((1-q^k)+q^k(1-q^{n-k})\bigr)
 =(q;q)_{n-1}(1-q^n)=(q;q)_n.
\]
The proof of \eqref{eq:qpascal2} is the same, using
\(
 q^{n-k}(1-q^k)+(1-q^{n-k})=1-q^n.
\)
\end{proof}

\begin{corollary}[Polynomiality and positivity]\label{cor:positivity}
For all $0\le k\le n$, the Gaussian coefficient $\qbinom nk$ belongs to $\N_0[q]$.
\end{corollary}

\begin{proof}
Induct on $n$.  The boundary values are $1$.  Either recurrence in \Cref{thm:qpascal} expresses an interior coefficient as a sum of two polynomials with nonnegative integral coefficients.
\end{proof}

\section{Symmetry, reciprocity, and degree}

\begin{theorem}[Elementary structure]\label{thm:qbinom-structure}
For $0\le k\le n$,
\begin{enumerate}[label=\textup{(\roman*)}]
\item $\qbinom nk=\qbinom n{n-k}$;
\item $\displaystyle \Qbinom nk{q^{-1}}=q^{-k(n-k)}\qbinom nk$;
\item $\qbinom nk$ has degree $k(n-k)$, constant term $1$, and leading coefficient $1$;
\item its coefficient sequence is palindromic.
\end{enumerate}
\end{theorem}

\begin{proof}
Symmetry is immediate from the factorial definition.  For reciprocity, use
\[
 (q^{-1};q^{-1})_m=(-1)^mq^{-\binom{m+1}{2}}(q;q)_m
\]
for $m=n,k,n-k$.  The signs cancel and the exponent is
\[
 -\binom{n+1}{2}+\binom{k+1}{2}+\binom{n-k+1}{2}=-k(n-k).
\]
Since $\qbinom nk$ is a polynomial with constant term $1$, reciprocity says
\[
 \qbinom nk=q^{k(n-k)}\Qbinom nk{q^{-1}}.
\]
Hence its degree is at most $k(n-k)$ and its leading coefficient equals its constant coefficient.  The product form shows that the order at infinity is exactly $k(n-k)$, so the degree is exact.  The same reciprocal identity equates coefficients at complementary degrees and proves palindromicity.
\end{proof}

\begin{corollary}[Classical limit]\label{cor:qbinom-classical}
For fixed $n,k$,
\[
 \lim_{q\to1}\qbinom nk=\binom nk.
\]
\end{corollary}

\begin{proof}
Use the $q$-factorial form and $[j]_q\to j$.
\end{proof}

\section{Multinomial coefficients and flags}

\begin{definition}
If $n_1,\dots,n_r\ge0$ and $n_1+\cdots+n_r=n$, define
\[
 \genfrac{[}{]}{0pt}{}{n}{n_1,\ldots,n_r}_q
 =\frac{(q;q)_n}{(q;q)_{n_1}\cdots(q;q)_{n_r}}.
\]
\end{definition}

\begin{proposition}[Factorization into binomials]\label{prop:qmult-factor}
Let $N_j=n_j+n_{j+1}+\cdots+n_r$.  Then
\[
 \genfrac{[}{]}{0pt}{}{n}{n_1,\ldots,n_r}_q
 =\prod_{j=1}^{r-1}\genfrac{[}{]}{0pt}{}{N_j}{n_j}_q.
\]
In particular, every $q$-multinomial coefficient is a polynomial with nonnegative integral coefficients.
\end{proposition}

\begin{proof}
Write out the factorial ratios.  Consecutive numerator and denominator factors telescope:
\[
 \prod_{j=1}^{r-1}
 \frac{(q;q)_{N_j}}{(q;q)_{n_j}(q;q)_{N_{j+1}}}
 =\frac{(q;q)_n}{\prod_{j=1}^{r}(q;q)_{n_j}}.
\]
Positivity follows from \Cref{cor:positivity}.
\end{proof}

\section{Adjacent ratios and efficient recurrences}

\begin{proposition}\label{prop:adjacent-ratio}
For $0\le k<n$,
\[
 \frac{\qbinom n{k+1}}{\qbinom nk}
 =\frac{1-q^{n-k}}{1-q^{k+1}}.
\]
For $0\le k\le n-1$,
\[
 \frac{\qbinom n k}{\qbinom{n-1}k}
 =\frac{1-q^n}{1-q^{n-k}}.
\]
\end{proposition}

\begin{proof}
Cancel common factors in the product or factorial definitions.
\end{proof}

These ratios are useful numerically, but near roots of unity they hide cancellations.  The polynomial recurrences of \Cref{thm:qpascal} are then preferable.

\section{Coefficient moments}

Write
\[
 \qbinom nk=\sum_{j=0}^{D}c_jq^j,
 \qquad D=k(n-k),
\]
and choose an exponent $X$ with probability
\(
 \Pr(X=j)=c_j/\binom nk.
\)
The combinatorial meaning of $X$ will be established in \Cref{chap:combinatorics}; its moments can already be found algebraically.

\begin{theorem}[Mean and variance]\label{thm:qbinom-moments}
For the preceding distribution,
\[
 \mathbb E X=\frac{k(n-k)}2,
 \qquad
 \operatorname{Var}(X)=\frac{k(n-k)(n+1)}{12}.
\]
Equivalently,
\[
 \left.\frac{\dd}{\dd q}\qbinom nk\right|_{q=1}
 =\frac{k(n-k)}2\binom nk.
\]
\end{theorem}

\begin{proof}
The palindromic identity $c_j=c_{D-j}$ immediately gives $\mathbb EX=D/2$ and the derivative formula.

For the variance, use the moment-generating function
\[
 M(t)=\frac{\Qbinom nk{\e^t}}{\binom nk}
 =\prod_{i=1}^{k}
 \frac{\e^{(n-k+i)t}-1}{(n-k+i)t}
 \frac{it}{\e^{it}-1}.
\]
The local expansion
\[
 \log\frac{\e^{u}-1}{u}=\frac u2+\frac{u^2}{24}+O(u^3)
\]
gives
\[
 \log M(t)
 =\frac t2\sum_{i=1}^k(n-k)
 +\frac{t^2}{24}\sum_{i=1}^k\bigl((n-k+i)^2-i^2\bigr)+O(t^3).
\]
The first sum is $k(n-k)$.  The second equals
\[
 k(n-k)^2+2(n-k)\frac{k(k+1)}2
 =k(n-k)(n+1).
\]
Since the first and second derivatives of $\log M$ at $0$ are the mean and variance, respectively, the result follows.
\end{proof}

\chapter{Finite \texorpdfstring{$q$}{q}-binomial theorems and inversion}\label{chap:finite-qbinomial}

\section{Euler's finite product expansions}

\begin{theorem}[Finite $q$-binomial theorem]\label{thm:finite-qbinomial}
For $n\ge0$,
\begin{align}
 (z;q)_n
 &=\sum_{k=0}^{n}(-1)^kq^{\binom k2}\qbinom nk z^k,\label{eq:finite-qbinomial-minus}\\
 (-z;q)_n
 &=\sum_{k=0}^{n}q^{\binom k2}\qbinom nk z^k.\label{eq:finite-qbinomial-plus}
\end{align}
\end{theorem}

\begin{proof}
It is enough to prove the first identity; replace $z$ by $-z$ for the second.  The case $n=0$ is clear.  Suppose the formula holds for $n$.  Then
\[
 (z;q)_{n+1}=(1-zq^n)(z;q)_n.
\]
The coefficient of $z^k$ on the right is
\[
 (-1)^kq^{\binom k2}
 \left(\qbinom nk+q^{n-k+1}\qbinom n{k-1}\right).
\]
By the second $q$-Pascal recurrence, the parenthesis is $\qbinom{n+1}k$.  This completes the induction.
\end{proof}

\begin{corollary}[Alternating sums]\label{cor:alternating-qbinom}
For $n>0$,
\[
 \sum_{k=0}^{n}(-1)^kq^{\binom k2}\qbinom nk=0.
\]
More generally,
\[
 \sum_{k=0}^{n}(-1)^kq^{\binom k2+mk}\qbinom nk=(q^m;q)_n.
\]
\end{corollary}

\begin{proof}
Set $z=1$ and $z=q^m$, respectively, in \eqref{eq:finite-qbinomial-minus}.
\end{proof}

\begin{corollary}[Weighted subset sum]\label{cor:weighted-subsets}
For $0\le k\le n$,
\[
 \sum_{0\le i_1<\cdots<i_k\le n-1}q^{i_1+\cdots+i_k}
 =q^{\binom k2}\qbinom nk.
\]
Equivalently,
\[
 \sum_{1\le j_1<\cdots<j_k\le n}q^{j_1+\cdots+j_k}
 =q^{\binom{k+1}{2}}\qbinom nk.
\]
\end{corollary}

\begin{proof}
Expand $\prod_{i=0}^{n-1}(1+zq^i)$ by choosing either $1$ or $zq^i$ from each factor, and compare the coefficient of $z^k$ with \eqref{eq:finite-qbinomial-plus}.
\end{proof}

\section{Reciprocal finite products}

\begin{theorem}[Reciprocal finite $q$-binomial theorem]\label{thm:reciprocal-finite}
For $n\ge0$ and $\abs z<1$,
\[
 \frac{1}{(z;q)_{n+1}}
 =\sum_{k=0}^{\infty}\qbinom{n+k}{k}z^k.
\]
The identity also holds formally in $\Z[q][[z]]$.
\end{theorem}

\begin{proof}
Let
\(
 H_n(z)=\sum_{k\ge0}\qbinom{n+k}{k}z^k.
\)
The second $q$-Pascal recurrence gives
\[
 \qbinom{n+k}{k}-q^n\qbinom{n+k-1}{k-1}
 =\qbinom{n+k-1}{k}.
\]
Therefore
\[
 (1-zq^n)H_n(z)=H_{n-1}(z),
\]
where $H_0(z)=\sum_{k\ge0}z^k=(1-z)^{-1}$.  Iterating yields
\[
 H_n(z)=\frac1{(1-z)(1-zq)\cdots(1-zq^n)}.
\]
All operations are valid formally; analytically they are justified for $\abs z<1$.
\end{proof}

\section{\texorpdfstring{$q$}{q}-Vandermonde convolution}

\begin{theorem}[$q$-Vandermonde convolution]\label{thm:q-vandermonde}
For $m,n,r\in\N_0$,
\begin{align}
 \qbinom{m+n}{r}
 &=\sum_{j\in\Z}q^{(m-j)(r-j)}
 \qbinom mj\qbinom n{r-j},\label{eq:qvand1}\\
 &=\sum_{j\in\Z}q^{j(n-r+j)}
 \qbinom mj\qbinom n{r-j}.\label{eq:qvand2}
\end{align}
Only finitely many terms are nonzero.
\end{theorem}

\begin{proof}
Factor
\[
 (-z;q)_{m+n}=(-z;q)_m(-zq^m;q)_n
\]
and use \eqref{eq:finite-qbinomial-plus}.  The coefficient of $z^r$ on the left is
\(
 q^{\binom r2}\qbinom{m+n}{r}.
\)
The term with $j$ chosen from the first factor and $r-j$ from the second has exponent
\[
 \binom j2+m(r-j)+\binom{r-j}{2}.
\]
Subtracting $\binom r2$ gives $(m-j)(r-j)$, proving \eqref{eq:qvand1}.  Splitting the product in the reverse order gives \eqref{eq:qvand2}.
\end{proof}

\begin{corollary}[Central form]\label{cor:central-vandermonde}
For $m,n\ge0$,
\[
 \qbinom{m+n}{m}
 =\sum_j q^{j^2}\qbinom mj\qbinom n j.
\]
\end{corollary}

\begin{proof}
Set $r=m$ in \eqref{eq:qvand1} and put $\ell=m-j$.  Then the exponent becomes $\ell^2$, while symmetry gives $\qbinom m{m-\ell}=\qbinom m\ell$ and the second factor is $\qbinom n\ell$.
\end{proof}

\section{\texorpdfstring{$q$}{q}-binomial inversion}

\begin{theorem}[$q$-binomial inversion]\label{thm:qbinomial-inversion}
For sequences $(a_n)$ and $(b_n)$ in a commutative ring containing $q$, the relations
\[
 b_n=\sum_{k=0}^{n}\qbinom nk a_k
\]
are equivalent to
\[
 a_n=\sum_{k=0}^{n}(-1)^{n-k}q^{\binom{n-k}{2}}\qbinom nk b_k.
\]
\end{theorem}

\begin{proof}
Substitute the first relation into the proposed inverse.  The coefficient of $a_j$ becomes
\[
 \sum_{k=j}^{n}(-1)^{n-k}q^{\binom{n-k}{2}}
 \qbinom nk\qbinom kj.
\]
The factorial definition gives
\[
 \qbinom nk\qbinom kj=\qbinom nj\qbinom{n-j}{k-j}.
\]
Put $r=n-k$.  The coefficient is
\[
 \qbinom nj
 \sum_{r=0}^{n-j}(-1)^rq^{\binom r2}\qbinom{n-j}{r}.
\]
By \Cref{cor:alternating-qbinom}, this equals $0$ when $n>j$ and $1$ when $n=j$.  Thus the composition is the identity.  The same argument in the opposite order proves equivalence.
\end{proof}

\begin{corollary}[Matrix orthogonality]\label{cor:qbinomial-orthogonality}
For $0\le j\le n$,
\[
 \sum_{k=j}^{n}(-1)^{n-k}q^{\binom{n-k}{2}}
 \qbinom nk\qbinom kj=\delta_{n,j}.
\]
\end{corollary}

\begin{proof}
This is the kernel computation in the preceding proof.
\end{proof}

\section{A \texorpdfstring{$q$}{q}-Cauchy product identity}

\begin{theorem}[$q$-Cauchy identity]\label{thm:q-cauchy}
For $u,v$ in a commutative ring,
\[
 (uv;q)_n
 =\sum_{k=0}^{n}\qbinom nk (u;q)_k v^k(v;q)_{n-k}.
\]
\end{theorem}

\begin{proof}
Induct on $n$.  The assertion is clear at $n=0$.  Let the right side be $C_n(u,v)$.  Use
\[
 \qbinom{n+1}{k}=q^{n+1-k}\qbinom n{k-1}+\qbinom nk
\]
and split $C_{n+1}$ accordingly.  In the first part shift $k\mapsto k+1$, use
\((u;q)_{k+1}=(u;q)_k(1-uq^k)\), and combine with the second part after using
\((v;q)_{n+1-k}=(v;q)_{n-k}(1-vq^{n-k})\).
The bracket that remains is
\[
 (1-vq^{n-k})+vq^{n-k}(1-uq^k)=1-uvq^n,
\]
independent of $k$.  Therefore
\[
 C_{n+1}(u,v)=(1-uvq^n)C_n(u,v)=(uv;q)_{n+1}.
\]
\end{proof}

\begin{remark}
Taking $u=0$ recovers the tautology $(0;q)_k=1$ inside a weighted expansion; taking special limits of $u$ or $v$ produces several finite convolution formulas used later in the proof of the $q$-Pfaff--Saalsch\"utz summation.
\end{remark}


\chapter{Weighted binomial coefficients and symmetric functions}\label{chap:weighted}
\index{weighted binomial coefficient}
\index{symmetric function}

The ordinary and Gaussian binomial coefficients are instances of a broader construction in which each available object carries a weight.  This point of view, developed systematically by Konvalina, places subset selection, multiset selection, Stirling numbers, and Gaussian coefficients in one framework.

\section{Elementary and complete weighted coefficients}

\begin{definition}
Let $\mathbf w=(w_1,\dots,w_n)$ be elements of a commutative ring.  Define
\begin{align*}
 C_k^n(\mathbf w)
 &=\sum_{1\le i_1<\cdots<i_k\le n}w_{i_1}\cdots w_{i_k},\\
 S_k^n(\mathbf w)
 &=\sum_{1\le i_1\le\cdots\le i_k\le n}w_{i_1}\cdots w_{i_k}.
\end{align*}
We set $C_0^n=S_0^n=1$.  Thus $C_k^n$ is the elementary symmetric function $e_k(w_1,\dots,w_n)$ and $S_k^n$ is the complete homogeneous symmetric function $h_k(w_1,\dots,w_n)$.
\end{definition}

\begin{theorem}[Weighted Pascal recurrences]\label{thm:weighted-pascal}
For $n,k\ge1$,
\begin{align}
 C_k^n(\mathbf w)&=C_k^{n-1}(\mathbf w)+w_nC_{k-1}^{n-1}(\mathbf w),\label{eq:weighted-C}\\
 S_k^n(\mathbf w)&=S_k^{n-1}(\mathbf w)+w_nS_{k-1}^{n}(\mathbf w).\label{eq:weighted-S}
\end{align}
\end{theorem}

\begin{proof}
A $k$-element subset of $\{1,\dots,n\}$ either omits $n$, contributing $C_k^{n-1}$, or contains $n$, in which case removal of $n$ leaves a $(k-1)$-element subset and the weight acquires a factor $w_n$.  This proves \eqref{eq:weighted-C}.

For multisets, separate those with no occurrence of $n$ from those with at least one occurrence.  Removing one occurrence of $n$ from the latter leaves an arbitrary multiset of size $k-1$ drawn from all $n$ types.  This proves \eqref{eq:weighted-S}.
\end{proof}

\begin{theorem}[Generating products]\label{thm:weighted-generating}
One has
\begin{align}
 \prod_{i=1}^{n}(1+w_i z)
 &=\sum_{k=0}^{n}C_k^n(\mathbf w)z^k,\label{eq:weighted-e-gen}\\
 \prod_{i=1}^{n}\frac1{1-w_i z}
 &=\sum_{k=0}^{\infty}S_k^n(\mathbf w)z^k.\label{eq:weighted-h-gen}
\end{align}
The second identity is formal in $z$, or analytic when $\abs{w_i z}<1$ for all $i$.
\end{theorem}

\begin{proof}
In the first product, choosing $w_i z$ from exactly the factors indexed by a subset $I$ produces $z^{\abs I}\prod_{i\in I}w_i$.  Summing over subsets gives \eqref{eq:weighted-e-gen}.  For the second, expand each factor geometrically:
\[
 \frac1{1-w_i z}=\sum_{r_i\ge0}w_i^{r_i}z^{r_i}.
\]
The coefficient of $z^k$ sums over weak compositions $r_1+\cdots+r_n=k$, which are equivalent to multisets of size $k$.
\end{proof}

\begin{corollary}[Elementary--complete orthogonality]\label{cor:e-h-orthogonality}
For $m\ge0$,
\[
 \sum_{j=0}^{m}(-1)^jC_j^n(\mathbf w)S_{m-j}^n(\mathbf w)=\delta_{m,0}.
\]
\end{corollary}

\begin{proof}
Multiply
\(
 \prod_i(1-w_i z)
\)
by its reciprocal and compare the coefficient of $z^m$, using \Cref{thm:weighted-generating} with $z$ replaced by $-z$ in the elementary generating function.
\end{proof}

\section{Classical, Stirling, and Gaussian specializations}

\begin{theorem}[A specialization table]\label{thm:weighted-specializations}
The following identities hold.
\begin{enumerate}[label=\textup{(\roman*)}]
\item If $w_i=1$, then
\[
 C_k^n=\binom nk,
 \qquad
 S_k^n=\binom{n+k-1}{k}.
\]
\item If $w_i=q^{i-1}$, then
\[
 C_k^n=q^{\binom k2}\qbinom nk,
 \qquad
 S_k^n=\qbinom{n+k-1}{k}.
\]
\item For the unsigned Stirling numbers of the first kind $c(n,k)$ and Stirling numbers of the second kind $\left\{\begin{smallmatrix}n\\k\end{smallmatrix}\right\}$,
\[
 c(n+1,k+1)=C_{n-k}^{n}(1,2,\dots,n),
\]
\[
 \left\{\begin{matrix}n+k\\n\end{matrix}\right\}
 =S_k^n(1,2,\dots,n).
\]
\end{enumerate}
\end{theorem}

\begin{proof}
For $w_i=1$, $C_k^n$ counts $k$-subsets.  The multiset count is the stars-and-bars number $\binom{n+k-1}{k}$.

For $w_i=q^{i-1}$, the elementary identity is \Cref{cor:weighted-subsets}.  For the complete identity, \Cref{thm:weighted-generating} gives
\[
 \prod_{i=0}^{n-1}\frac1{1-zq^i}
 =\frac1{(z;q)_n}.
\]
The coefficient of $z^k$ is $\qbinom{n+k-1}{k}$ by \Cref{thm:reciprocal-finite}.

Finally,
\[
 x(x+1)\cdots(x+n)=\sum_{j=0}^{n+1}c(n+1,j)x^j.
\]
The coefficient of $x^{k+1}$ is $e_{n-k}(1,2,\dots,n)$.  The standard ordinary generating function
\[
 \prod_{i=1}^{n}\frac1{1-iz}
 =\sum_{k\ge0}\left\{\begin{matrix}n+k\\n\end{matrix}\right\}z^k
\]
follows from the recurrence for Stirling numbers of the second kind, or by induction on $n$; comparison with \eqref{eq:weighted-h-gen} proves the last formula.
\end{proof}

\section{Principal specialization of Schur polynomials}
\index{Schur polynomial}
\index{principal specialization}

The two Gaussian entries of \Cref{thm:weighted-specializations}(ii) are the
one-column and the one-row case of a single product formula.  Stating it
requires no symmetric-function theory beyond determinants, because a Schur
polynomial can be \emph{defined} by its bialternant.

\begin{definition}\label{def:schur-bialternant}
Let $m\ge1$ and let $\lambda=(\lambda_1\ge\lambda_2\ge\cdots\ge\lambda_m\ge0)$
be a partition with at most $m$ parts, padded by zeros.  Put
\[
 \mu_j=\lambda_j+m-j\quad(1\le j\le m),
 \qquad
 \delta=(m-1,m-2,\ldots,0),
\]
so that $\mu_1>\mu_2>\cdots>\mu_m\ge0$, and $\delta$ is the value of $\mu$ at
$\lambda=0$.  For indeterminates $x_1,\ldots,x_m$ set
\[
 a_\mu(x_1,\ldots,x_m)=\det\bigl(x_i^{\mu_j}\bigr)_{i,j=1}^{m},
\]
and define the \emph{Schur polynomial}
\[
 s_\lambda(x_1,\ldots,x_m)
 =\frac{a_\mu(x_1,\ldots,x_m)}{a_\delta(x_1,\ldots,x_m)}.
\]
Finally put $n(\lambda)=\sum_{j=1}^{m}(j-1)\lambda_j$.
\end{definition}

\begin{lemma}[Vandermonde alternant]\label{lem:vandermonde-alternant}
For indeterminates $y_1,\ldots,y_m$,
\[
 \det\bigl(y_i^{\,j-1}\bigr)_{i,j=1}^{m}
 =\prod_{1\le i<j\le m}(y_j-y_i).
\]
Consequently $a_\delta(x_1,\ldots,x_m)=\prod_{1\le i<j\le m}(x_i-x_j)$, and
$a_\delta$ divides $a_\mu$ in $\Z[x_1,\ldots,x_m]$, so $s_\lambda$ is a
polynomial with integer coefficients.
\end{lemma}

\begin{proof}
Induct on $m$; for $m=1$ both sides equal $1$.  Regard
$D=\det(y_i^{\,j-1})_{i,j=1}^{m}$ as a polynomial in $y_m$ of degree at most
$m-1$.  If $y_m=y_i$ for some $i<m$, two rows coincide and $D=0$; hence
$\prod_{i<m}(y_m-y_i)$ divides $D$, and by the degree bound
$D=c\prod_{i<m}(y_m-y_i)$ with $c$ free of $y_m$.  Expanding along the last
row, the coefficient of $y_m^{\,m-1}$ in $D$ is the cofactor at position
$(m,m)$, namely $\det(y_i^{\,j-1})_{i,j=1}^{m-1}$, which by induction equals
$\prod_{1\le i<j\le m-1}(y_j-y_i)$.  Multiplying proves the formula.

Reversing the order of the columns of $\bigl(x_i^{\,m-j}\bigr)_{i,j}$ turns
the exponents $m-j$ into $j-1$ and multiplies the determinant by the sign
$(-1)^{\binom m2}$ of the reversal permutation, so
\[
 a_\delta=(-1)^{\binom m2}\prod_{1\le i<j\le m}(x_j-x_i)
 =\prod_{1\le i<j\le m}(x_i-x_j).
\]
For the divisibility, interchanging $x_i$ and $x_j$ interchanges two rows of
$a_\mu$, so $a_\mu$ changes sign; setting $x_i=x_j$ therefore gives
$a_\mu=0$, and $x_i-x_j$ divides $a_\mu$ in $\Z[x_1,\ldots,x_m]$.  The
polynomials $x_i-x_j$ with $i<j$ are pairwise non-associate primes in that
unique factorization domain, so their product, which is $a_\delta$, divides
$a_\mu$.
\end{proof}

\begin{theorem}[Principal specialization of a Schur polynomial]\label{thm:schur-principal}
Let $\lambda$ have at most $m$ parts.  Then, as an identity in $\Q(q)$,
\[
 s_\lambda(1,q,q^2,\ldots,q^{m-1})
 =q^{n(\lambda)}
 \prod_{1\le i<j\le m}
 \frac{1-q^{\,\lambda_i-\lambda_j+j-i}}{1-q^{\,j-i}}.
\]
The left side lies in $\Z[q]$, so by \Cref{cor:safe-specialization} the
identity persists at every complex $q$ once the right-hand quotient has been
reduced; as displayed it is to be read at those $q$ with $q^{\,d}\ne1$ for
$1\le d\le m-1$.
\end{theorem}

\begin{proof}
Specialize $x_i=q^{\,i-1}$, so that $x_i^{\mu_j}=q^{(i-1)\mu_j}$.  Writing
$y_j=q^{\mu_j}$, the matrix $\bigl(x_i^{\mu_j}\bigr)_{i,j}$ has entries
$y_j^{\,i-1}$, hence is the transpose of $\bigl(y_i^{\,j-1}\bigr)_{i,j}$.
A determinant is invariant under transposition, so
\Cref{lem:vandermonde-alternant} gives
\[
 a_\mu(1,q,\ldots,q^{m-1})
 =\prod_{1\le i<j\le m}\bigl(q^{\mu_j}-q^{\mu_i}\bigr).
\]
Because $\lambda$ is weakly decreasing we have $\mu_1>\cdots>\mu_m\ge0$, so
$\mu_i-\mu_j>0$ whenever $i<j$, and we extract from each difference the
\emph{smaller} power:
% ed.: extracting q^{\mu_j} rather than q^{\mu_i} produces no sign at all.
% ed.: The incoming survey's `factor the lower powers of q from each
% ed.: difference' left \binom m2 minus signs unaccounted for.
\[
 q^{\mu_j}-q^{\mu_i}=q^{\mu_j}\bigl(1-q^{\,\mu_i-\mu_j}\bigr).
\]
Hence
\[
 a_\mu(1,q,\ldots,q^{m-1})
 =q^{A(\mu)}\prod_{1\le i<j\le m}\bigl(1-q^{\,\mu_i-\mu_j}\bigr),
 \qquad
 A(\mu)=\sum_{1\le i<j\le m}\mu_j=\sum_{j=1}^{m}(j-1)\mu_j,
\]
the last equality because $\mu_j$ is contributed once for each $i<j$.
Applying this to $\mu=\delta$, where $\delta_i-\delta_j=j-i$, gives
\[
 a_\delta(1,q,\ldots,q^{m-1})
 =q^{A(\delta)}\prod_{1\le i<j\le m}\bigl(1-q^{\,j-i}\bigr),
\]
a nonzero element of $\Q(q)$.  Dividing,
\[
 s_\lambda(1,q,\ldots,q^{m-1})
 =q^{\,A(\mu)-A(\delta)}
 \prod_{1\le i<j\le m}
 \frac{1-q^{\,\mu_i-\mu_j}}{1-q^{\,j-i}}.
\]
Finally $\mu_i-\mu_j=\lambda_i-\lambda_j+j-i$ and
\[
 A(\mu)-A(\delta)=\sum_{j=1}^{m}(j-1)(\mu_j-\delta_j)
 =\sum_{j=1}^{m}(j-1)\lambda_j=n(\lambda).
\]
The left side is a polynomial with integer coefficients by
\Cref{lem:vandermonde-alternant}; the transfer to arbitrary complex $q$ is
\Cref{lem:polynomial-identity-principle} together with
\Cref{cor:safe-specialization}.
\end{proof}

\begin{corollary}[The one-column and one-row shadows]\label{cor:schur-e-h}
Let $m=n\ge1$.
\begin{enumerate}[label=\textup{(\roman*)}]
\item For $0\le k\le n$ and $\lambda=(1^k)=(1,\ldots,1,0,\ldots,0)$,
\[
 s_{(1^k)}(1,q,\ldots,q^{n-1})=q^{\binom k2}\qbinom nk.
\]
\item For $k\ge0$ and $\lambda=(k)$,
\[
 s_{(k)}(1,q,\ldots,q^{n-1})=\qbinom{n+k-1}{k}.
\]
\end{enumerate}
Since $s_{(1^k)}=e_k$ and $s_{(k)}=h_k$~\cite{Macdonald1995}, these are
exactly the two Gaussian entries of \Cref{thm:weighted-specializations}(ii),
that is, $C_k^n(\mathbf w)$ and $S_k^n(\mathbf w)$ at $w_i=q^{\,i-1}$.
\end{corollary}

\begin{proof}
(i) Here $\lambda_i=1$ for $i\le k$ and $\lambda_i=0$ for $i>k$, so
$n(\lambda)=\sum_{i\le k}(i-1)=\binom k2$, and $\lambda_i-\lambda_j$ vanishes
except when $i\le k<j$, where it equals $1$.  The product of
\Cref{thm:schur-principal} therefore collapses to
\[
 \prod_{i=1}^{k}\prod_{j=k+1}^{n}
 \frac{1-q^{\,j-i+1}}{1-q^{\,j-i}}
 =\prod_{i=1}^{k}\frac{1-q^{\,n-i+1}}{1-q^{\,k-i+1}},
\]
the inner product telescoping.  The numerators are
$(1-q^{n})(1-q^{n-1})\cdots(1-q^{\,n-k+1})=(q;q)_n/(q;q)_{n-k}$ and the
denominators are $(1-q^k)\cdots(1-q)=(q;q)_k$, so the product is $\qbinom nk$.

(ii) Here $\lambda_1=k$ and $\lambda_j=0$ for $j\ge2$, so $n(\lambda)=0$ and
only the pairs $(1,j)$ contribute:
\[
 \prod_{j=2}^{n}\frac{1-q^{\,k+j-1}}{1-q^{\,j-1}}
 =\frac{(q;q)_{n+k-1}\big/(q;q)_k}{(q;q)_{n-1}}
 =\qbinom{n+k-1}{k}.
 \qedhere
\]
\end{proof}

\begin{remark}\label{rem:schur-hypotheses}
Both hypotheses of \Cref{thm:schur-principal} are needed as stated.  If
$\lambda$ had more than $m$ parts, $\mu$ would not be strictly decreasing and
$a_\mu$ would vanish identically.  And the displayed quotient is meaningful
only where $q^{\,j-i}\ne1$ for all $1\le i<j\le m$; at the excluded roots of
unity, and at $q=1$, one must cancel first, exactly as prescribed by
\Cref{cor:safe-specialization}.  After cancellation the limit $q\to1$ is
\[
 s_\lambda(1,1,\ldots,1)
 =\prod_{1\le i<j\le m}\frac{\lambda_i-\lambda_j+j-i}{j-i},
\]
by $\lim_{q\to1}(1-q^a)/(1-q^b)=a/b$; this is Weyl's dimension formula for
$\mathrm{GL}_m$, and it stands to \Cref{thm:schur-principal} exactly as
\Cref{cor:qbinom-classical} stands to the Gaussian coefficient.
\end{remark}

\section{Weighted inversion}

\begin{theorem}[Symmetric-function inversion]\label{thm:weighted-inversion}
Let $(a_m)$ and $(b_m)$ be sequences.  The transformations
\[
 b_m=\sum_{j=0}^{m}S_{m-j}^n(\mathbf w)a_j
\]
and
\[
 a_m=\sum_{j=0}^{m}(-1)^{m-j}C_{m-j}^n(\mathbf w)b_j
\]
are inverse to each other.
\end{theorem}

\begin{proof}
The generating functions $A(z)=\sum a_mz^m$ and $B(z)=\sum b_mz^m$ satisfy
\[
 B(z)=A(z)\prod_{i=1}^{n}(1-w_i z)^{-1}.
\]
Multiplication by $\prod_i(1-w_i z)$ gives the proposed inverse coefficient formula.  Equivalently, use the orthogonality in \Cref{cor:e-h-orthogonality}.
\end{proof}

\section{The subset--subspace analogy}

The choice $w_i=q^{i-1}$ is more than a convenient specialization.  It replaces the Boolean lattice of subsets by the lattice of subspaces over a finite field.  Ordinary binomial coefficients count subsets by cardinality; Gaussian coefficients count subspaces by dimension.  The two-term weighted recurrence mirrors the decomposition of a subspace according to its relation with a fixed hyperplane.  This analogy is made exact in \Cref{chap:finite-geometry}.

\chapter{Partitions, words, paths, and inversion statistics}\label{chap:combinatorics}
\index{partition!inside a rectangle}
\index{inversion statistic}
\index{binary word}

\section{Binary words}

Let $\mathcal W_{n,k}$ be the set of binary words of length $n$ containing $k$ ones.  Define
\[
 \inv(w)=\#\set{(i,j):i<j,\ w_i=1,\ w_j=0}.
\]

\begin{theorem}[Inversion generating function]\label{thm:binary-inversions}
For $0\le k\le n$,
\[
 \qbinom nk=\sum_{w\in\mathcal W_{n,k}}q^{\inv(w)}.
\]
\end{theorem}

\begin{proof}
Let $F_{n,k}(q)$ denote the right side.  If the last letter is $1$, deletion of it gives a word in $\mathcal W_{n-1,k-1}$ and creates no new inversion.  If the last letter is $0$, deletion gives a word in $\mathcal W_{n-1,k}$ and the final zero forms an inversion with each of the $k$ ones, contributing a factor $q^k$.  Hence
\[
 F_{n,k}=F_{n-1,k-1}+q^kF_{n-1,k}.
\]
The boundary values are $1$, so $F_{n,k}$ and $\qbinom nk$ obey the same recurrence and initial conditions.
\end{proof}

\begin{corollary}
The random variable in \Cref{thm:qbinom-moments} is the inversion number of a uniformly random word in $\mathcal W_{n,k}$.
\end{corollary}

\begin{proof}
The coefficient of $q^j$ in \Cref{thm:binary-inversions} counts words with $j$ inversions, while $\abs{\mathcal W_{n,k}}=\binom nk$.
\end{proof}

\section{Partitions inside a rectangle}

A partition $\lambda=(\lambda_1,\dots,\lambda_k)$ is said to fit in a $k\times m$ rectangle if
\[
 m\ge\lambda_1\ge\cdots\ge\lambda_k\ge0.
\]
Its size is $\abs\lambda=\lambda_1+\cdots+\lambda_k$.

\begin{theorem}[Rectangular partition generating function]\label{thm:rectangle-partitions}
For $m,k\ge0$,
\[
 \Qbinom{m+k}{k}{q}
 =\sum_{\lambda\subseteq k\times m}q^{\abs\lambda}.
\]
\end{theorem}

\begin{proof}
Take a word $w\in\mathcal W_{m+k,k}$ and let the positions of its ones be
\(
 1\le i_1<\cdots<i_k\le m+k.
\)
Define
\[
 \lambda_j=m+j-i_j\qquad(1\le j\le k).
\]
Then $0\le\lambda_j\le m$ and
\[
 \lambda_j-\lambda_{j+1}=i_{j+1}-i_j-1\ge0.
\]
Thus $\lambda$ fits in the rectangle.  Conversely,
\(
 i_j=m+j-\lambda_j
\)
recovers a strictly increasing sequence of positions, so this is a bijection.  The number $\lambda_j$ is exactly the number of zeros to the right of the $j$th one, hence
\(
 \abs\lambda=\inv(w).
\)
Apply \Cref{thm:binary-inversions}.
\end{proof}

\begin{corollary}[Conjugation and complement]\label{cor:partition-symmetries}
Conjugating Ferrers diagrams proves
\(
 \qbinom{m+k}{k}=\qbinom{m+k}{m}.
\)
Taking the complement in the rectangle proves
\[
 \Qbinom{m+k}{k}{q^{-1}}=q^{-mk}\Qbinom{m+k}{k}{q}.
\]
\end{corollary}

\begin{proof}
Conjugation exchanges a $k\times m$ rectangle with an $m\times k$ rectangle without changing size.  Complementation sends size $r$ to $mk-r$, so the generating polynomial is reciprocal of degree $mk$.
\end{proof}

\section{Lattice paths and area}

Encode a binary word by a path from $(0,0)$ to $(m,k)$, reading $0$ as an east step and $1$ as a north step.  With a suitable choice of side, the number of unit squares between the path and the boundary is the inversion number.

\begin{corollary}[Area generating function]\label{cor:path-area}
The polynomial $\Qbinom{m+k}{k}{q}$ is the generating function of monotone lattice paths from $(0,0)$ to $(m,k)$ by the area under the path:
\[
 \Qbinom{m+k}{k}{q}=\sum_P q^{\area(P)}.
\]
\end{corollary}

\begin{proof}
Under the word--path encoding, every north step contributes the number of later east steps, which is precisely the number of squares in its row between the path and the chosen boundary.  Summing over north steps gives $\inv(w)$.
\end{proof}

\section{Durfee squares and the central convolution}

The largest square in the northwest corner of a partition diagram is its Durfee square.

\begin{theorem}[Durfee-square decomposition]\label{thm:durfee-qvand}
For $m,n\ge0$,
\[
 \qbinom{m+n}{m}
 =\sum_{j=0}^{\min(m,n)}q^{j^2}\qbinom mj\qbinom nj.
\]
\end{theorem}

\begin{proof}
This is \Cref{cor:central-vandermonde}; we give the partition proof.  A partition in an $m\times n$ rectangle with Durfee square of side $j$ consists of the $j\times j$ square, a partition to its right fitting in a $j\times(n-j)$ rectangle, and a partition below it fitting in an $(m-j)\times j$ rectangle.  By \Cref{thm:rectangle-partitions}, their generating functions are
\[
 q^{j^2}\Qbinom{n}{j}{q}\Qbinom{m}{j}{q}.
\]
The Durfee size is unique, so summing over $j$ counts every partition exactly once.
\end{proof}

\section{Words with several letters}

Let $\mathcal W(n_1,\dots,n_r)$ be the words containing $n_i$ copies of the letter $i$, ordered by $1<\cdots<r$.  Let $\inv(w)$ count pairs $i<j$ of positions with $w_i>w_j$.

\begin{theorem}[Multiset inversion enumerator]\label{thm:qmult-inversions}
If $n=n_1+\cdots+n_r$, then
\[
 \genfrac{[}{]}{0pt}{}{n}{n_1,\dots,n_r}_q
 =\sum_{w\in\mathcal W(n_1,\dots,n_r)}q^{\inv(w)}.
\]
\end{theorem}

\begin{proof}
Choose the positions of the largest letter $r$.  If their binary indicator has inversion number $s$, then $s$ counts precisely the inversions in which the larger letter is $r$.  The remaining positions contain a word on $1,\dots,r-1$, and its inversions are independent.  Therefore the generating function factors as
\[
 \qbinom n{n_r}
 \genfrac{[}{]}{0pt}{}{n-n_r}{n_1,\dots,n_{r-1}}_q.
\]
Iterating and using \Cref{prop:qmult-factor} proves the formula.
\end{proof}

\section{Rogers--Szeg\H{o} polynomials}

\begin{definition}
The Rogers--Szeg\H{o} polynomial is
\[
 H_n(z;q)=\sum_{k=0}^{n}\qbinom nk z^k.
\]
\end{definition}

\begin{proposition}[Basic recurrence]\label{prop:rogers-szego-recurrence}
One has
\[
 H_{n+1}(z;q)=zH_n(z;q)+H_n(qz;q),
 \qquad H_0(z;q)=1.
\]
Moreover, if $\abs q<1$, then
\[
 \sum_{n=0}^{\infty}\frac{H_n(z;q)}{(q;q)_n}t^n
 =\frac1{(t;q)_\infty(zt;q)_\infty}
\]
for $\abs t<1$ and $\abs{zt}<1$.
\end{proposition}

\begin{proof}
Insert the first $q$-Pascal recurrence into $H_{n+1}$:
\[
 \sum_k\left(\qbinom n{k-1}+q^k\qbinom nk\right)z^k
 =zH_n(z;q)+H_n(qz;q).
\]
For the generating function, use
\[
 \frac{\qbinom nk}{(q;q)_n}=\frac1{(q;q)_k(q;q)_{n-k}}.
\]
Writing $n=k+j$ and separating the double sum gives
\[
 \left(\sum_{j\ge0}\frac{t^j}{(q;q)_j}\right)
 \left(\sum_{k\ge0}\frac{(zt)^k}{(q;q)_k}\right).
\]
Each factor equals the reciprocal infinite product by Euler's identity, proved later in \Cref{cor:euler-infinite}; alternatively, the equality follows already by taking $n\to\infty$ in \Cref{thm:reciprocal-finite}.  This gives the stated product.
\end{proof}

\begin{proposition}[Three-term recurrence]\label{prop:rogers-szego-three-term}
\index{Rogers--Szeg\H{o} polynomial!three-term recurrence}
For $n\ge1$,
\[
 H_{n+1}(z;q)=(1+z)H_n(z;q)-z\,(1-q^n)H_{n-1}(z;q).
\]
\end{proposition}

\begin{proof}
% ed.: the incoming survey proved this by quoting a Pascal II instance in the
% ed.: indices (n-1,k-1),(n-1,k-2), which does not effect the reduction, and
% ed.: then appealing to ``a common factorial denominator''.  The argument below
% ed.: replaces that step by the ratio identity displayed second, after which
% ed.: the comparison is exactly eq:qpascal2.
Both sides are polynomials in $z$, so it suffices to compare the coefficient
of $z^k$ for every $k\in\Z$, using the convention $\qbinom nk=0$ for $k<0$ or
$k>n$.  On the left that coefficient is $\qbinom{n+1}k$; on the right it is
\[
 \qbinom nk+\qbinom n{k-1}-(1-q^n)\qbinom{n-1}{k-1}.
\]
We claim that for every $k\in\Z$ and every $n\ge1$,
\[
 (1-q^n)\qbinom{n-1}{k-1}=\bigl(1-q^{\,n+1-k}\bigr)\qbinom n{k-1}.
\]
Granting the claim, the right-hand coefficient equals
\[
 \qbinom nk+\Bigl(1-\bigl(1-q^{\,n+1-k}\bigr)\Bigr)\qbinom n{k-1}
 =\qbinom nk+q^{\,n+1-k}\qbinom n{k-1},
\]
which is the second $q$-Pascal recurrence \eqref{eq:qpascal2} with $n$
replaced by $n+1$; by \Cref{thm:qpascal} that recurrence is valid for all $k$
under the zero convention, and the right-hand coefficient is therefore
$\qbinom{n+1}k$, as required.

To prove the claim, let first $1\le k\le n$, so that $k-1\ge0$ and
$n-k\ge0$ and both Gaussian coefficients are given by factorial quotients.
From $(q;q)_n=(1-q^n)(q;q)_{n-1}$ and
$(q;q)_{n-k+1}=(1-q^{\,n-k+1})(q;q)_{n-k}$,
\[
 \frac{\qbinom n{k-1}}{\qbinom{n-1}{k-1}}
 =\frac{(q;q)_n\,(q;q)_{n-k}}{(q;q)_{n-1}\,(q;q)_{n-k+1}}
 =\frac{1-q^n}{1-q^{\,n-k+1}},
\]
and cross-multiplication gives the claim.  For $k\le0$ both sides vanish,
since then $k-1<0$ and both Gaussian coefficients are zero.  For $k=n+1$ the
left side vanishes because $\qbinom{n-1}{n}=0$ and the right side vanishes
because $1-q^{\,n+1-k}=1-q^0=0$.  For $k>n+1$ every term is zero.  This
exhausts all $k$ and completes the proof.
\end{proof}

\begin{corollary}[Dilation formula]\label{cor:rogers-szego-dilation}
For $n\ge1$,
\[
 H_n(qz;q)=H_n(z;q)-z\,(1-q^n)H_{n-1}(z;q).
\]
\end{corollary}

\begin{proof}
By \Cref{prop:rogers-szego-recurrence}, $H_n(qz;q)=H_{n+1}(z;q)-zH_n(z;q)$.
Substituting \Cref{prop:rogers-szego-three-term} for $H_{n+1}(z;q)$ and
cancelling $zH_n(z;q)$ gives the formula.
\end{proof}

\begin{remark}\label{rem:rogers-szego-three-term}
The recurrence of \Cref{prop:rogers-szego-recurrence} moves the argument, from
$z$ to $qz$; the one above does not.  It is therefore a three-term recurrence
in the index $n$ alone, of the shape familiar from orthogonal polynomial
sequences, and it evaluates $H_0,H_1,\ldots,H_N$ at a fixed numerical $z$ in
$\OO(N)$ ring operations, without ever forming the Gaussian coefficients.  At
$q=0$ every $\qbinom nk$ equals $1$ by \Cref{prop:extreme-specializations},
the recurrence becomes $H_{n+1}=(1+z)H_n-zH_{n-1}$, and its solution is the
geometric sum $1+z+\cdots+z^{n}$, as it must be.
\end{remark}

\chapter{Finite fields, Grassmannians, and rank enumeration}\label{chap:finite-geometry}
\index{finite field}
\index{Grassmannian}
\index{Schubert cell}

\section{Counting subspaces}

\begin{theorem}[Gaussian coefficients count subspaces]\label{thm:subspace-count}
Let $Q$ be a prime power.  The number of $k$-dimensional linear subspaces of $\F_Q^n$ is
\[
 \Qbinom nkQ.
\]
\end{theorem}

\begin{proof}
Count ordered linearly independent $k$-tuples in $\F_Q^n$.  The first vector may be any nonzero vector, the second any vector outside the span of the first, and so on, giving
\[
 (Q^n-1)(Q^n-Q)\cdots(Q^n-Q^{k-1}).
\]
Each $k$-dimensional subspace has
\[
 (Q^k-1)(Q^k-Q)\cdots(Q^k-Q^{k-1})
\]
ordered bases.  Dividing and cancelling $Q^i$ in corresponding factors yields
\[
 \prod_{i=0}^{k-1}\frac{Q^{n-i}-1}{Q^{k-i}-1}
 =\prod_{j=1}^{k}\frac{1-Q^{n-k+j}}{1-Q^j}
 =\Qbinom nkQ.
\]
\end{proof}

\begin{theorem}[Geometric Pascal recurrence]\label{thm:geometric-pascal}
Let $H$ be a fixed hyperplane in $\F_Q^n$.  Then
\[
 \Qbinom nkQ
 =\Qbinom{n-1}{k}Q+Q^{n-k}\Qbinom{n-1}{k-1}Q.
\]
\end{theorem}

\begin{proof}
A $k$-subspace $U$ is either contained in $H$ or not.  The contained subspaces contribute $\Qbinom{n-1}{k}Q$.  If $U\not\subseteq H$, then $W=U\cap H$ has dimension $k-1$.  Fixing $W$, the quotient $\F_Q^n/W$ has dimension $n-k+1$, and $H/W$ is a hyperplane.  The subspace $U/W$ is a line not contained in $H/W$.  The total number of lines in $\F_Q^{n-k+1}$ minus those in its hyperplane is
\[
 \frac{Q^{n-k+1}-1}{Q-1}-\frac{Q^{n-k}-1}{Q-1}=Q^{n-k}.
\]
There are $\Qbinom{n-1}{k-1}Q$ choices for $W$, proving the formula.
\end{proof}

\section{Flags and \texorpdfstring{$q$}{q}-multinomial coefficients}

\begin{theorem}[Partial flags]\label{thm:flag-count}
Let $d_i=n_1+\cdots+n_i$.  The number of flags
\[
 0=V_0\subset V_1\subset\cdots\subset V_r=\F_Q^n,
 \qquad \dim(V_i/V_{i-1})=n_i,
\]
is
\[
 \genfrac{[}{]}{0pt}{}{n}{n_1,\dots,n_r}_{Q}.
\]
\end{theorem}

\begin{proof}
Choose $V_1$ in $\Qbinom n{n_1}Q$ ways.  In the quotient $\F_Q^n/V_1$, choose $V_2/V_1$ in $\Qbinom{n-n_1}{n_2}Q$ ways, and continue.  The product is the factorization in \Cref{prop:qmult-factor}.
\end{proof}

\section{Matrices of fixed rank}

\begin{theorem}[Rank enumeration]\label{thm:rank-matrices}
The number of $m\times n$ matrices of rank $r$ over $\F_Q$ is
\[
 N_{m,n,r}(Q)
 =\Qbinom mrQ\Qbinom nrQ\abs{\mathrm{GL}_r(\F_Q)},
\]
where
\[
 \abs{\mathrm{GL}_r(\F_Q)}=\prod_{i=0}^{r-1}(Q^r-Q^i).
\]
Equivalently,
\[
 N_{m,n,r}(Q)
 =\frac{\prod_{i=0}^{r-1}(Q^m-Q^i)(Q^n-Q^i)}
 {\prod_{i=0}^{r-1}(Q^r-Q^i)}.
\]
\end{theorem}

\begin{proof}
Every rank-$r$ matrix $A$ can be factored as $A=UV$, where $U$ is an $m\times r$ matrix of full column rank and $V$ is an $r\times n$ matrix of full row rank.  There are
\(
 \prod_{i=0}^{r-1}(Q^m-Q^i)
\)
choices for $U$ and
\(
 \prod_{i=0}^{r-1}(Q^n-Q^i)
\)
choices for $V$.  Two pairs $(U,V)$ and $(U',V')$ give the same $A$ exactly when
\[
 U'=UG,\qquad V'=G^{-1}V
\]
for a unique $G\in\mathrm{GL}_r(\F_Q)$.  Thus every matrix is counted $\abs{\mathrm{GL}_r}$ times.  Substituting the subspace-count formulas gives the first expression.
\end{proof}

\section{Schubert cells and a topological meaning}

Let $\mathrm{Gr}(k,n)$ be the complex Grassmannian of $k$-planes in $\C^n$.  Its Schubert cell decomposition is indexed by partitions inside a $k\times(n-k)$ rectangle.

\begin{theorem}[Poincar\'e polynomial of the Grassmannian]\label{thm:grassmann-poincare}
The integral cohomology of $\mathrm{Gr}(k,n)$ is concentrated in even degrees, and its Poincar\'e polynomial is
\[
 P_{\mathrm{Gr}(k,n)}(t)
 =\sum_{j\ge0}\operatorname{rank}H^{2j}(\mathrm{Gr}(k,n);\Z)t^{2j}
 =\Qbinom nk{t^2}.
\]
\end{theorem}

\begin{proof}
Fix the standard flag in $\C^n$.  Every $k$-plane has a unique row-reduced echelon matrix.  Its pivot positions determine a binary word with $k$ ones, hence a partition $\lambda$ inside the $k\times(n-k)$ rectangle.  Once the pivots are fixed, the entries allowed to vary freely occupy exactly the boxes of $\lambda$; the corresponding stratum is therefore an affine complex space $\C^{\abs\lambda}$.  Ordering the strata by inclusion of diagrams yields a CW decomposition with one real cell of dimension $2\abs\lambda$ for each $\lambda$.  Since all cells have even dimension, every cellular boundary map is zero.  Thus the rank of $H^{2j}$ is the number of such partitions of size $j$.  Apply \Cref{thm:rectangle-partitions} with $q=t^2$.
\end{proof}

\begin{remark}[The ``same polynomial'' principle]
For a prime power $Q$, evaluating $\qbinom nk$ at $q=Q$ counts rational points of a Grassmannian over $\F_Q$.  Evaluating it at $q=t^2$ records the even Betti numbers of the complex Grassmannian.  This is an elementary instance of the broad principle that polynomial point counts over finite fields often shadow topological invariants over $\C$.
\end{remark}

\section{Growth over finite fields}

\begin{proposition}[Dimension-dominant bounds]\label{prop:qbinom-growth}
For $Q>1$ and $0\le k\le n$,
\[
 Q^{k(n-k)}
 \le \Qbinom nkQ
 \le \frac{Q^{k(n-k)}}{(Q^{-1};Q^{-1})_\infty}.
\]
\end{proposition}

\begin{proof}
Rewrite the product as
\[
 \Qbinom nkQ
 =Q^{k(n-k)}
 \prod_{j=1}^{k}\frac{1-Q^{-(n-k+j)}}{1-Q^{-j}}.
\]
Each ratio is at least $1$ because $n-k+j\ge j$, proving the lower bound.  For the upper bound, discard the numerator factors and note
\[
 \prod_{j=1}^{k}(1-Q^{-j})\ge\prod_{j=1}^{\infty}(1-Q^{-j})=(Q^{-1};Q^{-1})_\infty.
\]
\end{proof}

\chapter{Symmetry and unimodality}\label{chap:unimodality}
\index{unimodality}
\index{$\mathfrak{sl}_2$}

Palindromicity follows immediately from reciprocity.  Unimodality is subtler: the coefficients rise toward the middle and then fall.  A concise conceptual proof comes from the representation theory of $\mathfrak{sl}_2$.

\section{A representation-theoretic lemma}

Let $\mathfrak{sl}_2(\C)$ have generators $E,F,H$ with
\[
 [H,E]=2E,\qquad [H,F]=-2F,\qquad [E,F]=H.
\]

\begin{lemma}[Structure of finite-dimensional $\mathfrak{sl}_2$-modules]\label{lem:sl2-structure}
Every finite-dimensional complex $\mathfrak{sl}_2$-module is a direct sum of irreducibles $L_d$ indexed by $d\in\N_0$.  The module $L_d$ has dimension $d+1$ and $H$-weights
\[
 d,d-2,\dots,-d,
\]
each with multiplicity one.
\end{lemma}

\begin{proof}
We include the standard argument.  Restrict a representation to the compact real form $\mathfrak{su}_2$.  A finite-dimensional Lie-algebra representation integrates to the simply connected compact group $\mathrm{SU}(2)$ by exponentiating the representing matrices; the Lie relations ensure consistency of the local exponentials, and simple connectedness removes monodromy.  Starting from any Hermitian inner product, average it over $\mathrm{SU}(2)$ with normalized Haar measure.  The result is invariant.  Hence the orthogonal complement of every invariant subspace is invariant, proving complete reducibility.

It remains to classify an irreducible module $V$.  In the invariant Hermitian structure, a suitable scalar multiple of $H$ is self-adjoint, so $H$ is diagonalizable.  Choose an eigenvector $v$ whose eigenvalue $d$ has maximal real part.  Since $E$ raises an $H$-weight by $2$, maximality gives $Ev=0$.  The commutation relations imply by induction
\[
 H F^jv=(d-2j)F^jv,
 \qquad
 E F^jv=j(d-j+1)F^{j-1}v.
\]
Because $V$ is finite-dimensional, let $m$ be maximal with $F^mv\ne0$.  Applying $E$ to $F^{m+1}v=0$ gives
\[
 (m+1)(d-m)F^mv=0,
\]
so $d=m\in\N_0$.  The vectors $v,Fv,\dots,F^dv$ have distinct weights and are linearly independent.  Their span is invariant under $E,F,H$ and is nonzero; irreducibility makes it all of $V$.  The displayed weights and multiplicities follow.
\end{proof}

\section{Exterior powers and Gaussian coefficients}

Let $L_{n-1}$ have a weight basis $v_0,\dots,v_{n-1}$ with
\(
 Hv_i=(n-1-2i)v_i.
\)

\begin{proposition}[Exterior-power character]\label{prop:exterior-character}
For $0\le k\le n$,
\[
 \ch\!\left(\bigwedge^kL_{n-1}\right)(x)
 =x^{k(n-k)}\Qbinom nk{x^{-2}}.
\]
\end{proposition}

\begin{proof}
A basis of the exterior power consists of
\(
 v_{i_1}\wedge\cdots\wedge v_{i_k}
\)
with $0\le i_1<\cdots<i_k\le n-1$.  Its weight is
\[
 k(n-1)-2(i_1+\cdots+i_k).
\]
Therefore
\[
 \ch=x^{k(n-1)}
 \sum_{i_1<\cdots<i_k}(x^{-2})^{i_1+\cdots+i_k}.
\]
Use \Cref{cor:weighted-subsets} with $q=x^{-2}$.  The prefactor becomes
\[
 x^{k(n-1)}x^{-2\binom k2}=x^{k(n-k)}.
\]
\end{proof}

\begin{theorem}[Unimodality of Gaussian coefficients]\label{thm:qbinom-unimodal}
If
\[
 \qbinom nk=c_0+c_1q+\cdots+c_Dq^D,
 \qquad D=k(n-k),
\]
then
\[
 c_0\le c_1\le\cdots\le c_{\floor{D/2}}
\]
and $c_j=c_{D-j}$.  Thus the coefficient sequence is symmetric and unimodal.
\end{theorem}

\begin{proof}
Symmetry was proved in \Cref{thm:qbinom-structure}.  By \Cref{lem:sl2-structure}, decompose
\[
 \bigwedge^kL_{n-1}\cong\bigoplus_{d\ge0}m_dL_d.
\]
All weights in this exterior power have the same parity as $D$.  The character of one summand $L_d$, placed inside the ambient sequence of weights
\[
 D,D-2,\dots,-D,
\]
contributes a consecutive block of ones centered at weight $0$: it is zero until weight $d$ is reached, then one at every step down to $-d$, then zero again.  As one moves from weight $D$ toward $0$, such a block can turn on but cannot turn off.  A nonnegative sum of these centered blocks therefore has weakly increasing multiplicities toward the middle.

By \Cref{prop:exterior-character}, the multiplicity of weight $D-2j$ is exactly $c_j$.  Hence $c_0\le c_1\le\cdots\le c_{\floor{D/2}}$.
\end{proof}

\begin{remark}
The proof explains more than the inequality: the successive differences $c_j-c_{j-1}$ record multiplicities of irreducible summands with prescribed highest weight.  Stronger strict-unimodality results require additional arguments and have a small family of exceptional cases; ordinary unimodality is exactly what the $\mathfrak{sl}_2$ decomposition supplies for free.
\end{remark}


\part{Infinite series and basic hypergeometric summation}

\chapter{The infinite \texorpdfstring{$q$}{q}-binomial theorem}\label{chap:infinite-qbinomial}
\index{$q$-binomial theorem!infinite}

\section{A functional-equation proof}

\begin{theorem}[Infinite $q$-binomial theorem]\label{thm:infinite-qbinomial}
Let $\abs q<1$ and $\abs z<1$.  Then
\[
 \sum_{n=0}^{\infty}\frac{(a;q)_n}{(q;q)_n}z^n
 =\frac{(az;q)_\infty}{(z;q)_\infty}.
\]
The convergence is locally uniform in $(a,z)$ on compact subsets of $\C\times\set{\abs z<1}$.
\end{theorem}

\begin{proof}
Set
\[
 F(z)=\sum_{n\ge0}c_nz^n,
 \qquad c_n=\frac{(a;q)_n}{(q;q)_n}.
\]
Since $c_{n}/c_{n-1}=(1-aq^{n-1})/(1-q^n)\to1$, the radius of convergence is $1$ unless the series terminates; in every case it converges for $\abs z<1$.  The coefficient recurrence is equivalent to
\[
 (1-q^n)c_n=(1-aq^{n-1})c_{n-1}.
\]
Comparing coefficients shows that
\[
 (1-z)F(z)=(1-az)F(qz),\qquad F(0)=1.\label{eq:qbinom-fe}
\]
Now put
\(
 G(z)=(az;q)_\infty/(z;q)_\infty.
\)
Using $(z;q)_\infty=(1-z)(qz;q)_\infty$ gives the same functional equation and $G(0)=1$.

To prove uniqueness, iterate the functional equation:
\[
 F(z)=\prod_{j=0}^{N-1}\frac{1-azq^j}{1-zq^j}F(q^Nz).
\]
As $N\to\infty$, $q^Nz\to0$, so $F(q^Nz)\to1$.  The finite product tends to $G(z)$.  Hence $F=G$.
\end{proof}

\begin{corollary}[Euler's two infinite identities]\label{cor:euler-infinite}
For $\abs q<1$,
\begin{align}
 (z;q)_\infty
 &=\sum_{n=0}^{\infty}\frac{(-1)^nq^{\binom n2}}{(q;q)_n}z^n,
 \qquad z\in\C,\label{eq:euler-product-series}\\
 \frac1{(z;q)_\infty}
 &=\sum_{n=0}^{\infty}\frac{z^n}{(q;q)_n},
 \qquad \abs z<1.\label{eq:euler-reciprocal-series}
\end{align}
\end{corollary}

\begin{proof}
Set $a=0$ in \Cref{thm:infinite-qbinomial} to obtain \eqref{eq:euler-reciprocal-series}.  For \eqref{eq:euler-product-series}, start with the finite identity
\[
 (z;q)_N=\sum_{n=0}^{N}(-1)^nq^{\binom n2}\Qbinom Nnq z^n.
\]
For fixed $n$, $\Qbinom Nnq\to1/(q;q)_n$ as $N\to\infty$.  The limiting series converges for every $z$ because its successive-term ratio has magnitude asymptotic to $\abs z\abs q^n$.  To justify passage to the limit uniformly on $\abs z\le R$, note that
\[
 \left|\Qbinom Nnq\right|
 \le
 \frac{\prod_{j\ge1}(1+\abs q^j)}
      {\bigl(\inf_{r\ge0}\abs{(q;q)_r}\bigr)^{2}},
% ed.: squared the infimum --- the Gaussian coefficient has two
% ed.: shifted-factorial factors in its denominator, so the derivation
% ed.: bounds each by the infimum separately.
\]
where the denominator is positive because $(q;q)_r\to(q;q)_\infty\ne0$.  Thus the summands, extended by zero for $n>N$, are dominated by a constant multiple of $\abs q^{\binom n2}R^n$.  Dominated convergence, together with local uniform convergence of the products, now permits passage to the limit.
\end{proof}

\begin{corollary}[A ratio expansion]\label{cor:ratio-expansion}
For $\abs z<1$,
\[
 \frac{(a_1z;q)_\infty\cdots(a_rz;q)_\infty}
 {(b_1z;q)_\infty\cdots(b_sz;q)_\infty}
\]
can be expanded by multiplying $r$ copies of \eqref{eq:euler-product-series} and $s$ copies of \eqref{eq:euler-reciprocal-series}.  In particular, every coefficient is a finite convolution of $q$-Pochhammer ratios.
\end{corollary}

\begin{proof}
Each numerator factor has an entire series expansion and each reciprocal denominator factor converges for $\abs{b_jz}<1$.  Cauchy multiplication is valid in their common disk of absolute convergence.
\end{proof}

\section{Formal and analytic interpretations}

\begin{proposition}[Formal version]\label{prop:formal-qbinomial}
If $q$ is an indeterminate, then
\[
 \frac{(az;q)_\infty}{(z;q)_\infty}
 =\sum_{n\ge0}\frac{(a;q)_n}{(q;q)_n}z^n
\]
holds in the $q$-adically complete ring $\Q(a)[[q]][[z]]$.  More generally, it remains valid after any specialization into a complete topological ring for which the displayed products and series converge.
\end{proposition}

\begin{proof}
The products are well defined $q$-adically: apart from their initial factors, the $j$th factors differ from $1$ by a multiple of $q^j$.  For a formal power series $F(z)=\sum c_nz^n$, the equation
\[
 (1-z)F(z)=(1-az)F(qz),\qquad c_0=1,
\]
determines $c_n$ recursively by
\(
 (1-q^n)c_n=(1-aq^{n-1})c_{n-1},
\)
because $1-q^n$ is a unit in $\Q(a)[[q]]$.  Both sides satisfy this equation, so their coefficients agree.  The same uniqueness argument works in every stated complete specialization.
\end{proof}

\section{Limits in the number of factors}

\begin{proposition}\label{prop:fixed-k-limit}
For fixed $k$ and $\abs q<1$,
\[
 \lim_{n\to\infty}\qbinom nk=\frac1{(q;q)_k}.
\]
More generally,
\[
 \lim_{n\to\infty}\qbinom{n+k}{k}=\frac1{(q;q)_k}.
\]
\end{proposition}

\begin{proof}
Use
\[
 \qbinom nk=\frac{(q^{n-k+1};q)_k}{(q;q)_k}.
\]
Every numerator factor tends to $1$.
\end{proof}

\begin{remark}
The two Euler identities are therefore the stable limits of the finite product and reciprocal-product formulas.  This stability is one reason finite $q$-binomial identities so often serve as polynomial approximations to infinite $q$-series identities.
\end{remark}

\chapter{Basic hypergeometric series}\label{chap:basic-hypergeometric}
\index{basic hypergeometric series}

\section{Definition and convergence}

\begin{definition}
The basic hypergeometric series is
\[
 {}_r\phi_s\!\left[
 \begin{matrix}a_1,\dots,a_r\\ b_1,\dots,b_s\end{matrix};q,z\right]
 =\sum_{n=0}^{\infty}
 \frac{(a_1,\dots,a_r;q)_n}{(q,b_1,\dots,b_s;q)_n}
 \left((-1)^nq^{\binom n2}\right)^{1+s-r}z^n,
\]
provided no denominator factor vanishes.  If one numerator parameter is $q^{-N}$, the series terminates at $n=N$.
\end{definition}

\begin{proposition}[Elementary convergence classification]\label{prop:phi-convergence}
Assume the series does not terminate and the denominator parameters avoid poles.
\begin{enumerate}[label=\textup{(\roman*)}]
\item If $r\le s$, the series converges for every $z$.
\item If $r=s+1$, its radius of convergence is $1$.
\item If $r>s+1$, the terms generally fail to tend to zero for $z\ne0$; such series are normally interpreted only when terminating or formally.
\end{enumerate}
\end{proposition}

\begin{proof}
Let $T_n$ be the $n$th summand.  Since $(1-aq^n)\to1$ for each fixed parameter,
\[
 \frac{T_{n+1}}{T_n}
 =z\,\frac{\prod_{i=1}^{r}(1-a_iq^n)}
 {(1-q^{n+1})\prod_{j=1}^{s}(1-b_jq^n)}
 (-1)^{1+s-r}q^{n(1+s-r)}.
\]
The ratio tends to $0$ when $r\le s$, to $z$ when $r=s+1$, and has unbounded magnitude when $r>s+1$ and $z\ne0$.  Apply the ratio test.
\end{proof}

The infinite $q$-binomial theorem is the evaluation
\[
 {}_1\phi_0\!\left[\begin{matrix}a\\-\end{matrix};q,z\right]
 =\frac{(az;q)_\infty}{(z;q)_\infty}.
\]

\section{Heine's transformation}

\begin{theorem}[Heine transformation]\label{thm:heine}
Initially for $\abs z<1$ and $\abs b<1$,
\[
 {}_2\phi_1\!\left[\begin{matrix}a,b\\c\end{matrix};q,z\right]
 =\frac{(b,az;q)_\infty}{(c,z;q)_\infty}
 {}_2\phi_1\!\left[\begin{matrix}c/b,z\\az\end{matrix};q,b\right].
\]
The identity extends wherever both sides are defined by analytic continuation.
\end{theorem}

\begin{proof}
Use the tail-factor identity to write
\[
 \frac{(b;q)_n}{(c;q)_n}
 =\frac{(b;q)_\infty}{(c;q)_\infty}
 \frac{(cq^n;q)_\infty}{(bq^n;q)_\infty}.
\]
By the infinite $q$-binomial theorem,
\[
 \frac{(cq^n;q)_\infty}{(bq^n;q)_\infty}
 =\sum_{k\ge0}\frac{(c/b;q)_k}{(q;q)_k}b^kq^{nk}.
\]
Insert this into the ${}_2\phi_1$ series and interchange the absolutely convergent sums:
\[
 \frac{(b;q)_\infty}{(c;q)_\infty}
 \sum_{k\ge0}\frac{(c/b;q)_k}{(q;q)_k}b^k
 \sum_{n\ge0}\frac{(a;q)_n}{(q;q)_n}(zq^k)^n.
\]
The inner sum is
\[
 \frac{(azq^k;q)_\infty}{(zq^k;q)_\infty}
 =\frac{(az;q)_\infty}{(z;q)_\infty}
 \frac{(z;q)_k}{(az;q)_k}.
\]
Collecting factors yields the claimed transformation.  Both sides are meromorphic functions of the parameters, so equality extends from the initial open domain by the identity theorem.
\end{proof}

\section{The \texorpdfstring{$q$}{q}-Gauss sum}

\begin{theorem}[$q$-Gauss summation]\label{thm:q-gauss}
If $\abs{c/(ab)}<1$, then
\[
 {}_2\phi_1\!\left[
 \begin{matrix}a,b\\c\end{matrix};q,\frac{c}{ab}\right]
 =\frac{(c/a,c/b;q)_\infty}{(c,c/(ab);q)_\infty}.
\]
\end{theorem}

\begin{proof}
Set $z=c/(ab)$ in \Cref{thm:heine}.  Then $az=c/b$, so one numerator parameter of the transformed series cancels its denominator parameter:
\[
 {}_2\phi_1\!\left[\begin{matrix}c/b,c/(ab)\\c/b\end{matrix};q,b\right]
 ={}_1\phi_0\!\left[\begin{matrix}c/(ab)\\-\end{matrix};q,b\right]
 =\frac{(c/a;q)_\infty}{(b;q)_\infty}.
\]
Multiplication by the prefactor in Heine's transformation gives the result.  The initially stronger restrictions are removed by analytic continuation within the convergence domain.
\end{proof}

\section{A second finite Cauchy identity}

\begin{lemma}[Cauchy convolution II]\label{lem:cauchy-II}
For $a,b,c$ in a commutative ring,
\[
 (ac;q)_n(bc;q)_n
 =\sum_{k=0}^{n}\qbinom nk
 (a;q)_k(b;q)_k c^k(c;q)_{n-k}(abcq^k;q)_{n-k}.
\]
\end{lemma}

\begin{proof}
Apply the $q$-Cauchy identity \Cref{thm:q-cauchy} with $u=bq^k$, $v=ac$, and length $n-k$:
\[
 (abcq^k;q)_{n-k}
 =\sum_{r=0}^{n-k}\qbinom{n-k}{r}
 (bq^k;q)_r(ac)^r(ac;q)_{n-k-r}.
\]
Insert this in the right side of the proposed identity.  Use
\[
 \qbinom nk\qbinom{n-k}{r}
 =\qbinom n{k+r}\qbinom{k+r}k,
 \qquad
 (b;q)_k(bq^k;q)_r=(b;q)_{k+r}.
\]
Put $m=k+r$ and use
\[
 (c;q)_{n-k}=(c;q)_{n-m}(cq^{n-m};q)_{m-k}.
\]
The double sum becomes
\[
 \sum_{m=0}^{n}\qbinom nm(b;q)_m c^m(c;q)_{n-m}(ac;q)_{n-m}
 \sum_{k=0}^{m}\qbinom mk(a;q)_k a^{m-k}
 (cq^{n-m};q)_{m-k}.
\]
In the inner sum replace $j=m-k$ and apply \Cref{thm:q-cauchy} with $u=cq^{n-m}$ and $v=a$; it equals $(acq^{n-m};q)_m$.  Hence
\[
 (ac;q)_{n-m}(acq^{n-m};q)_m=(ac;q)_n
\]
can be factored out.  The remaining sum is
\[
 (ac;q)_n\sum_{m=0}^{n}\qbinom nm(b;q)_m c^m(c;q)_{n-m}
 =(ac;q)_n(bc;q)_n
\]
by one final application of the $q$-Cauchy identity.
\end{proof}

\section{The \texorpdfstring{$q$}{q}-Pfaff--Saalsch\"utz sum}

\begin{theorem}[$q$-Pfaff--Saalsch\"utz summation]\label{thm:q-pfaff}
For $n\in\N_0$,
\[
 {}_3\phi_2\!\left[
 \begin{matrix}a,b,q^{-n}\\c,abq^{1-n}/c\end{matrix};q,q\right]
 =\frac{(c/a,c/b;q)_n}{(c,c/(ab);q)_n}.
\]
\end{theorem}

\begin{proof}
In \Cref{lem:cauchy-II}, replace its parameters by
\[
 a\mapsto a,\qquad b\mapsto b,\qquad c\mapsto \frac{c}{ab}.
\]
This gives
\[
 \frac{(c/a,c/b;q)_n}{(c,c/(ab);q)_n}
 =\sum_{k=0}^{n}\qbinom nk(a,b;q)_k
 \left(\frac{c}{ab}\right)^k
 \frac{(c/(ab);q)_{n-k}(cq^k;q)_{n-k}}
 {(c,c/(ab);q)_n}.
\]
We show that the $k$th summand is the $k$th term of the stated ${}_3\phi_2$.  The tail identities give
\[
 \frac{(cq^k;q)_{n-k}}{(c;q)_n}=\frac1{(c;q)_k}.
\]
Also, reversal of the two finite products gives
\[
 (q^{-n};q)_k
 =(-1)^kq^{-nk+\binom k2}(q^{n-k+1};q)_k,
\]
\[
 (abq^{1-n}/c;q)_k
 =\left(-\frac{ab}{c}\right)^k
 q^{k(1-n)+\binom k2}
 (cq^{n-k}/(ab);q)_k.
\]
Therefore
\begin{align*}
 &\frac{(q^{-n};q)_k q^k}{(q;q)_k(abq^{1-n}/c;q)_k}
 \frac{(c/(ab);q)_n}{(c/(ab);q)_{n-k}}\\
 &\quad=\left(\frac{c}{ab}\right)^k
 \frac{(q^{n-k+1};q)_k}{(q;q)_k}
 \frac{(c/(ab);q)_n}
 {(c/(ab);q)_{n-k}(cq^{n-k}/(ab);q)_k}\\
 &\quad=\qbinom nk\left(\frac{c}{ab}\right)^k.
\end{align*}
Substitution yields exactly
\[
 \sum_{k=0}^{n}
 \frac{(a,b,q^{-n};q)_k}{(q,c,abq^{1-n}/c;q)_k}q^k.
\]
\end{proof}

\section{The two \texorpdfstring{$q$}{q}-Chu--Vandermonde sums}

\begin{corollary}[$q$-Chu--Vandermonde]\label{cor:q-chu}
For $n\in\N_0$,
\begin{align}
 {}_2\phi_1\!\left[
 \begin{matrix}q^{-n},a\\c\end{matrix};q,\frac{cq^n}{a}\right]
 &=\frac{(c/a;q)_n}{(c;q)_n},\label{eq:qchu1}\\
 {}_2\phi_1\!\left[
 \begin{matrix}q^{-n},a\\c\end{matrix};q,q\right]
 &=a^n\frac{(c/a;q)_n}{(c;q)_n}.
 \label{eq:qchu2}
\end{align}
\end{corollary}

\begin{proof}
For \eqref{eq:qchu1}, put $b=q^{-n}$ in the $q$-Gauss sum.  The infinite products telescope:
\[
 \frac{(cq^n,c/a;q)_\infty}{(c,cq^n/a;q)_\infty}
 =\frac{(c/a;q)_n}{(c;q)_n}.
\]
Because the series terminates, no convergence restriction remains.

For \eqref{eq:qchu2}, let $b\to0$ in \Cref{thm:q-pfaff}, with its other numerator parameter equal to $a$.  Then $(b;q)_k\to1$ and $(abq^{1-n}/c;q)_k\to1$, so the left side becomes the displayed ${}_2\phi_1$.  On the right,
\[
 \frac{(c/b;q)_n}{(c/(ab);q)_n}\longrightarrow a^n
\]
by comparing the leading factors of the two degree-$n$ polynomials in $1/b$.
\end{proof}

\section{A hierarchy of mechanisms}

The proofs in this chapter form a hierarchy:
\[
 \text{finite $q$-Pascal}
 \Longrightarrow \text{$q$-Cauchy}
 \Longrightarrow \text{$q$-Pfaff--Saalsch\"utz},
\]
while
\[
 \text{infinite $q$-binomial}
 \Longrightarrow \text{Heine}
 \Longrightarrow \text{$q$-Gauss}
 \Longrightarrow \text{first $q$-Chu--Vandermonde}.
\]
The finite chain is valid in a commutative ring after denominators are cleared.  The infinite chain requires convergence or a suitable complete formal topology.  This distinction becomes important in computer formalization; see \Cref{chap:formalization}.

\chapter{Bilateral series and Ramanujan's \texorpdfstring{${}_1\psi_1$}{1psi1} sum}\label{chap:bilateral}
\index{Ramanujan ${}_1\psi_1$ summation}
\index{bilateral basic hypergeometric series}

\section{Definition}

For integral $n$, negative $q$-shifted factorials are understood as in \Cref{prop:negative-index}.  Define
\[
 {}_1\psi_1\!\left[\begin{matrix}a\\b\end{matrix};q,z\right]
 =\sum_{n\in\Z}\frac{(a;q)_n}{(b;q)_n}z^n.
\]
The positive tail behaves like a geometric series in $z$, while the negative tail behaves like a geometric series in $b/(az)$.

\begin{proposition}[Convergence annulus]\label{prop:1psi1-convergence}
If denominator factors do not vanish, the ${}_1\psi_1$ series converges absolutely in
\[
 \abs{b/a}<\abs z<1.
\]
\end{proposition}

\begin{proof}
For $n\to+\infty$, the ratio of consecutive terms tends to $z$.  For the negative tail put $n=-m$ and use the negative-index formula.  The ratio of successive terms in $m$ tends to $b/(az)$.  The ratio test gives the annulus.
\end{proof}

\begin{theorem}[Ramanujan's ${}_1\psi_1$ summation]\label{thm:1psi1}
For $\abs{b/a}<\abs z<1$,
\[
 {}_1\psi_1\!\left[\begin{matrix}a\\b\end{matrix};q,z\right]
 =\frac{(q,b/a,az,q/(az);q)_\infty}
 {(b,q/a,z,b/(az);q)_\infty}.
\]
\end{theorem}

\begin{proof}
We first prove the formula for $b=q^{m+1}$, where $m\in\N_0$.  In that case $1/(b;q)_n=0$ for $n<-m$ in the limiting sense, so shift $n=k-m$:
\begin{align*}
 \sum_{n\in\Z}\frac{(a;q)_n}{(q^{m+1};q)_n}z^n
 &=z^{-m}\frac{(a;q)_{-m}}{(q^{m+1};q)_{-m}}
 \sum_{k=0}^{\infty}\frac{(aq^{-m};q)_k}{(q;q)_k}z^k\\
 &=z^{-m}\frac{(q;q)_m}{(aq^{-m};q)_m}
 \frac{(aq^{-m}z;q)_\infty}{(z;q)_\infty}.
\end{align*}
The second equality is the infinite $q$-binomial theorem.  Reverse the finite products:
\[
 (aq^{-m};q)_m=(-a)^mq^{-m(m+1)/2}(q/a;q)_m,
\]
\[
 (aq^{-m}z;q)_m=(-az)^mq^{-m(m+1)/2}(q/(az);q)_m.
\]
Since
\(
 (aq^{-m}z;q)_\infty=(aq^{-m}z;q)_m(az;q)_\infty,
\)
we obtain
\[
 \frac{(q;q)_m}{(q/a;q)_m}
 \frac{(q/(az);q)_m(az;q)_\infty}{(z;q)_\infty}.
\]
On the other hand, substituting $b=q^{m+1}$ in the proposed product and cancelling tails gives exactly the same expression:
\[
 \frac{(q;q)_\infty}{(q^{m+1};q)_\infty}
 \frac{(q^{m+1}/a;q)_\infty}{(q/a;q)_\infty}
 \frac{(q/(az);q)_\infty}{(q^{m+1}/(az);q)_\infty}
 \frac{(az;q)_\infty}{(z;q)_\infty}.
\]

For fixed $a,z$ with $\abs z<1$, both sides are analytic functions of $b$ in a neighborhood of $0$ contained in $\abs b<\abs{az}$.  Absolute convergence is locally uniform there by \Cref{prop:1psi1-convergence}.  For all sufficiently large $m$, the points $b=q^{m+1}$ lie in that neighborhood, and they accumulate at $0$.  The identity theorem therefore proves equality near $0$.  Both sides are meromorphic functions of $b$ throughout the disk $\abs b<\abs{az}$: the series gives the meromorphic continuation termwise on compacta avoiding denominator zeros, and the product quotient is visibly meromorphic.  Uniqueness of meromorphic continuation proves the identity everywhere in that disk where both sides are defined, which is precisely the asserted annulus for fixed $a$ and $z$.
\end{proof}

\begin{remark}
Ramanujan's sum contains the infinite $q$-binomial theorem and several theta-product identities as limiting or specialized cases.  It is a model example of a bilateral identity: the second inequality in the convergence annulus is controlled by the positive tail, while the first is controlled by the negative tail.
\end{remark}


\chapter{\texorpdfstring{$q$}{q}-difference calculus and the quantum binomial}\label{chap:qcalculus}
\index{$q$-derivative}
\index{Jackson integral}
\index{quantum plane}

\section{The \texorpdfstring{$q$}{q}-derivative}

Assume $q\ne1$.  For $x\ne0$ define
\[
 D_qf(x)=\frac{f(x)-f(qx)}{(1-q)x}.
\]
At $x=0$ one uses the continuous extension when it exists.

\begin{proposition}[Elementary rules]\label{prop:qderivative-rules}
For $n\ge1$,
\[
 D_qx^n=[n]_qx^{n-1}.
\]
The product rules are
\begin{align*}
 D_q(fg)(x)&=g(x)D_qf(x)+f(qx)D_qg(x),\\
 &=f(x)D_qg(x)+g(qx)D_qf(x).
\end{align*}
Moreover,
\[
 D_q\bigl(f(q^rx)\bigr)=q^r(D_qf)(q^rx)
\]
for integral $r$.
\end{proposition}

\begin{proof}
The monomial identity follows from
\(
 (x^n-q^nx^n)/((1-q)x)=[n]_qx^{n-1}.
\)
For the first product rule, insert and subtract $f(qx)g(x)$ in
\(
 f(x)g(x)-f(qx)g(qx).
\)
For the second, insert and subtract $f(x)g(qx)$.  The scaling rule follows directly from the definition.
\end{proof}

\begin{theorem}[$q$-Leibniz rule]\label{thm:q-leibniz}
For $n\in\N_0$,
\[
 D_q^n(fg)(x)
 =\sum_{k=0}^{n}\qbinom nk
 (D_q^kf)(q^{n-k}x)(D_q^{n-k}g)(x).
\]
\end{theorem}

\begin{proof}
Induct on $n$.  The assertion is the ordinary product at $n=0$.  Apply $D_q$ to the $k$th term at level $n$.  By the product and scaling rules, it contributes
\[
 q^{n-k}\qbinom nk
 (D_q^{k+1}f)(q^{n-k}x)(D_q^{n-k}g)(x)
\]
to index $k+1$, and
\[
 \qbinom nk(D_q^kf)(q^{n-k+1}x)(D_q^{n-k+1}g)(x)
\]
to index $k$.  The total coefficient of index $k$ is
\[
 \qbinom nk+q^{n-k+1}\qbinom n{k-1}
 =\qbinom{n+1}k
\]
by the second $q$-Pascal recurrence.
\end{proof}

\section{\texorpdfstring{$q$}{q}-Taylor and Newton interpolation}

Define the $q$-falling power based at $a$ by
\[
 (x-a)_q^m=\prod_{j=0}^{m-1}(x-q^ja),
 \qquad (x-a)_q^0=1.
\]

\begin{lemma}\label{lem:qfalling-derivative}
For $m\ge1$,
\[
 D_q(x-a)_q^m=[m]_q(x-a)_q^{m-1}.
\]
Consequently,
\[
 D_q^j(x-a)_q^m
 =\frac{[m]_q!}{[m-j]_q!}(x-a)_q^{m-j}
\]
for $0\le j\le m$.
\end{lemma}

\begin{proof}
Let $P_m(x)=(x-a)_q^m$.  Then
\[
 P_m(qx)=q^m(x-q^{-1}a)\prod_{j=0}^{m-2}(x-q^ja).
\]
Subtracting from
\(
 P_m(x)=(x-q^{m-1}a)\prod_{j=0}^{m-2}(x-q^ja)
\)
gives
\[
 P_m(x)-P_m(qx)=(1-q^m)xP_{m-1}(x).
\]
Divide by $(1-q)x$.  Iteration proves the second formula.
\end{proof}

\begin{theorem}[$q$-Taylor formula]\label{thm:q-taylor}
Assume $[j]_q\ne0$ for $1\le j\le N$.  If $f$ is a polynomial of degree at most $N$, then
\[
 f(x)=\sum_{k=0}^{N}\frac{(D_q^kf)(a)}{[k]_q!}(x-a)_q^k.
\]
\end{theorem}

\begin{proof}
The polynomials $(x-a)_q^k$, $0\le k\le N$, are monic of distinct degrees and hence form a basis.  Write
\(
 f=\sum c_k(x-a)_q^k.
\)
Apply $D_q^j$ and set $x=a$.  By \Cref{lem:qfalling-derivative}, every term with $k>j$ still contains the factor $x-a$, terms with $k<j$ vanish after too many derivatives, and the $k=j$ term equals $c_j[j]_q!$.  Thus $c_j=(D_q^jf)(a)/[j]_q!$.
\end{proof}

\begin{remark}
When the nodes $a,aq,aq^2,\ldots$ are distinct, this is a Newton interpolation formula on that geometric grid.  In every case covered by the theorem, its basis zeros follow a multiplicative rather than an additive progression.
\end{remark}

\section{The quantum plane}

\begin{theorem}[Noncommutative $q$-binomial theorem]\label{thm:quantum-binomial}
Let $X,Y$ belong to an associative algebra and satisfy
\[
 YX=qXY.
\]
Then
\[
 (X+Y)^n=\sum_{k=0}^{n}\qbinom nkX^{n-k}Y^k.
\]
\end{theorem}

\begin{proof}
Induct on $n$ and multiply on the right by $X+Y$.  Since $Y^kX=q^kXY^k$, the coefficient of $X^{n+1-k}Y^k$ becomes
\[
 q^k\qbinom nk+\qbinom n{k-1}=\qbinom{n+1}k
\]
by the first $q$-Pascal recurrence.
\end{proof}

\index{quantum affine space}
\index{$q$-multinomial coefficient!as structure constants}

\begin{corollary}[Noncommutative $q$-multinomial theorem]\label{cor:quantum-multinomial}
Let $r\ge1$ and let $x_1,\dots,x_r$ be elements of an associative algebra over a commutative ring in which $q$ is a central element, and suppose that
\[
 x_jx_i=qx_ix_j
 \qquad(1\le i<j\le r).
\]
Then, for every $n\ge0$,
\[
 (x_1+\cdots+x_r)^n
 =\sum_{\substack{n_1,\dots,n_r\ge0\\ n_1+\cdots+n_r=n}}
 \qmultinom{n}{n_1,\dots,n_r}\,
 x_1^{n_1}x_2^{n_2}\cdots x_r^{n_r},
\]
where on the right the factors are written in increasing order of index.  Equivalently, the $q$-multinomial coefficients are exactly the structure constants that expand powers of $x_1+\cdots+x_r$ in the ordered monomial basis of the quantum affine $r$-space.
\end{corollary}

% ed.: the incoming survey states only that ``iterating the proof gives a
% ed.: $q$-multinomial theorem''.  The ordering of the monomials is part of the
% ed.: statement, not a convention, because $x_jx_i=qx_ix_j$ is imposed only for
% ed.: $i<j$; centrality of $q$ is likewise recorded explicitly here.

\begin{proof}
Induct on $r$.  For $r=1$ there is nothing to prove, and $r=2$ is \Cref{thm:quantum-binomial}.

Let $r\ge2$ and set $X=x_1$ and $Y=x_2+\cdots+x_r$.  Because $x_jx_1=qx_1x_j$ for every $j\ge2$ and $q$ is central,
\[
 YX=\sum_{j=2}^{r}x_jx_1
 =q\sum_{j=2}^{r}x_1x_j
 =qXY,
\]
so \Cref{thm:quantum-binomial} applies to the pair $(X,Y)$ and gives
\[
 (x_1+\cdots+x_r)^n
 =\sum_{k=0}^{n}\qbinom nk\,x_1^{\,n-k}\,(x_2+\cdots+x_r)^{k}.
\]
Substituting $n_1=n-k$ and using the symmetry $\qbinom nk=\qbinom n{n-k}$ of \Cref{thm:qbinom-structure}, this reads
\[
 (x_1+\cdots+x_r)^n
 =\sum_{n_1=0}^{n}\qbinom n{n_1}\,x_1^{\,n_1}\,(x_2+\cdots+x_r)^{\,n-n_1}.
\]
The elements $x_2,\dots,x_r$ satisfy the same commutation relations among themselves, so the inductive hypothesis applies to them and gives, with $m=n-n_1$,
\[
 (x_2+\cdots+x_r)^{m}
 =\sum_{\substack{n_2,\dots,n_r\ge0\\ n_2+\cdots+n_r=m}}
 \qmultinom{m}{n_2,\dots,n_r}\,x_2^{n_2}\cdots x_r^{n_r}.
\]
Since the coefficients are central, substitution produces precisely the ordered monomials $x_1^{n_1}x_2^{n_2}\cdots x_r^{n_r}$, the coefficient of $x_1^{n_1}\cdots x_r^{n_r}$ being
\[
 \qbinom n{n_1}\qmultinom{n-n_1}{n_2,\dots,n_r}
 =\frac{\poch{q}{q}{n}}{\poch{q}{q}{n_1}\poch{q}{q}{n-n_1}}
 \cdot
 \frac{\poch{q}{q}{n-n_1}}{\poch{q}{q}{n_2}\cdots\poch{q}{q}{n_r}}
 =\qmultinom{n}{n_1,\dots,n_r}
\]
by the definition of the $q$-multinomial coefficient preceding \Cref{prop:qmult-factor}; the factor $\poch{q}{q}{n-n_1}$ telescopes away, exactly as in the proof of that proposition.  No factor was transposed at any stage after the expansion, so the monomials are produced in increasing order of index, as asserted.
\end{proof}

\begin{remark}[The same coefficient read combinatorially]\label{rem:quantum-multinomial-inversions}
Expanding $(x_1+\cdots+x_r)^n$ without using any relation produces the $r^n$ words $x_{w_1}x_{w_2}\cdots x_{w_n}$ with $w\in\{1,\dots,r\}^n$.  Sorting such a word into increasing order uses one adjacent transposition for each inversion of $w$, and each transposition contributes exactly one factor $q$; hence
\[
 x_{w_1}x_{w_2}\cdots x_{w_n}=q^{\inv(w)}x_1^{n_1}\cdots x_r^{n_r},
\]
where $n_i$ is the multiplicity of the letter $i$ in $w$.  Comparing coefficients with \Cref{cor:quantum-multinomial} recovers the multiset inversion enumerator of \Cref{thm:qmult-inversions}, and conversely that theorem proves \Cref{cor:quantum-multinomial} again.  The quantum affine space is therefore a purely algebraic encoding of the inversion statistic, just as \Cref{thm:quantum-binomial} encodes \Cref{thm:binary-inversions}.
\end{remark}

\begin{corollary}[Factorization of $q$-exponentials]\label{cor:qexp-factor}
Assume $\abs q<1$.  Define
\[
 e_q(x)=\sum_{n=0}^{\infty}\frac{x^n}{[n]_q!}
 =\frac1{((1-q)x;q)_\infty},
 \qquad \abs{(1-q)x}<1,
\]
and
\[
 E_q(x)=\sum_{n=0}^{\infty}\frac{q^{\binom n2}x^n}{[n]_q!}
 =(-(1-q)x;q)_\infty,
 \qquad x\in\C.
\]
Then $e_q(x)E_q(-x)=1$ wherever the first series is defined.  If $YX=qXY$, then, as an identity of formal series in total degree,
\[
 e_q(X+Y)=e_q(X)e_q(Y).
\]
\end{corollary}

\begin{proof}
The product formulas and the stated convergence domains follow from Euler's identities after replacing $z$ by $(1-q)x$.  Their product is $1$.  For the noncommutative factorization, apply \Cref{thm:quantum-binomial} degree by degree and use
\[
 \frac1{[n]_q!}\qbinom nk
 =\frac1{[n-k]_q![k]_q!}.
\]
Reindex with $r=n-k$ to obtain the Cauchy product of $e_q(X)$ and $e_q(Y)$, with all $X$ factors preceding all $Y$ factors.  Since every total degree receives only finitely many contributions, the formal rearrangement is valid.
\end{proof}

\begin{proposition}[Eigenfunction properties]\label{prop:qexp-eigen}
Under the hypotheses and within the convergence domains of \Cref{cor:qexp-factor}, for scalar $\lambda$,
\[
 D_qe_q(\lambda x)=\lambda e_q(\lambda x),
 \qquad
 D_qE_q(\lambda x)=\lambda E_q(q\lambda x).
\]
\end{proposition}

\begin{proof}
Differentiate the series term by term and use $D_qx^n=[n]_qx^{n-1}$.
\end{proof}

\section{Jackson's \texorpdfstring{$q$}{q}-integral}

Assume $0<q<1$.  Define
\[
 \int_0^a f(t)\,d_qt
 =(1-q)a\sum_{n=0}^{\infty}q^nf(aq^n),
\]
whenever the series converges.

\begin{theorem}[Fundamental theorem of $q$-calculus]\label{thm:q-fundamental}
Let
\[
 F(x)=\int_0^x f(t)\,d_qt.
\]
Whenever the defining series converges locally uniformly,
\[
 D_qF(x)=f(x).
\]
Conversely, if $G(aq^n)\to G(0)$, then
\[
 \int_0^aD_qG(t)\,d_qt=G(a)-G(0).
\]
\end{theorem}

\begin{proof}
From the definition,
\[
 F(x)=(1-q)x\sum_{n\ge0}q^nf(xq^n),
\]
while
\[
 F(qx)=(1-q)x\sum_{n\ge1}q^nf(xq^n).
\]
Their difference is $(1-q)xf(x)$, proving the first identity.  For the second,
\begin{align*}
 \int_0^aD_qG(t)\,d_qt
 &=(1-q)a\sum_{n\ge0}q^n
 \frac{G(aq^n)-G(aq^{n+1})}{(1-q)aq^n}\\
 &=\sum_{n\ge0}\bigl(G(aq^n)-G(aq^{n+1})\bigr),
\end{align*}
which telescopes to $G(a)-G(0)$.
\end{proof}

\begin{corollary}[Jackson integration by parts]\label{cor:q-integration-parts}
Under convergence assumptions,
\[
 \int_0^a f(t)D_qg(t)\,d_qt
 =f(a)g(a)-f(0)g(0)
 -\int_0^a g(qt)D_qf(t)\,d_qt.
\]
\end{corollary}

\begin{proof}
Use the product rule
\(
 D_q(fg)=fD_qg+g(q\,\cdot)D_qf
\)
and integrate with \Cref{thm:q-fundamental}.
\end{proof}

\begin{example}[Monomials]
For $\Re(s)>-1$,
\[
 \int_0^a t^s\,d_qt
 =(1-q)a^{s+1}\sum_{n\ge0}q^{n(s+1)}
 =\frac{a^{s+1}}{[s+1]_q},
\]
where $[u]_q=(1-q^u)/(1-q)$ for complex $u$.
\end{example}

\chapter{The \texorpdfstring{$q$}{q}-gamma and \texorpdfstring{$q$}{q}-beta functions}\label{chap:qgamma}
\index{$q$-gamma function}
\index{$q$-beta function}

Throughout this chapter $0<q<1$ unless stated otherwise.

\section{Definition and functional equation}

\begin{definition}
For $x>0$, define
\[
 \Gamma_q(x)
 =(1-q)^{1-x}\frac{(q;q)_\infty}{(q^x;q)_\infty}.
\]
With $q^z=\exp(z\log q)$, the same product quotient defines a
meromorphic function of $z\in\C$; its possible poles occur where
$q^{z+j}=1$ for some $j\in\N_0$.  Whenever complex arguments of
$\Gamma_q$ occur below, this meromorphic continuation is intended.
\end{definition}

\begin{theorem}[Basic properties]\label{thm:qgamma-basic}
For $x>0$,
\[
 \Gamma_q(1)=1,
 \qquad
 \Gamma_q(x+1)=[x]_q\Gamma_q(x).
\]
Consequently, for $n\in\N$,
\[
 \Gamma_q(n)=[n-1]_q!.
\]
The recurrence extends as an identity of meromorphic functions of a
complex argument.
\end{theorem}

\begin{proof}
The value at $1$ is immediate.  The tail identity
\(
 (q^x;q)_\infty=(1-q^x)(q^{x+1};q)_\infty
\)
gives
\[
 \frac{\Gamma_q(x+1)}{\Gamma_q(x)}
 =(1-q)^{-1}(1-q^x)=[x]_q.
\]
Iteration gives the integral values.  The same tail identity is valid for
complex $x$, and hence proves the meromorphic recurrence wherever both
sides are finite; equality elsewhere follows in the sense of
meromorphic continuation.
\end{proof}

\begin{proposition}[Logarithmic derivatives and convexity]\label{prop:qgamma-convex}
For $x>0$,
\[
 \frac{\dd}{\dd x}\log\Gamma_q(x)
 =-\log(1-q)+(\log q)\sum_{n=0}^{\infty}
 \frac{q^{n+x}}{1-q^{n+x}},
\]
\[
 \frac{\dd^2}{\dd x^2}\log\Gamma_q(x)
 = (\log q)^2\sum_{n=0}^{\infty}
 \frac{q^{n+x}}{(1-q^{n+x})^2}>0.
\]
Thus $\Gamma_q$ is strictly log-convex.
\end{proposition}

\begin{proof}
Differentiate the locally uniformly convergent logarithm of the defining infinite product.  A second differentiation gives the positive series.
\end{proof}

\section{A \texorpdfstring{$q$}{q}-Bohr--Mollerup theorem}

\begin{theorem}[$q$-Bohr--Mollerup characterization]\label{thm:q-bohr-mollerup}
The function $\Gamma_q$ is the unique function $f:(0,\infty)\to(0,\infty)$ satisfying
\[
 f(1)=1,
 \qquad
 f(x+1)=[x]_qf(x),
\]
and such that $\log f$ is convex.
\end{theorem}

\begin{proof}
Existence follows from \Cref{thm:qgamma-basic,prop:qgamma-convex}.  For uniqueness let $f$ have the stated properties and fix $0<x\le1$.  Convexity on the interval from $n$ to $n+1$ gives
\[
 f(n+x)\le f(n)^{1-x}f(n+1)^x.
\]
A comparison of the adjacent secant slopes on $[n+x,n+1]$ and $[n+1,n+2]$ gives
\[
 f(n+x)\ge\frac{f(n+1)}{[n+1]_q^{1-x}}.
\]
Using the recurrence,
\[
 f(n+x)=f(x)\prod_{j=0}^{n-1}[x+j]_q,
 \quad
 f(n)=[n-1]_q!,
 \quad
 f(n+1)=[n]_q!,
\]
we obtain
\[
 \frac{[n]_q!}
 {[n+1]_q^{1-x}\prod_{j=0}^{n-1}[x+j]_q}
 \le f(x)\le
 \frac{[n-1]_q![n]_q^x}
 {\prod_{j=0}^{n-1}[x+j]_q}.
\]
The ratio of the upper bound to the lower bound is
\[
 \left(\frac{[n+1]_q}{[n]_q}\right)^{1-x}\longrightarrow1.
\]
Hence the bounds determine $f(x)$ uniquely for $0<x\le1$, and the recurrence determines it on all of $(0,\infty)$.  Since $\Gamma_q$ satisfies the same conditions, $f=\Gamma_q$.
\end{proof}

\section{Classical limit}

\begin{theorem}[Limit to Euler's gamma function]\label{thm:qgamma-classical-limit}
For every $x>0$,
\[
 \lim_{q\to1^-}\Gamma_q(x)=\Gamma(x).
\]
The convergence is locally uniform on $(0,\infty)$.
\end{theorem}

\begin{proof}
It suffices by the functional equations to consider $0<x\le1$.  The bounds in the proof of \Cref{thm:q-bohr-mollerup} hold for $\Gamma_q(x)$.  Fix $n$ and let $q\to1^-$.  Since $[u]_q\to u$, they become
\[
 \frac{n!}{(n+1)^{1-x}x(x+1)\cdots(x+n-1)}
 \le \liminf_{q\to1^-}\Gamma_q(x)
\]
\[
 \le\limsup_{q\to1^-}\Gamma_q(x)
 \le\frac{(n-1)!n^x}{x(x+1)\cdots(x+n-1)}.
\]
Both outer expressions tend to $\Gamma(x)$ as $n\to\infty$.  To see this directly, use Euler's beta integral:
\[
 B(x,n)=\int_0^1t^{x-1}(1-t)^{n-1}\,dt
 =\frac{(n-1)!}{x(x+1)\cdots(x+n-1)}.
\]
After $t=u/n$, dominated convergence gives $n^xB(x,n)\to\int_0^\infty u^{x-1}\e^{-u}\,du=\Gamma(x)$.  The lower expression differs from the upper by a factor tending to $1$.  This proves pointwise convergence.  We record the uniformity argument.  On a compact interval $x\in[\delta,1]$, the upper outer bound can be written as
\[
 n^xB(x,n)=\int_0^n u^{x-1}\left(1-\frac un\right)^{n-1}\dd u.
\]
For $n\ge2$, its integrand is dominated uniformly in $x$ by $u^{\delta-1}$ on $(0,1]$ and by $e^{-u/2}$ on $[1,\infty)$.  More explicitly, after extending the kernel by $0$ for $u>n$,
\[
 \sup_{\delta\le x\le1}
 \left|u^{x-1}\left(1-\frac un\right)_{+}^{n-1}
       -u^{x-1}e^{-u}\right|
 \longrightarrow0
\]
for every $u>0$, while this supremum is bounded by an integrable
multiple of $u^{\delta-1}\mathbf 1_{(0,1]}(u)+e^{-u/2}\mathbf 1_{[1,\infty)}(u)$.
The dominated convergence theorem therefore gives
\[
 \sup_{\delta\le x\le1}
 \left|n^xB(x,n)-\Gamma(x)\right|\longrightarrow0.
\]
The lower outer bound differs by the factor $(n/(n+1))^{1-x}$, which also tends uniformly to $1$.  For fixed $n$, the original $q$-bounds converge uniformly as $q\to1^-$ because they contain only finitely many continuous factors.  A uniform squeeze proves local uniform convergence on $(0,1]$, and the recurrence extends it to arbitrary compact subsets of $(0,\infty)$.
\end{proof}

\section{The Jackson \texorpdfstring{$q$}{q}-beta integral}

\begin{definition}
For $x,y>0$, define
\[
 B_q(x,y)=\int_0^1 t^{x-1}
 \frac{(qt;q)_\infty}{(q^yt;q)_\infty}\,d_qt.
\]
\end{definition}

\begin{theorem}[$q$-beta evaluation]\label{thm:qbeta}
For $x,y>0$,
\[
 B_q(x,y)=\frac{\Gamma_q(x)\Gamma_q(y)}{\Gamma_q(x+y)}.
\]
\end{theorem}

\begin{proof}
Expand the Jackson integral:
\begin{align*}
 B_q(x,y)
 &=(1-q)\sum_{n\ge0}q^{nx}
 \frac{(q^{n+1};q)_\infty}{(q^{n+y};q)_\infty}\\
 &=(1-q)\frac{(q;q)_\infty}{(q^y;q)_\infty}
 \sum_{n\ge0}\frac{(q^y;q)_n}{(q;q)_n}q^{nx}.
\end{align*}
The infinite $q$-binomial theorem evaluates the sum as
\(
 (q^{x+y};q)_\infty/(q^x;q)_\infty.
\)
Thus
\[
 B_q(x,y)
 =(1-q)\frac{(q;q)_\infty(q^{x+y};q)_\infty}
 {(q^x;q)_\infty(q^y;q)_\infty},
\]
which is exactly the gamma quotient after substituting the definition of $\Gamma_q$.
\end{proof}

\begin{corollary}[Recurrences]\label{cor:qbeta-recurrence}
For $x,y>0$,
\[
 B_q(x+1,y)=\frac{[x]_q}{[x+y]_q}B_q(x,y),
 \qquad
 B_q(x,y+1)=\frac{[y]_q}{[x+y]_q}B_q(x,y).
\]
\end{corollary}

\begin{proof}
Apply the functional equation for $\Gamma_q$ to the quotient in \Cref{thm:qbeta}.
\end{proof}

\section{Multiplication and reflection-type formulas}

\begin{theorem}[$q$-multiplication formula]\label{thm:qgamma-multiplication}
For an integer $m\ge1$ and $x>0$,
\[
 \Gamma_q(mx)\prod_{j=1}^{m-1}\Gamma_{q^m}\!\left(\frac jm\right)
 =[m]_q^{mx-1}
 \prod_{j=0}^{m-1}\Gamma_{q^m}\!\left(x+\frac jm\right).
\]
\end{theorem}

\begin{proof}
Multiply the definitions on the right.  By dissection,
\[
 \prod_{j=0}^{m-1}(q^{mx+j};q^m)_\infty=(q^{mx};q)_\infty.
\]
Hence
\[
 \prod_{j=0}^{m-1}\Gamma_{q^m}\!\left(x+\frac jm\right)
 =(1-q^m)^{(m+1)/2-mx}
 \frac{(q^m;q^m)_\infty^m}{(q^{mx};q)_\infty}.
\]
Similarly,
\[
 \prod_{j=1}^{m-1}\Gamma_{q^m}\!\left(\frac jm\right)
 =(1-q^m)^{(m-1)/2}
 \frac{(q^m;q^m)_\infty^m}{(q;q)_\infty}.
\]
Taking the ratio and using
\(
 [m]_q=(1-q^m)/(1-q)
\)
produces the formula.
\end{proof}

\begin{proposition}[Reflection-type product]\label{prop:qgamma-reflection}
For $0<x<1$,
\[
 \Gamma_q(x)\Gamma_q(1-x)
 =(1-q)\frac{(q;q)_\infty^2}
 {(q^x,q^{1-x};q)_\infty}.
\]
\end{proposition}

\begin{proof}
Multiply the two defining products; the powers of $1-q$ combine to one.
\end{proof}

\begin{remark}
Unlike Euler's reflection formula, the right side is naturally a theta quotient rather than an elementary trigonometric function.  Jacobi's triple product in the next part makes that theta structure explicit.
\end{remark}


\part{Theta products and partition identities}

\chapter{Jacobi's triple product}\label{chap:jtp}
\index{Jacobi triple product}
\index{theta function}

\section{A bilateral theta series}

For $z\ne0$ and $\abs q<1$, define
\[
 \Theta(z;q)=\sum_{n\in\Z}q^{\binom n2}z^n.
\]
The series converges locally uniformly on $\C^*$ because its terms decay quadratically in both directions.

\begin{theorem}[Jacobi triple product]\label{thm:jtp}
For $z\ne0$ and $\abs q<1$,
\begin{align}
 \sum_{n\in\Z}q^{\binom n2}z^n
 &=(-z,-q/z,q;q)_\infty,\label{eq:jtp-plus}\\
 \sum_{n\in\Z}(-1)^nq^{\binom n2}z^n
 &=(z,q/z,q;q)_\infty.\label{eq:jtp-minus}
\end{align}
\end{theorem}

\begin{proof}
It suffices to prove \eqref{eq:jtp-plus}; replace $z$ by $-z$ for \eqref{eq:jtp-minus}.  Put
\[
 F(z)=(-z;q)_\infty(-q/z;q)_\infty(q;q)_\infty.
\]
The products converge locally uniformly on $\C^*$, so $F$ has a Laurent expansion
\(
 F(z)=\sum_{n\in\Z}c_nz^n.
\)
The shift identities give
\[
 F(qz)=\frac{1+z^{-1}}{1+z}F(z)=z^{-1}F(z).
\]
Comparing Laurent coefficients yields
\[
 c_{n+1}=q^nc_n,
 \qquad
 c_n=q^{\binom n2}c_0
\]
for every integer $n$.

It remains to find $c_0$.  Euler's product expansions give
\[
 (-z;q)_\infty=\sum_{r\ge0}\frac{q^{\binom r2}z^r}{(q;q)_r},
 \qquad
 (-q/z;q)_\infty=\sum_{s\ge0}\frac{q^{\binom{s+1}{2}}z^{-s}}{(q;q)_s}.
\]
Hence
\[
 c_0=(q;q)_\infty\sum_{r\ge0}\frac{q^{r^2}}{(q;q)_r^2}.
\]
The sum is the generating function obtained by decomposing an arbitrary partition according to the side $r$ of its Durfee square: the square contributes $q^{r^2}$, while the diagrams to its right and below are independent partitions with at most $r$ parts, each having generating function $1/(q;q)_r$.  Therefore
\[
 \sum_{r\ge0}\frac{q^{r^2}}{(q;q)_r^2}
 =\frac1{(q;q)_\infty},
\]
so $c_0=1$.  Thus
\(
 F(z)=\sum_{n\in\Z}q^{\binom n2}z^n.
\)
\end{proof}

\begin{corollary}[Quasi-periodicity and zeros]\label{cor:theta-quasi}
The theta function satisfies
\[
 \Theta(qz;q)=z^{-1}\Theta(z;q).
\]
Its zeros are the simple points $z=-q^m$, $m\in\Z$.
\end{corollary}

\begin{proof}
The functional equation follows either from reindexing the bilateral series or from the proof above.  The product in \eqref{eq:jtp-plus} vanishes exactly when $1+zq^j=0$ or $1+q^{j+1}/z=0$, which together give $z=-q^m$ for $m\in\Z$.  Exactly one factor vanishes at each such point.
\end{proof}

\section{Euler's pentagonal number theorem}

\begin{theorem}[Pentagonal number theorem]\label{thm:pentagonal}
For $\abs q<1$,
\[
 (q;q)_\infty
 =\sum_{n\in\Z}(-1)^nq^{n(3n-1)/2}
 =1+\sum_{n=1}^{\infty}(-1)^n
 \left(q^{n(3n-1)/2}+q^{n(3n+1)/2}\right).
\]
\end{theorem}

\begin{proof}
In \eqref{eq:jtp-minus}, replace the base $q$ by $q^3$ and set $z=q$.  The product side is
\[
 (q;q^3)_\infty(q^2;q^3)_\infty(q^3;q^3)_\infty=(q;q)_\infty
\]
by dissection.  The series exponent is
\[
 3\binom n2+n=\frac{n(3n-1)}2.
\]
Pair the terms $n$ and $-n$ to obtain the second form.
\end{proof}

\begin{corollary}[Euler's partition recurrence]\label{cor:euler-partition-recurrence}
Let $p(0)=1$ and $p(n)=0$ for $n<0$.  Then for $n\ge1$,
\[
 p(n)=\sum_{j=1}^{\infty}(-1)^{j-1}
 \left(p\!\left(n-\frac{j(3j-1)}2\right)
 +p\!\left(n-\frac{j(3j+1)}2\right)\right).
\]
\end{corollary}

\begin{proof}
The generating function for partitions is
\(
 \sum_{n\ge0}p(n)q^n=1/(q;q)_\infty.
\)
Multiply this by the pentagonal expansion and compare the coefficient of $q^n$.
\end{proof}

\section{Jacobi's cubic identity}

\begin{theorem}[Jacobi identity]\label{thm:jacobi-cubic}
For $\abs q<1$,
\[
 (q;q)_\infty^3
 =\sum_{n=0}^{\infty}(-1)^n(2n+1)q^{n(n+1)/2}.
\]
\end{theorem}

\begin{proof}
Differentiate \eqref{eq:jtp-minus} with respect to $z$ and set $z=1$.  Since
\[
 (z;q)_\infty=(1-z)(qz;q)_\infty,
\]
the derivative of the product side at $z=1$ is $-(q;q)_\infty^3$.  On the series side it is
\[
 \sum_{m\in\Z}m(-1)^mq^{\binom m2}.
\]
Pair $m=-n$ with $m=n+1$ for $n\ge0$.  Their sum is
\[
 -(2n+1)(-1)^nq^{n(n+1)/2}.
\]
Cancelling the common minus sign proves the identity.
\end{proof}

\section{A theta form of the \texorpdfstring{$q$}{q}-gamma reflection product}

\begin{corollary}\label{cor:qgamma-theta}
For $0<x<1$,
\[
 \Gamma_q(x)\Gamma_q(1-x)
 =(1-q)\frac{(q;q)_\infty^3}
 {\displaystyle\sum_{n\in\Z}(-1)^n
 q^{\binom n2+xn}}.
\]
\end{corollary}

\begin{proof}
Set $z=q^x$ in \eqref{eq:jtp-minus}:
\[
 \sum_{n\in\Z}(-1)^nq^{\binom n2+xn}
 =(q^x,q^{1-x},q;q)_\infty.
\]
Substitute this in \Cref{prop:qgamma-reflection}.
\end{proof}

\chapter{Partitions and product--sum identities}\label{chap:partitions}
\index{partition generating function}

\section{Unrestricted, distinct, and odd parts}

\begin{theorem}[Basic partition products]\label{thm:partition-products}
Let $p(n)$ count unrestricted partitions, $d(n)$ partitions into distinct parts, and $o(n)$ partitions into odd parts.  Then
\begin{align*}
 \sum_{n\ge0}p(n)q^n&=\frac1{(q;q)_\infty},\\
 \sum_{n\ge0}d(n)q^n&=(-q;q)_\infty,\\
 \sum_{n\ge0}o(n)q^n&=\frac1{(q;q^2)_\infty}.
\end{align*}
Moreover, $d(n)=o(n)$ for every $n$.
\end{theorem}

\begin{proof}
For unrestricted partitions, the part $j$ may occur $0,1,2,\dots$ times and contributes the geometric factor $(1-q^j)^{-1}$.  For distinct partitions, the part $j$ occurs zero or one time and contributes $1+q^j$.  For odd parts, retain only the odd geometric factors.  Finally,
\[
 (-q;q)_\infty
 =\prod_{j\ge1}\frac{1-q^{2j}}{1-q^j}
 =\frac{(q^2;q^2)_\infty}{(q;q)_\infty}
 =\frac1{(q;q^2)_\infty},
\]
so the two generating functions agree coefficientwise.
\end{proof}

\begin{remark}[A bijection]
The equality $d(n)=o(n)$ also has a direct proof.  In a partition into distinct parts, write each part uniquely as $2^ru$ with $u$ odd.  Replace that part by $2^r$ copies of $u$.  Distinct powers of two ensure that the process is reversible by binary expansion of the multiplicity of every odd part.
\end{remark}

\section{Bounded partitions}

\begin{proposition}\label{prop:bounded-partitions}
The following generating functions are equal:
\[
 \frac1{(q;q)_r}
 =\sum_{n\ge0}p_{\le r}(n)q^n,
\]
where $p_{\le r}(n)$ counts partitions with largest part at most $r$, or equivalently partitions with at most $r$ parts.
\end{proposition}

\begin{proof}
The product allows parts $1,\dots,r$ with arbitrary multiplicity.  Conjugation of Ferrers diagrams exchanges the condition ``largest part at most $r$'' with ``at most $r$ parts.''
\end{proof}

\begin{corollary}[Durfee decomposition of all partitions]\label{cor:durfee-all}
\[
 \frac1{(q;q)_\infty}
 =\sum_{r=0}^{\infty}\frac{q^{r^2}}{(q;q)_r^2}.
\]
\end{corollary}

\begin{proof}
Decompose every partition by the side $r$ of its Durfee square.  The diagrams to the right and below are independent partitions with at most $r$ parts.  This is the identity used in the proof of Jacobi's triple product.
\end{proof}

\section{Finite approximations and stabilization}

\begin{proposition}\label{prop:partition-stabilization}
For fixed $N$, the coefficient of $q^N$ in $\qbinom{m+k}{k}$ equals $p(N)$ whenever $m,k\ge N$.
\end{proposition}

\begin{proof}
The Gaussian coefficient counts partitions inside a $k\times m$ rectangle.  Any partition of $N$ has at most $N$ parts and largest part at most $N$, so for $m,k\ge N$ the rectangle restriction is vacuous.
\end{proof}

\begin{remark}
Thus $\qbinom{m+k}{k}$ converges coefficientwise to $1/(q;q)_\infty$ as both sides of the rectangle tend to infinity.  This finite-to-infinite passage is a recurring bridge between Gaussian polynomials and partition identities.
\end{remark}

\chapter{Bailey pairs and the Rogers--Ramanujan identities}\label{chap:bailey}
\index{Bailey pair}
\index{Rogers--Ramanujan identities}

\section{Bailey pairs}

For brevity, write $(a)_n=(a;q)_n$ in this chapter.

\begin{definition}
Sequences $(\alpha_n)_{n\ge0}$ and $(\beta_n)_{n\ge0}$ form a \emph{Bailey pair relative to $a$} if
\[
 \beta_n=\sum_{r=0}^{n}\frac{\alpha_r}{(q)_{n-r}(aq)_{n+r}}
 \qquad(n\ge0).
\]
\end{definition}

\section{Bailey inversion}

\begin{lemma}[$q$-difference annihilation]\label{lem:qdiff-annihilation}
If $P(x)$ is a polynomial of degree less than $N$, then
\[
 \sum_{s=0}^{N}(-1)^{N-s}q^{\binom{N-s}{2}}
 \qbinom Ns P(q^s)=0.
\]
\end{lemma}

\begin{proof}
By linearity it suffices to take $P(x)=x^m$ with $0\le m<N$.  Put $t=N-s$ and use the finite $q$-binomial theorem:
\begin{align*}
 \sum_{s=0}^{N}(-1)^{N-s}q^{\binom{N-s}{2}}\qbinom Ns q^{ms}
 &=q^{mN}\sum_{t=0}^{N}(-1)^tq^{\binom t2}\qbinom Ntq^{-mt}\\
 &=q^{mN}(q^{-m};q)_N=0,
\end{align*}
for the product contains the factor $1-q^0$.
\end{proof}

\begin{theorem}[Bailey inversion]\label{thm:bailey-inversion}
Assume $a\ne1$ initially.  The Bailey-pair relation is equivalent to
\[
 \alpha_n=\frac{1-aq^{2n}}{1-a}
 \sum_{j=0}^{n}
 \frac{(a)_{n+j}}{(q)_{n-j}}
 (-1)^{n-j}q^{\binom{n-j}{2}}\beta_j.
\]
At $a=1$, the right side is interpreted by its removable limit.
\end{theorem}

\begin{proof}
Substitute the definition of $\beta_j$ into the proposed inverse.  The coefficient of $\alpha_r$ is
\[
 \frac{1-aq^{2n}}{1-a}
 \sum_{j=r}^{n}
 \frac{(a)_{n+j}}{(q)_{n-j}(q)_{j-r}(aq)_{j+r}}
 (-1)^{n-j}q^{\binom{n-j}{2}}.
\]
If $n=r$, only $j=n$ occurs, and
\[
 \frac{1-aq^{2n}}{1-a}\frac{(a)_{2n}}{(aq)_{2n}}=1.
\]
Suppose $N=n-r>0$ and put $j=r+s$.  Then
\[
 \frac{(a)_{n+r+s}}{(aq)_{2r+s}}
 =(1-a)(aq^{2r+s+1};q)_{N-1}.
\]
The coefficient becomes
\[
 \frac{1-aq^{2n}}{(q)_N}
 \sum_{s=0}^{N}(-1)^{N-s}q^{\binom{N-s}{2}}
 \qbinom Ns P(q^s),
\]
where
\(
 P(x)=(aq^{2r+1}x;q)_{N-1}
\)
has degree $N-1$.  It vanishes by \Cref{lem:qdiff-annihilation}.  Thus the two triangular matrices are inverse.  Since every expression is rational in $a$ and the singularity at $a=1$ is removable, the limiting statement follows.
\end{proof}

\begin{corollary}[Two unit Bailey pairs]\label{cor:unit-bailey-pairs}
Let $\beta_n=\delta_{n,0}$.
\begin{enumerate}[label=\textup{(\roman*)}]
\item Relative to $a=1$,
\[
 \alpha_0=1,
 \qquad
 \alpha_n=(-1)^nq^{\binom n2}(1+q^n)\quad(n>0).
\]
\item Relative to $a=q$,
\[
 \alpha_n=(-1)^nq^{\binom n2}[2n+1]_q\quad(n\ge0).
\]
\end{enumerate}
\end{corollary}

\begin{proof}
With $\beta_j=\delta_{j,0}$, Bailey inversion gives
\[
 \alpha_n=\frac{1-aq^{2n}}{1-a}
 \frac{(a)_n}{(q)_n}(-1)^nq^{\binom n2}.
\]
For $a\to1$ and $n>0$, $(a)_n/(1-a)\to(q)_{n-1}$, so the prefactor reduces to
\(
 (1-q^{2n})/(1-q^n)=1+q^n.
\)
For $a=q$, the ratio $(a)_n/(q)_n$ is $1$ and the remaining factor is $[2n+1]_q$.
\end{proof}

\section{The limiting Bailey lemma}

\begin{lemma}[Auxiliary finite identity]\label{lem:bailey-aux}
For $N\ge0$,
\[
 \sum_{k=0}^{N}z^kq^{k(k-1)}\qbinom Nk
 (zq^k;q)_{N-k}=1.
\]
\end{lemma}

\begin{proof}
Call the left side $S_N(z)$.  Using the second $q$-Pascal recurrence and splitting the last factor of $(zq^k;q)_{N+1-k}$ gives
\[
 S_{N+1}(z)=(1-zq^N)S_N(z)+zq^NS_N(qz).
\]
The identity is immediate for $N=0$.  If $S_N$ is identically $1$, the recurrence gives $S_{N+1}=1-zq^N+zq^N=1$.  Induction completes the proof.
\end{proof}

\begin{theorem}[Limiting Bailey lemma, finite form]\label{thm:bailey-limit-finite}
If $(\alpha_n,\beta_n)$ is a Bailey pair relative to $a$, define
\[
 \alpha_n'=a^nq^{n^2}\alpha_n,
 \qquad
 \beta_n'=\sum_{j=0}^{n}\frac{a^jq^{j^2}}{(q)_{n-j}}\beta_j.
\]
Then $(\alpha_n',\beta_n')$ is also a Bailey pair relative to $a$.
\end{theorem}

\begin{proof}
Substitute the Bailey relation for $\beta_j$ and interchange the finite sums.  The coefficient of $\alpha_r$ in $\beta_n'$ is
\[
 a^rq^{r^2}\sum_{s=0}^{N}
 \frac{a^sq^{s^2+2rs}}
 {(q)_{N-s}(q)_s(aq)_{2r+s}},
 \qquad N=n-r.
\]
Let $A=aq^{2r+1}$.  Since
\[
 a^sq^{s^2+2rs}=A^sq^{s(s-1)},
 \qquad
 (aq)_{2r+s}=(aq)_{2r}(A)_s,
\]
the desired Bailey coefficient
\(
 a^rq^{r^2}/((q)_N(aq)_{n+r})
\)
is equivalent, after clearing factors, to
\[
 \sum_{s=0}^{N}\qbinom Ns A^sq^{s(s-1)}(Aq^s)_{N-s}=1.
\]
This is \Cref{lem:bailey-aux}.
\end{proof}

\begin{corollary}[Limiting Bailey lemma, infinite form]\label{cor:bailey-limit-infinite}
If the series converge absolutely, then every Bailey pair relative to $a$ satisfies
\[
 \sum_{n=0}^{\infty}a^nq^{n^2}\beta_n
 =\frac1{(aq;q)_\infty}
 \sum_{n=0}^{\infty}a^nq^{n^2}\alpha_n.
\]
\end{corollary}

\begin{proof}
The transformed pair in \Cref{thm:bailey-limit-finite} satisfies
\[
 \beta_N'=\sum_{r=0}^{N}
 \frac{\alpha_r'}{(q)_{N-r}(aq)_{N+r}}.
\]
Let $N\to\infty$.  On the left,
\[
 \beta_N'\longrightarrow\frac1{(q)_\infty}
 \sum_{j\ge0}a^jq^{j^2}\beta_j.
\]
On the right,
\[
 \beta_N'\longrightarrow
 \frac1{(q)_\infty(aq)_\infty}
 \sum_{r\ge0}a^rq^{r^2}\alpha_r.
\]
Absolute convergence permits dominated convergence.  Cancel $(q)_\infty^{-1}$.
\end{proof}

\section{The Rogers--Ramanujan identities}

\begin{theorem}[Rogers--Ramanujan]\label{thm:rogers-ramanujan}
For $\abs q<1$,
\begin{align}
 \sum_{n=0}^{\infty}\frac{q^{n^2}}{(q;q)_n}
 &=\frac1{(q;q^5)_\infty(q^4;q^5)_\infty},\label{eq:RR1}\\
 \sum_{n=0}^{\infty}\frac{q^{n(n+1)}}{(q;q)_n}
 &=\frac1{(q^2;q^5)_\infty(q^3;q^5)_\infty}.
 \label{eq:RR2}
\end{align}
\end{theorem}

\begin{proof}
For the first identity, begin with the unit pair relative to $a=1$ in \Cref{cor:unit-bailey-pairs}.  Apply the finite limiting lemma once.  The new pair is
\[
 \alpha_n'=q^{n^2}\alpha_n,
 \qquad
 \beta_n'=\frac1{(q)_n}.
\]  Applying the infinite limiting lemma to this new pair gives
\begin{align*}
 \sum_{n\ge0}\frac{q^{n^2}}{(q)_n}
 &=\frac1{(q)_\infty}
 \left(1+\sum_{r\ge1}(-1)^r
 q^{2r^2+\binom r2}(1+q^r)\right)\\
 &=\frac1{(q)_\infty}
 \sum_{r\in\Z}(-1)^rq^{(5r^2-r)/2}.
\end{align*}
By Jacobi's triple product with base $q^5$ and $z=q^2$,
\[
 \sum_{r\in\Z}(-1)^rq^{(5r^2-r)/2}
 =(q^2,q^3,q^5;q^5)_\infty.
\]
Dissection gives
\[
 (q;q)_\infty=(q,q^2,q^3,q^4,q^5;q^5)_\infty,
\]
which proves \eqref{eq:RR1}.

For the second identity use the unit pair relative to $a=q$.  After one finite limiting step,
\[
 \alpha_n'=(-1)^nq^{n^2+n+\binom n2}[2n+1]_q,
 \qquad
 \beta_n'=\frac1{(q)_n}.
\]
The infinite limiting lemma gives
\begin{align*}
 \sum_{n\ge0}\frac{q^{n^2+n}}{(q)_n}
 &=\frac1{(q^2;q)_\infty}
 \sum_{r\ge0}q^{r^2+r}\alpha_r'\\
 &=\frac1{(q;q)_\infty}
 \sum_{r\ge0}(-1)^rq^{(5r^2+3r)/2}(1-q^{2r+1})\\
 &=\frac1{(q;q)_\infty}
 \sum_{r\in\Z}(-1)^rq^{(5r^2+3r)/2}.
\end{align*}
In the last step, the subtracted term indexed by $r\ge0$ is the term with bilateral index $-r-1$.  Jacobi's triple product with base $q^5$ and $z=q^4$ gives
\[
 \sum_{r\in\Z}(-1)^rq^{(5r^2+3r)/2}
 =(q,q^4,q^5;q^5)_\infty.
\]
Divide by the fivefold dissection of $(q;q)_\infty$ to obtain \eqref{eq:RR2}.
\end{proof}

\begin{remark}[Partition interpretation]
The product side of \eqref{eq:RR1} counts partitions into parts congruent to $1$ or $4$ modulo $5$; the product side of \eqref{eq:RR2} counts parts congruent to $2$ or $3$ modulo $5$.  The sum sides count partitions with successive parts differing by at least $2$, with an additional smallest-part condition in the second identity.  Establishing the latter combinatorial interpretations requires a separate bijective or finite-polynomial argument; the proof above instead reveals the hypergeometric engine.
\end{remark}


\part{Cyclotomic arithmetic and extended indices}

\chapter{Cyclotomic factorization and congruences}\label{chap:cyclotomic}

The arithmetic of Gaussian coefficients is governed by cyclotomic polynomials.  This is the precise sense in which a primitive root of unity plays the role of a prime.  The point of this chapter is that the essential congruences do not require mysterious cancellation: they follow from product factorization, coefficient extraction, and a careful accounting of blocks.

\section{Cyclotomic valuations and carries}

Let $\Phi_d(q)$ denote the $d$th cyclotomic polynomial, normalized by
\[
 q^n-1=\prod_{d\mid n}\Phi_d(q).
\]
For a nonzero polynomial $F\in\Z[q]$, write $\nu_d(F)$ for the exponent of $\Phi_d(q)$ in its factorization.

\begin{theorem}[Cyclotomic factorization of a shifted factorial]\label{thm:cyclotomic-pochhammer}
For every $n\ge0$,
\[
 (q;q)_n=(-1)^n\prod_{d=1}^{n}\Phi_d(q)^{\floor{n/d}}.
\]
Consequently,
\[
 \qbinom nk
 =\prod_{d=1}^{n}\Phi_d(q)^{e_d(n,k)},
 \qquad
 e_d(n,k)=\floor{n/d}-\floor{k/d}-\floor{(n-k)/d}.
 \tag{16.1}\label{eq:cyclotomic-gaussian}
\]
Moreover $e_d(n,k)$ is always either $0$ or $1$.
\end{theorem}

\begin{proof}
Since $1-q^j=-(q^j-1)$ and $q^j-1=\prod_{d\mid j}\Phi_d(q)$,
\[
 (q;q)_n
 =\prod_{j=1}^{n}(1-q^j)
 =(-1)^n\prod_{j=1}^{n}\prod_{d\mid j}\Phi_d(q).
\]
For fixed $d$, the multiples of $d$ among $1,\ldots,n$ are $d,2d,\ldots,\floor{n/d}d$, so the exponent of $\Phi_d$ is $\floor{n/d}$.  Dividing the factorizations of $(q;q)_n$, $(q;q)_k$, and $(q;q)_{n-k}$ gives \eqref{eq:cyclotomic-gaussian}; the signs cancel.

Write $k=ad+r$ and $n-k=bd+s$ with $0\le r,s<d$.  Then
\[
 e_d(n,k)=\floor{(r+s)/d},
\]
which is $0$ or $1$.
\end{proof}

\begin{corollary}[Square-free cyclotomic structure]\label{cor:squarefree-cyclotomic}
Every Gaussian polynomial is a product of distinct cyclotomic polynomials.  More precisely,
\[
 \Phi_d(q)\mid\qbinom nk
 \quad\Longleftrightarrow\quad
 (k\bmod d)> (n\bmod d).
 \tag{16.2}\label{eq:carry-criterion}
\]
Equivalently, $\Phi_d$ occurs exactly when adding $k$ and $n-k$ produces a carry in the units place in base $d$.
\end{corollary}

\begin{proof}
The first assertion follows because every exponent in \eqref{eq:cyclotomic-gaussian} belongs to $\{0,1\}$.  If $n=cd+t$ and $k=ad+r$ with $0\le r,t<d$, then the residue of $n-k$ is $t-r$ when $r\le t$ and $t-r+d$ when $r>t$.  Thus the two residues add to $t$ in the first case and to $t+d$ in the second.  The latter is precisely the carry case, and by the proof of \Cref{thm:cyclotomic-pochhammer} it is equivalent to $e_d(n,k)=1$.
\end{proof}

\begin{example}
For $n=10$ and $k=4$, the values of $d$ for which $4\bmod d>10\bmod d$ are $d=5,7,8,9$.  Hence
\[
 \qbinom{10}{4}=\Phi_5(q)\Phi_7(q)\Phi_8(q)\Phi_9(q).
\]
The degrees add to $4+6+4+6=20=4(10-4)$, as required by reciprocity.
\end{example}

\begin{corollary}[MacMahon's $q$-Catalan polynomial]\label{cor:q-catalan}
\index{Catalan number!$q$-analogue}
For $n\ge0$ put
\[
 C_n(q)\defeq\frac1{\qint{n+1}}\qbinom{2n}{n}.
\]
Then $C_n(q)\in\Z[q]$; it has degree $n(n-1)$, and $C_n(1)=\frac1{n+1}\binom{2n}{n}$.
\end{corollary}

\begin{proof}
% ed.: the incoming survey proved this in one sentence -- `subtracting the
% ed.: exponent of each \Phi_d in [n+1]_q from that in the central Gaussian
% ed.: coefficient gives a nonnegative integer' -- which is the assertion, not
% ed.: an argument.  The valuation comparison is carried out below, and the
% ed.: cases d=1 and d\nmid n+1 are treated separately, as they must be.
Since $q^{N}-1=\prod_{d\mid N}\Phi_d(q)$ and $\Phi_1(q)=q-1$,
\[
 \qint{n+1}=\frac{q^{\,n+1}-1}{q-1}
 =\prod_{\substack{d\mid n+1\\ d>1}}\Phi_d(q).
\]
Hence $\nu_d(\qint{n+1})=1$ exactly for those $d>1$ dividing $n+1$, and
$\nu_d(\qint{n+1})=0$ for every other $d$, the case $d=1$ included: indeed
$\qint{n+1}$ takes the value $n+1\ne0$ at $q=1$, so $q-1$ does not divide it.

Let $d>1$ divide $n+1$.  Then $n\equiv-1\pmod d$, so $n\bmod d=d-1$, while
$2n\equiv-2\pmod d$ together with $0\le d-2<d$ gives $2n\bmod d=d-2$.
Therefore
\[
 n\bmod d=d-1>d-2=2n\bmod d,
\]
and the carry criterion \eqref{eq:carry-criterion} of
\Cref{cor:squarefree-cyclotomic}, applied with upper index $2n$ and lower
index $n$, gives $\Phi_d(q)\mid\qbinom{2n}{n}$, that is,
$\nu_d\bigl(\qbinom{2n}{n}\bigr)=1$.  Combining the two paragraphs,
\[
 \nu_d\bigl(\qbinom{2n}{n}\bigr)\ge\nu_d\bigl(\qint{n+1}\bigr)
 \qquad\text{for every }d\ge1,
\]
so $\qint{n+1}$ divides $\qbinom{2n}{n}$ in $\Q[q]$ and $C_n\in\Q[q]$.
Moreover $\qint{n+1}=1+q+\cdots+q^{n}$ is monic and
$\qbinom{2n}{n}\in\Z[q]$ by \Cref{cor:positivity}, so division with remainder
in $\Z[q]$ produces a quotient in $\Z[q]$; by uniqueness of the quotient in
$\Q[q]$ that quotient is $C_n$, whence $C_n\in\Z[q]$.

Finally, $\deg\qbinom{2n}{n}=n\cdot n=n^2$ by
\Cref{thm:qbinom-structure}(iii) and $\deg\qint{n+1}=n$, so
$\deg C_n=n^2-n$.  Evaluating the polynomial identity
$\qint{n+1}\,C_n(q)=\qbinom{2n}{n}$ at $q=1$ and using
\Cref{cor:qbinom-classical} gives $(n+1)C_n(1)=\binom{2n}{n}$.
\end{proof}

\begin{remark}\label{rem:q-catalan-positivity}
The argument above establishes integrality and says nothing about signs.  The
coefficients of $C_n(q)$ are in fact nonnegative --- $C_n$ enumerates the
Dyck paths of semilength $n$ by a suitable major-index statistic, for which
see \cite{Stanley1999} --- but that is a combinatorial theorem and not a
consequence of cyclotomic divisibility.  Compare \Cref{cor:path-area}, where
the corresponding statistic for unrestricted monotone lattice paths is the
area and the enumeration is the unconstrained $\qbinom{2n}{n}$.
\end{remark}

\section{The \texorpdfstring{$q$}{q}-Lucas congruence}

The following theorem is often called the $q$-Lucas theorem.  It is naturally proved by grouping a finite product into complete root-of-unity blocks.

\begin{lemma}[A complete root-of-unity block]\label{lem:root-block}
If $\zeta$ is a primitive $d$th root of unity, then
\[
 \prod_{j=0}^{d-1}(1+x\zeta^j)=1-(-x)^d.
 \tag{16.3}\label{eq:root-block}
\]
\end{lemma}

\begin{proof}
The roots in $x$ of the left side are $-\zeta^{-j}$, which are exactly the roots of $1-(-x)^d$.  Both polynomials have constant term $1$, so they are equal.
\end{proof}

\begin{theorem}[$q$-Lucas]\label{thm:q-lucas}
Let $a,r\ge0$, let $d\ge1$, and let $0\le b,s<d$.  Then
\[
 \Qbinom{ad+b}{rd+s}{q}
 \equiv
 \binom ar\Qbinom bsq
 \pmod{\Phi_d(q)}.
 \tag{16.4}\label{eq:q-lucas}
\]
Here the right side is understood as $0$ when $r>a$ or $s>b$.
\end{theorem}

\begin{proof}
It suffices to prove equality after substituting an arbitrary primitive $d$th root $\zeta$ for $q$, because $\Phi_d$ is the minimal polynomial of $\zeta$ over $\Q$.

The finite $q$-binomial theorem gives
\[
 \prod_{j=0}^{ad+b-1}(1+x\zeta^j)
 =\sum_{h=0}^{ad+b}\zeta^{\binom h2}
   \Qbinom{ad+b}{h}{\zeta}x^h.
 \tag{16.5}\label{eq:lucas-coeff-left}
\]
Group the product into $a$ complete blocks and one incomplete block.  By \Cref{lem:root-block},
\[
 \prod_{j=0}^{ad+b-1}(1+x\zeta^j)
 =\bigl(1-(-x)^d\bigr)^a
   \prod_{j=0}^{b-1}(1+x\zeta^j).
\]
The coefficient of $x^{rd+s}$ on the right is
\[
 \binom ar(-1)^{r(d+1)}
 \zeta^{\binom s2}\Qbinom bs{\zeta}.
 \tag{16.6}\label{eq:lucas-coeff-right}
\]
Comparing \eqref{eq:lucas-coeff-left} and \eqref{eq:lucas-coeff-right} leaves only the phase.  We claim that
\[
 (-1)^{r(d+1)}
 \zeta^{\binom s2-\binom{rd+s}{2}}=1.
 \tag{16.7}\label{eq:lucas-phase}
\]
Indeed,
\[
 2\left(\binom{rd+s}{2}-\binom s2\right)
 =rd(rd+2s-1).
\]
If $d$ is odd, $r(rd+2s-1)$ is even, so the exponent difference is a multiple of $d$, while $r(d+1)$ is even.  If $d$ is even, then $\zeta^{d/2}=-1$ and $rd+2s-1$ is odd; hence the second factor in \eqref{eq:lucas-phase} equals $(-1)^r$, as does the first.  Their product is $1$.  The coefficient comparison therefore gives \eqref{eq:q-lucas} at $q=\zeta$, completing the proof.
\end{proof}

\begin{corollary}[Evaluation at $q=-1$]\label{cor:qminusone}
For $0\le k\le n$,
\[
 \qbinom nk\bigg|_{q=-1}
 =\begin{cases}
 \displaystyle\binom{n/2}{k/2},&n\text{ even and }k\text{ even},\\[7pt]
 0,&n\text{ even and }k\text{ odd},\\[7pt]
 \displaystyle\binom{(n-1)/2}{\floor{k/2}},&n\text{ odd}.
 \end{cases}
\]
\end{corollary}

\begin{proof}
Apply \Cref{thm:q-lucas} with $d=2$.  Write $n=2a+b$ and $k=2r+s$, where $b,s\in\{0,1\}$.  Since $\Qbinom bs{-1}$ is $1$ when $s\le b$ and $0$ otherwise, the three cases follow.
\end{proof}

\section{Cyclic sieving for subsets}

A triple $(X,C,F(q))$, where a finite cyclic group $C=\angles c$ acts on a finite set $X$ and $F(q)\in\Z[q]$, exhibits the \emph{cyclic sieving phenomenon} when
\[
 \#\operatorname{Fix}(c^j)=F(\omega^j)
\]
for every $j$, where $\omega$ is a primitive $\abs C$th root of unity.

\begin{theorem}[Cyclic sieving for $k$-subsets]\label{thm:csp-subsets}
Let $X$ be the set of $k$-element subsets of $\Z/n\Z$, and let $C_n$ act by translation.  Then
\[
 \left(X,C_n,\qbinom nk\right)
\]
exhibits the cyclic sieving phenomenon.
\end{theorem}

\begin{proof}
Fix $j$, and let $d$ be the order of translation by $j$.  Its orbits on $\Z/n\Z$ all have size $d$, and there are $n/d$ such orbits.  A subset is fixed precisely when it is a union of these orbits.  Thus
\[
 \#\operatorname{Fix}(c^j)
 =\begin{cases}
 \displaystyle\binom{n/d}{k/d},&d\mid k,\\
 0,&d\nmid k.
 \end{cases}
 \tag{16.8}\label{eq:csp-fixed}
\]
The number $\omega^j$ is a primitive $d$th root.  Write $n=(n/d)d$ and $k=rd+s$ with $0\le s<d$.  By $q$-Lucas,
\[
 \qbinom nk\bigg|_{q=\omega^j}
 =\binom{n/d}{r}\Qbinom 0s{\omega^j}.
\]
This is $0$ if $s>0$, and it is $\binom{n/d}{k/d}$ if $s=0$, exactly as in \eqref{eq:csp-fixed}.
\end{proof}

\section{A \texorpdfstring{$q$}{q}-Babbage congruence modulo a cyclotomic square}

The modulus $\Phi_n(q)^2$ records not only the value at a primitive $n$th root, but also first-order contact there.  The next congruence is a particularly transparent example.  Its block proof is also a useful template for stronger $q$-congruences; see Andrews~\cite{Andrews1999Congruences} and Straub~\cite{Straub2011Ljunggren}.

\begin{theorem}[$q$-Babbage congruence]\label{thm:q-babbage}
For integers $a\ge b\ge0$ and $n\ge1$,
\[
 \Qbinom{an}{bn}{q}
 \equiv
 \Qbinom ab{q^{n^2}}
 \pmod{\Phi_n(q)^2}.
 \tag{16.9}\label{eq:q-babbage}
\]
\end{theorem}

\begin{proof}
Split the finite product into $a$ consecutive blocks of length $n$:
\[
 \prod_{j=0}^{an-1}(1+xq^j)
 =\prod_{i=0}^{a-1}\prod_{j=0}^{n-1}(1+xq^{in+j}).
\]
Apply the finite $q$-binomial theorem to each block.  Comparing the coefficient of $x^{bn}$ gives
\begin{align}
 q^{\binom{bn}{2}}\Qbinom{an}{bn}{q}
 &=\sum_{\substack{0\le k_0,\ldots,k_{a-1}\le n\\
                    k_0+\cdots+k_{a-1}=bn}}
 q^{\sum_i\binom{k_i}{2}+n\sum_i i k_i}
 \prod_{i=0}^{a-1}\Qbinom n{k_i}q.
 \tag{16.10}\label{eq:block-expansion}
\end{align}
If $0<k_i<n$, then $\Phi_n(q)$ divides $\Qbinom n{k_i}q$ by \Cref{cor:squarefree-cyclotomic}.  A tuple contributing to \eqref{eq:block-expansion} cannot have exactly one such intermediate entry: all other entries would be $0$ or $n$, making the sum of the $k_i$ nonzero modulo $n$.  Therefore every tuple containing an intermediate entry contributes a multiple of $\Phi_n(q)^2$.

Modulo $\Phi_n(q)^2$, only tuples with each $k_i\in\{0,n\}$ survive.  Such a tuple is specified by a $b$-element subset $S\subseteq\{0,1,\ldots,a-1\}$.  After division by $q^{\binom{bn}{2}}$, its exponent is
\begin{align*}
 n^2\sum_{i\in S}i+b\binom n2-\binom{bn}{2}
 &=n^2\left(\sum_{i\in S}i-\binom b2\right).
\end{align*}
Hence
\[
 \Qbinom{an}{bn}{q}
 \equiv
 \sum_{\substack{S\subseteq\{0,\ldots,a-1\}\\\abs S=b}}
 (q^{n^2})^{\sum_{i\in S}i-\binom b2}
 \pmod{\Phi_n(q)^2}.
\]
The subset form of the finite $q$-binomial theorem identifies the sum as $\Qbinom ab{q^{n^2}}$.
\end{proof}

\begin{remark}[Why the square appears]
The decisive observation in the proof is combinatorial: because the target degree is a multiple of the block length, nontrivial block choices occur in pairs.  Each nontrivial block supplies one factor $\Phi_n(q)$, so two of them supply the square.  Congruences modulo $\Phi_n^3$ require genuinely finer first- and second-order information; the $q$-Ljunggren congruences of Straub and subsequent refinements belong to that deeper layer.
\end{remark}

\begin{corollary}[Primitive-root value and derivative]\label{cor:babbage-derivative}
Let $\zeta$ be a primitive $n$th root of unity.  Then
\[
 \Qbinom{an}{bn}{\zeta}=\binom ab,
\]
and the derivatives at $q=\zeta$ of the two sides of \eqref{eq:q-babbage} are equal.
\end{corollary}

\begin{proof}
The first assertion follows either from $q$-Lucas or from \eqref{eq:q-babbage}, because $\zeta^{n^2}=1$.  If a polynomial is divisible by $\Phi_n^2$, it has a zero of order at least two at every primitive $n$th root, so both its value and first derivative vanish there.
\end{proof}

\chapter{Negative upper indices and geometric Newton interpolation}\label{chap:qnewton}

Ordinary binomial coefficients acquire a useful extension to negative upper arguments.  Gaussian coefficients do as well, but the powers of $q$ are essential: omitting them destroys the Pascal recurrences.  The same algebra leads naturally to interpolation on a geometric grid.

\section{Generalized Gaussian coefficients}

\begin{definition}
For an integer $N$ and $k\ge0$, define
\[
 \Qbinom Nkq
 :=\frac{(q^{N-k+1};q)_k}{(q;q)_k}.
 \tag{17.1}\label{eq:generalized-gaussian}
\]
For $N\ge0$ this agrees with the usual Gaussian coefficient, including the value $0$ when $k>N$.
\end{definition}

\begin{proposition}[Negative-index reflection]\label{prop:negative-gaussian}
For $m\ge1$ and $k\ge0$,
\[
 \Qbinom{-m}{k}q
 =(-1)^kq^{-mk-\binom k2}
   \Qbinom{m+k-1}{k}q.
 \tag{17.2}\label{eq:negative-reflection}
\]
\end{proposition}

\begin{proof}
From the definition,
\begin{align*}
 (q^{-m-k+1};q)_k
 &=\prod_{j=0}^{k-1}\left(1-q^{-(m+k-1-j)}\right)\\
 &=(-1)^kq^{-\sum_{j=0}^{k-1}(m+k-1-j)}
   \prod_{j=0}^{k-1}(1-q^{m+k-1-j})\\
 &=(-1)^kq^{-mk-\binom k2}(q^m;q)_k.
\end{align*}
After division by $(q;q)_k$, the final ratio is $\Qbinom{m+k-1}{k}q$.
\end{proof}

\begin{theorem}[Pascal recurrences for all integral upper indices]\label{thm:generalized-pascal}
For every integer $N$ and every $k\ge1$,
\begin{align}
 \Qbinom Nkq
 &=\Qbinom{N-1}{k-1}q+q^k\Qbinom{N-1}{k}q,
 \tag{17.3}\label{eq:general-pascal-a}\\
 &=q^{N-k}\Qbinom{N-1}{k-1}q+\Qbinom{N-1}{k}q.
 \tag{17.4}\label{eq:general-pascal-b}
\end{align}
Both formulas are identities in $\Q(q)$.
\end{theorem}

\begin{proof}
Put the two terms on the right of \eqref{eq:general-pascal-a} over the common denominator $(q;q)_k$.  Since
\[
 \Qbinom{N-1}{k-1}q
 =\frac{(q^{N-k+1};q)_{k-1}(1-q^k)}{(q;q)_k}
\]
the numerator of the right side is
\begin{align*}
 &(q^{N-k+1};q)_{k-1}
 \left[(1-q^k)+q^k(1-q^{N-k})\right]\\
 &\qquad=(q^{N-k+1};q)_{k-1}(1-q^N)
 =(q^{N-k+1};q)_k,
\end{align*}
which proves \eqref{eq:general-pascal-a}.  For \eqref{eq:general-pascal-b}, the corresponding bracket is
\[
 q^{N-k}(1-q^k)+(1-q^{N-k})=1-q^N.
\]
The same common-denominator calculation applies.
\end{proof}

\begin{corollary}[Reciprocal product series]\label{cor:reciprocal-product-negative}
For $m\ge1$ and $\abs z<1$,
\[
 \frac1{(z;q)_m}
 =\sum_{k=0}^{\infty}\Qbinom{m+k-1}{k}q z^k
 =\sum_{k=0}^{\infty}(-1)^kq^{mk+\binom k2}
   \Qbinom{-m}{k}qz^k.
 \tag{17.5}\label{eq:negative-reciprocal-series}
\]
The first equality also holds in the formal power-series ring $\Z[q][[z]]$.
\end{corollary}

\begin{proof}
The reciprocal finite $q$-binomial theorem, already proved in \Cref{thm:reciprocal-finite}, gives the first expansion.  Substitute \eqref{eq:negative-reflection} to obtain the second.  Formal validity follows because each coefficient of $z^k$ is a polynomial and multiplication by the finite product $(z;q)_m$ can be checked coefficientwise.
\end{proof}

\section{An arbitrary upper parameter}

When a branch of $q^\alpha$ has been fixed, it is often convenient to define
\[
 \Qbinom{\alpha}{k}q
 =\frac{(q^{\alpha-k+1};q)_k}{(q;q)_k}.
 \tag{17.6}\label{eq:complex-upper-gaussian}
\]
The proofs of the Pascal recurrences are purely algebraic in $q^\alpha$ and therefore remain valid.  A related, and usually more analytic, expansion is the following.

\begin{theorem}[Generalized upper-parameter series]\label{thm:upper-parameter-series}
For $\abs q<1$, $\abs z<1$, and any parameter $\alpha$ for which the factors are defined,
\[
 \frac{(q^{\alpha+1}z;q)_\infty}{(z;q)_\infty}
 =\sum_{k=0}^{\infty}
   \Qbinom{\alpha+k}{k}q z^k.
 \tag{17.7}\label{eq:upper-parameter-series}
\]
\end{theorem}

\begin{proof}
Apply the infinite $q$-binomial theorem with $a=q^{\alpha+1}$:
\[
 \frac{(q^{\alpha+1}z;q)_\infty}{(z;q)_\infty}
 =\sum_{k\ge0}\frac{(q^{\alpha+1};q)_k}{(q;q)_k}z^k.
\]
The coefficient is exactly $\Qbinom{\alpha+k}{k}q$ by \eqref{eq:complex-upper-gaussian}.
\end{proof}

\section{Two complex arguments and generalized binomial series}

The extension in \Cref{def:complex-qshift} also permits both arguments of a Gaussian coefficient to be complex.  Throughout this section take $0<q<1$ and set $q^z=\exp(z\log q)$, using the real logarithm of $q$.

\begin{definition}\label{def:complex-gaussian}
For $\alpha,\beta\in\C$, define
\[
 \genfrac{[}{]}{0pt}{}{\alpha}{\beta}_q
 :=\frac{(q^{\beta+1},q^{\alpha-\beta+1};q)_\infty}
 {(q,q^{\alpha+1};q)_\infty},
 \tag{17.8}\label{eq:complex-gaussian}
\]
whenever the denominator is nonzero.
\end{definition}

\begin{theorem}[Gamma representation and Pascal laws]\label{thm:complex-gaussian-properties}
Wherever the expressions are defined,
\begin{align}
 \genfrac{[}{]}{0pt}{}{\alpha}{\beta}_q
 &=\frac{\Gamma_q(\alpha+1)}
 {\Gamma_q(\beta+1)\Gamma_q(\alpha-\beta+1)},
 \tag{17.9}\label{eq:complex-gaussian-gamma}\\
 \genfrac{[}{]}{0pt}{}{\alpha}{\beta}_q
 &=\genfrac{[}{]}{0pt}{}{\alpha}{\alpha-\beta}_q,
 \tag{17.10}\label{eq:complex-gaussian-symmetry}\\
 \genfrac{[}{]}{0pt}{}{\alpha+1}{\beta}_q
 &=q^\beta\genfrac{[}{]}{0pt}{}{\alpha}{\beta}_q
   +\genfrac{[}{]}{0pt}{}{\alpha}{\beta-1}_q,
 \tag{17.11}\label{eq:complex-pascal-a}\\
 &=\genfrac{[}{]}{0pt}{}{\alpha}{\beta}_q
   +q^{\alpha+1-\beta}
    \genfrac{[}{]}{0pt}{}{\alpha}{\beta-1}_q.
 \tag{17.12}\label{eq:complex-pascal-b}
\end{align}
If $\beta=k\in\N_0$, this definition agrees with \eqref{eq:complex-upper-gaussian}.
\end{theorem}

\begin{proof}
Substitution of the product definition of $\Gamma_q$ proves \eqref{eq:complex-gaussian-gamma}; all powers of $1-q$ cancel.  Symmetry is immediate either from that quotient or directly from \eqref{eq:complex-gaussian}.  For $\beta=k\in\N_0$, tail cancellation gives
\[
 \frac{(q^{k+1};q)_\infty}{(q;q)_\infty}
 =\frac1{(q;q)_k},
 \qquad
 \frac{(q^{\alpha-k+1};q)_\infty}
 {(q^{\alpha+1};q)_\infty}
 =(q^{\alpha-k+1};q)_k,
\]
which is exactly \eqref{eq:complex-upper-gaussian}.

For the first Pascal law, divide every term by
$\genfrac{[}{]}{0pt}{}{\alpha}{\beta}_q$.  The gamma functional equation gives
\[
 \frac{\genfrac{[}{]}{0pt}{}{\alpha+1}{\beta}_q}
      {\genfrac{[}{]}{0pt}{}{\alpha}{\beta}_q}
 =\frac{[\alpha+1]_q}{[\alpha-\beta+1]_q},
 \qquad
 \frac{\genfrac{[}{]}{0pt}{}{\alpha}{\beta-1}_q}
      {\genfrac{[}{]}{0pt}{}{\alpha}{\beta}_q}
 =\frac{[\beta]_q}{[\alpha-\beta+1]_q}.
\]
The identity
\(
 [\alpha+1]_q=q^\beta[\alpha-\beta+1]_q+[\beta]_q
\)
proves \eqref{eq:complex-pascal-a}.  The alternative identity
\(
 [\alpha+1]_q=[\alpha-\beta+1]_q
 +q^{\alpha+1-\beta}[\beta]_q
\)
proves \eqref{eq:complex-pascal-b}.
\end{proof}

\begin{corollary}[Classical generalized binomial]\label{cor:complex-gaussian-classical}
If $\alpha>\beta>-1$ are real, then
\[
 \lim_{q\to1^-}\genfrac{[}{]}{0pt}{}{\alpha}{\beta}_q
 =\frac{\Gamma(\alpha+1)}
 {\Gamma(\beta+1)\Gamma(\alpha-\beta+1)}.
\]
\end{corollary}

\begin{proof}
Apply the locally uniform $q$-gamma limit in \Cref{thm:qgamma-classical-limit} to \eqref{eq:complex-gaussian-gamma}.
\end{proof}

\begin{theorem}[Generalized finite and reciprocal $q$-binomial theorems]\label{thm:generalized-qbinomial}
Let $0<q<1$.  If $\abs{zq^\alpha}<1$, then
\[
 (z;q)_\alpha
 =\sum_{k=0}^{\infty}(-1)^kq^{\binom k2}
   \genfrac{[}{]}{0pt}{}{\alpha}{k}_q z^k.
 \tag{17.13}\label{eq:generalized-finite-qbinomial}
\]
If $\abs z<1$, then
\[
 \frac1{(z;q)_\alpha}
 =\sum_{k=0}^{\infty}
   \genfrac{[}{]}{0pt}{}{\alpha+k-1}{k}_q z^k.
 \tag{17.14}\label{eq:generalized-reciprocal-qbinomial}
\]
For a nonnegative integral $\alpha$, the first series terminates and becomes Euler's finite theorem; in that case the restriction $\abs{zq^\alpha}<1$ can be removed because both sides are polynomials.  For a positive integral $\alpha$, the second becomes the reciprocal finite theorem.
\end{theorem}

\begin{proof}
Apply the infinite $q$-binomial theorem with
\(
 a=q^{-\alpha}
\)
and argument $zq^\alpha$:
\[
 \frac{(z;q)_\infty}{(zq^\alpha;q)_\infty}
 =\sum_{k\ge0}\frac{(q^{-\alpha};q)_k}{(q;q)_k}
 (zq^\alpha)^k.
\]
The left side is $(z;q)_\alpha$.  Reversing the finite numerator product gives
\[
 \frac{(q^{-\alpha};q)_k}{(q;q)_k}q^{\alpha k}
 =(-1)^kq^{\binom k2}
   \genfrac{[}{]}{0pt}{}{\alpha}{k}_q,
\]
which proves \eqref{eq:generalized-finite-qbinomial}.

For the reciprocal formula, apply the infinite theorem with $a=q^\alpha$ and argument $z$:
\[
 \frac{(q^\alpha z;q)_\infty}{(z;q)_\infty}
 =\sum_{k\ge0}\frac{(q^\alpha;q)_k}{(q;q)_k}z^k.
\]
The left side is $(z;q)_\alpha^{-1}$ and the coefficient is
$\genfrac{[}{]}{0pt}{}{\alpha+k-1}{k}_q$.  The integral specializations follow from the consistency results already proved.
\end{proof}

\section{Newton interpolation on the grid \texorpdfstring{$1,q,q^2,\ldots$}{1, q, q-squared, ...}}

The usual Newton basis is adapted to an arbitrary sequence of distinct nodes.  At geometric nodes its basis polynomials are $q$-Pochhammer symbols.

\begin{theorem}[Geometric Newton interpolation]\label{thm:geometric-newton}
Assume $q^i\ne q^j$ whenever $0\le i<j\le n$.  For arbitrary values $y_0,\ldots,y_n$ there is a unique polynomial $P$ of degree at most $n$ satisfying
\[
 P(q^j)=y_j\qquad(0\le j\le n).
\]
It has the Newton expansion
\[
 P(x)=\sum_{k=0}^{n}c_k\prod_{r=0}^{k-1}(x-q^r)
 =\sum_{k=0}^{n}c_k(-1)^kq^{\binom k2}(x;q^{-1})_k,
 \tag{17.15}\label{eq:geometric-newton}
\]
where
\begin{align}
 c_k
 &=\sum_{j=0}^{k}
 \frac{y_j}{\prod_{\substack{0\le r\le k\\r\ne j}}(q^j-q^r)}
 \tag{17.16}\label{eq:geometric-divided-difference}\\
 &=\sum_{j=0}^{k}
 \frac{(-1)^j q^{-jk+\binom{j+1}{2}}}{(q;q)_j(q;q)_{k-j}}\,y_j.
 \tag{17.17}\label{eq:geometric-divided-difference-explicit}
\end{align}
\end{theorem}

\begin{proof}
For distinct nodes $x_0,\ldots,x_n$, Lagrange interpolation gives existence and uniqueness.  The coefficient of the Newton basis polynomial
\(
 \prod_{r=0}^{k-1}(x-x_r)
\)
is the divided difference
\[
 y[x_0,\ldots,x_k]
 =\sum_{j=0}^{k}\frac{y_j}{\prod_{r\ne j}(x_j-x_r)};
\]
this follows by comparing the coefficient of $x^k$ in the Lagrange interpolant through the first $k+1$ nodes.  Taking $x_j=q^j$ proves \eqref{eq:geometric-divided-difference}.

It remains to simplify the denominator.  Splitting the factors below and above $j$ gives
\begin{align*}
 \prod_{\substack{0\le r\le k\\r\ne j}}(q^j-q^r)
 &=\left[\prod_{r=0}^{j-1}-q^r(1-q^{j-r})\right]
   \left[\prod_{r=j+1}^{k}q^j(1-q^{r-j})\right]\\
 &=(-1)^jq^{\binom j2+j(k-j)}(q;q)_j(q;q)_{k-j}\\
 &=(-1)^jq^{jk-\binom{j+1}{2}}(q;q)_j(q;q)_{k-j}.
\end{align*}
Taking reciprocals yields \eqref{eq:geometric-divided-difference-explicit}.  Finally,
\[
 \prod_{r=0}^{k-1}(x-q^r)
 =(-1)^kq^{\binom k2}\prod_{r=0}^{k-1}(1-xq^{-r})
 =(-1)^kq^{\binom k2}(x;q^{-1})_k,
\]
which proves the second form of \eqref{eq:geometric-newton}.
\end{proof}

\begin{corollary}[Triangular reconstruction]\label{cor:geometric-triangular}
The coefficients in \eqref{eq:geometric-newton} can be recovered successively by
\[
 c_0=y_0,
 \qquad
 c_k=\frac{y_k-\sum_{r=0}^{k-1}c_r\prod_{j=0}^{r-1}(q^k-q^j)}
 {\prod_{j=0}^{k-1}(q^k-q^j)}.
 \tag{17.18}\label{eq:triangular-newton}
\]
\end{corollary}

\begin{proof}
Evaluate \eqref{eq:geometric-newton} at $x=q^k$.  Every term with index greater than $k$ vanishes, while the coefficient multiplying $c_k$ is nonzero by the distinct-node hypothesis.  Solving for $c_k$ gives the formula.
\end{proof}

\begin{remark}
Kupershmidt's $q$-Newton viewpoint emphasizes that the classical binomial theorem, Newton interpolation, and $q$-difference calculus are not separate constructions: each is controlled by a basis whose successive factors form a geometric progression.  The finite formula above is algebraic and needs no convergence hypothesis; infinite Newton series require additional growth conditions on the data.
\end{remark}

\part{Central sums, limits, and effective methods}

\chapter{Central \texorpdfstring{$q$}{q}-binomial sums and Lambert series}\label{chap:central}

Central Gaussian coefficients frequently become tractable only after a seemingly extraneous factor is inserted.  The elementary identity in the next lemma shows why the factor $(-q;q)_{2k}$ is natural: it converts a central coefficient at base $q^2$ into a single quotient of shifted factorials.  This chapter develops both a regular $q$-beta evaluation and a deleted-singularity identity proved recently by Chern, Dilcher, and Jiu~\cite{ChernDilcherJiu2026}.

\section{The central-factor reduction}

\begin{lemma}\label{lem:central-reduction}
For every $k\ge0$,
\[
 \Qbinom{2k}{k}{q^2}\frac1{(-q;q)_{2k}}
 =\frac{(q;q^2)_k}{(q^2;q^2)_k}.
 \tag{18.1}\label{eq:central-reduction}
\]
\end{lemma}

\begin{proof}
The elementary product identity
\[
 (q;q)_{2k}(-q;q)_{2k}=(q^2;q^2)_{2k}
\]
gives
\begin{align*}
 \Qbinom{2k}{k}{q^2}\frac1{(-q;q)_{2k}}
 &=\frac{(q^2;q^2)_{2k}}{(q^2;q^2)_k^2}
   \frac{(q;q)_{2k}}{(q^2;q^2)_{2k}}\\
 &=\frac{(q;q^2)_k(q^2;q^2)_k}{(q^2;q^2)_k^2},
\end{align*}
which is the claimed quotient.
\end{proof}

For a complex parameter $z$ and $0<q<1$, we use the analytic extension
\[
 [z]_q=\frac{1-q^z}{1-q},
\]
with the real logarithm in $q^z=\exp(z\log q)$.

\section{A \texorpdfstring{$q$}{q}-beta evaluation}

\begin{theorem}[A regular central $q$-binomial sum]\label{thm:regular-central-sum}
Let $0<q<1$, and let $\alpha\in\C$ be such that none of the denominators below vanishes.  Then
\begin{align}
 &\sum_{k=0}^{\infty}
 \Qbinom{2k}{k}{q^2}
 \frac{q^k}{(-q;q)_{2k}[2k+1+\alpha]_q}
 \notag\\
 &\qquad=
 \frac{\Gamma_{q^2}(3/2)\,
       \Gamma_{q^2}((\alpha+1)/2)}
      {\Gamma_{q^2}((\alpha+2)/2)}.
 \tag{18.2}\label{eq:regular-central-sum}
\end{align}
In particular, for real $\alpha$ it is enough to exclude
$\alpha=-1,-3,-5,\ldots$.
\end{theorem}

\begin{proof}
By \Cref{lem:central-reduction} and the telescoping quotient
\[
 \frac1{[2k+1+\alpha]_q}
 =\frac1{[\alpha+1]_q}
   \frac{(q^{\alpha+1};q^2)_k}
        {(q^{\alpha+3};q^2)_k},
\]
the left side is
\[
 \frac1{[\alpha+1]_q}
 {}_2\phi_1\!\left[
 \begin{matrix}q,q^{\alpha+1}\\q^{\alpha+3}\end{matrix};
 q^2,q\right].
\]
The $q$-Gauss sum from \Cref{thm:q-gauss}, with base $q^2$, gives
\[
 \frac1{[\alpha+1]_q}
 \frac{(q^{\alpha+2},q^2;q^2)_\infty}
      {(q^{\alpha+3},q;q^2)_\infty}.
 \tag{18.3}\label{eq:regular-central-product}
\]
Now
\[
 (q^{\alpha+1};q^2)_\infty
 =(1-q^{\alpha+1})(q^{\alpha+3};q^2)_\infty,
 \qquad
 (q;q^2)_\infty=(1-q)(q^3;q^2)_\infty.
\]
Using these two relations to simplify \eqref{eq:regular-central-product}, we obtain
\[
 \frac{(q^2;q^2)_\infty(q^{\alpha+2};q^2)_\infty}
 {(q^3;q^2)_\infty(q^{\alpha+1};q^2)_\infty}.
 \tag{18.4}\label{eq:regular-central-product2}
\]
Finally insert
\[
 \Gamma_{q^2}(z)
 =(1-q^2)^{1-z}
 \frac{(q^2;q^2)_\infty}{(q^{2z};q^2)_\infty}.
\]
In the quotient on the right of \eqref{eq:regular-central-sum}, all powers of $1-q^2$ cancel, and the remaining products are exactly \eqref{eq:regular-central-product2}.
\end{proof}

\begin{corollary}[Classical limit of the regular sum]\label{cor:regular-central-classical}
For $\alpha\notin\{-1,-3,-5,\ldots\}$,
\[
 \sum_{k=0}^{\infty}\binom{2k}{k}
 \frac1{4^k(2k+1+\alpha)}
 =\frac{\Gamma(3/2)\Gamma((\alpha+1)/2)}
        {\Gamma((\alpha+2)/2)}.
 \tag{18.5}\label{eq:regular-central-classical}
\]
\end{corollary}

\begin{proof}
The series is
\[
 \frac1{\alpha+1}
 {}_2F_1\!\left(
 \begin{matrix}1/2,(\alpha+1)/2\\(\alpha+3)/2\end{matrix};1\right).
\]
Gauss's classical summation gives
\[
 \frac1{\alpha+1}
 \frac{\Gamma((\alpha+3)/2)\Gamma(1/2)}
      {\Gamma((\alpha+2)/2)}.
\]
Since $\Gamma((\alpha+3)/2)=((\alpha+1)/2)\Gamma((\alpha+1)/2)$ and $\Gamma(1/2)/2=\Gamma(3/2)$, this is \eqref{eq:regular-central-classical}.  The same identity is also the $q\to1^-$ limit of \Cref{thm:regular-central-sum} by the $q$-gamma limit in \Cref{thm:qgamma-classical-limit}.
\end{proof}

\section{Deleting a singular term}

The regular formula becomes singular when $\alpha=-2m-1$.  Removing the offending $k=m$ term does not merely produce a finite part of the same gamma quotient; it yields an alternating Lambert-series tail.

\begin{theorem}[Deleted-singularity central identity]\label{thm:singular-central}
Let $m\ge0$ and $0<q<1$.  Then
\begin{align}
 &\sum_{\substack{k\ge0\\k\ne m}}
 \Qbinom{2k}{k}{q^2}
 \frac{q^k}{(-q;q)_{2k}[k-m]_{q^2}}
 \notag\\
 &\quad=
 \Qbinom{2m}{m}{q^2}
 \frac{(1+q)q^m}{(-q;q)_{2m}}
 \sum_{r=2m+1}^{\infty}
 \frac{(-1)^{r-1}q^r}{[r]_q}.
 \tag{18.6}\label{eq:singular-central}
\end{align}
\end{theorem}

\begin{proof}
By \Cref{lem:central-reduction}, the left side is
\[
 (1-q^2)
 \sum_{\substack{k\ge0\\k\ne m}}
 \frac{(q;q^2)_kq^k}
 {(q^2;q^2)_k(1-q^{2k-2m})}.
 \tag{18.7}\label{eq:singular-central-reduced}
\]
Introduce an auxiliary parameter $a$ and define
\begin{equation}
\begin{aligned}
 F(a,q)
 :=\frac1{1-aq^{-2m}}\Bigg(&
 {}_2\phi_1\!\left[
 \begin{matrix}q,aq^{-2m}\\aq^{-2m+2}\end{matrix};q^2,q\right]\\
 &-\frac{(q,aq^{-2m};q^2)_m q^m}
 {(q^2,aq^{-2m+2};q^2)_m}
 \Bigg).
\end{aligned}
 \tag{18.8}\label{eq:F-auxiliary}
\end{equation}
The elementary cancellation
\[
 \frac{(aq^{-2m};q^2)_k}
      {(aq^{-2m+2};q^2)_k}
 =\frac{1-aq^{-2m}}{1-aq^{2k-2m}}
\]
shows that
\[
 F(a,q)=
 \sum_{\substack{k\ge0\\k\ne m}}
 \frac{(q;q^2)_kq^k}
 {(q^2;q^2)_k(1-aq^{2k-2m})}.
 \tag{18.9}\label{eq:F-series}
\]
Thus \eqref{eq:singular-central-reduced} equals
$(1-q^2)\lim_{a\to1}F(a,q)$.

Apply $q$-Gauss with base $q^2$ to the ${}_2\phi_1$ in \eqref{eq:F-auxiliary}.  We obtain
\begin{align*}
 F(a,q)
 &=\frac1{1-aq^{-2m}}
 \left(
 \frac{(q^2,aq^{-2m+1};q^2)_\infty}
 {(q,aq^{-2m+2};q^2)_\infty}
 -\frac{(q,aq^{-2m};q^2)_m q^m}
 {(q^2,aq^{-2m+2};q^2)_m}
 \right)\\
 &=\frac{H(a)-C_m}{1-a},
 \tag{18.10}\label{eq:F-H}
\end{align*}
where
\begin{align*}
 H(a)&=
 \frac{(q^2,aq;q^2)_\infty(aq^{-2m+1};q^2)_m}
 {(q,aq^2;q^2)_\infty(aq^{-2m};q^2)_m},\\
 C_m&=\frac{(q;q^2)_m q^m}{(q^2;q^2)_m}.
\end{align*}
The reversal identity for finite shifted factorials gives
\[
 \frac{(q^{-2m+1};q^2)_m}{(q^{-2m};q^2)_m}
 =q^m\frac{(q;q^2)_m}{(q^2;q^2)_m},
\]
so $H(1)=C_m$.  Hence L'Hospital's rule yields
\[
 \lim_{a\to1}F(a,q)=-H'(1)
 =-H(1)\left.\frac{\partial}{\partial a}\log H(a)\right|_{a=1}.
 \tag{18.11}\label{eq:F-logderivative}
\]
Differentiating the logarithms of the product factors gives
\begin{align*}
 \lim_{a\to1}F(a,q)
 =C_m\biggl(&
 \sum_{r=1}^{\infty}
 \frac{(-1)^{r-1}q^r}{1-q^r}
 +\sum_{r=1}^{2m}
 \frac{(-1)^{r-1}q^{-r}}{1-q^{-r}}
 \biggr).
 \tag{18.12}\label{eq:F-two-sums}
\end{align*}
For $r\ge1$,
\[
 \frac{q^{-r}}{1-q^{-r}}=-\frac1{1-q^r}.
\]
Consequently, the two summands in \eqref{eq:F-two-sums} combine, for $1\le r\le2m$, to $(-1)^r$.  Their sum is $0$ because $2m$ is even.  Therefore
\[
 \lim_{a\to1}F(a,q)
 =C_m\sum_{r=2m+1}^{\infty}
 \frac{(-1)^{r-1}q^r}{1-q^r}.
 \tag{18.13}\label{eq:F-tail}
\]
Multiply by $1-q^2=(1+q)(1-q)$, use
$(1-q)/(1-q^r)=1/[r]_q$, and finally use \eqref{eq:central-reduction} at $k=m$ to rewrite
\[
 C_m=q^m\frac{(q;q^2)_m}{(q^2;q^2)_m}
 =q^m\Qbinom{2m}{m}{q^2}\frac1{(-q;q)_{2m}}.
\]
The result is exactly \eqref{eq:singular-central}.
\end{proof}

\begin{remark}[The mechanism]
The deleted term is encoded by the auxiliary parameter $a$.  The $q$-Gauss product has a removable zero-over-zero singularity at $a=1$, and its logarithmic derivative turns products into Lambert series.  This parameter-differentiation mechanism recurs throughout the theory of harmonic-number and Lambert-series refinements of hypergeometric summations.
\end{remark}

\section{Separate finite and infinite components}

The deleted sum in \Cref{thm:singular-central} naturally splits at the missing
index.  Each component has a simpler one-sided Lambert-series evaluation.

\begin{theorem}[Infinite and finite components]\label{thm:central-components}
Let $m\ge0$ and $0<q<1$, and put
\[
 P_m(q)=\Qbinom{2m}{m}{q^2}
 \frac{(1+q)q^m}{(-q;q)_{2m}}.
\]
Then
\begin{align}
 &\sum_{k=m+1}^{\infty}
 \Qbinom{2k}{k}{q^2}
 \frac{q^k}{(-q;q)_{2k}[k-m]_{q^2}}
 \notag\\
 &\qquad=P_m(q)\left(
 \sum_{r=1}^{\infty}\frac{q^{2r-1}}{[2r-1]_q}
 -\sum_{r=1}^{\infty}\frac{q^{2r+2m}}{[2r+2m]_q}
 \right),
 \tag{18.14}\label{eq:central-infinite-component}
\end{align}
and
\[
 \sum_{k=0}^{m-1}
 \Qbinom{2k}{k}{q^2}
 \frac{q^k}{(-q;q)_{2k}[k-m]_{q^2}}
 =-P_m(q)\sum_{r=1}^{m}\frac{q^{2r-1}}{[2r-1]_q}.
 \tag{18.15}\label{eq:central-finite-component}
\]
For $m=0$, the finite sum and the sum on its right are both empty.
\end{theorem}

\begin{proof}
By \Cref{lem:central-reduction}, shifting $k\mapsto k+m$ in the
infinite component gives
\begin{align*}
 &\sum_{k=m+1}^{\infty}
 \Qbinom{2k}{k}{q^2}
 \frac{q^k}{(-q;q)_{2k}[k-m]_{q^2}}\\
 &\quad=(1-q^2)q^m\frac{(q;q^2)_m}{(q^2;q^2)_m}
 \sum_{k=1}^{\infty}
 \frac{(q^{2m+1};q^2)_kq^k}
 {(q^{2m+2};q^2)_k(1-q^{2k})}.
\end{align*}
Introduce
\[
 U(a,q)=\frac{1}{1-a}\left(
 {}_2\phi_1\!\left[
 \begin{matrix}a,q^{2m+1}\\q^{2m+2}\end{matrix};q^2,a^{-1}q
 \right]-1\right).
\]
For $a$ near $1$, absolute convergence permits termwise passage to the
limit.  Since
\[
 \frac{(a;q^2)_k}{1-a}=(aq^2;q^2)_{k-1}\qquad(k\ge1),
\]
we obtain
\[
 \lim_{a\to1}U(a,q)=
 \sum_{k=1}^{\infty}
 \frac{(q^{2m+1};q^2)_kq^k}
 {(q^{2m+2};q^2)_k(1-q^{2k})}.
\]
The $q$-Gauss summation, at base $q^2$, evaluates the auxiliary series as
\[
 U(a,q)=\frac{1}{1-a}\left(
 \frac{(q,a^{-1}q^{2m+2};q^2)_\infty}
 {(a^{-1}q,q^{2m+2};q^2)_\infty}-1\right).
\]
The product quotient tends to $1$ as $a\to1$.  L'Hospital's rule and
termwise logarithmic differentiation, justified locally uniformly, give
\begin{align*}
 \lim_{a\to1}U(a,q)
 &=\sum_{r=1}^{\infty}\frac{q^{2r-1}}{1-q^{2r-1}}
   -\sum_{r=1}^{\infty}\frac{q^{2r+2m}}{1-q^{2r+2m}}.
\end{align*}
Multiplication by the prefactor above, followed by
$(1-q)/(1-q^j)=1/[j]_q$ and \Cref{lem:central-reduction} at $k=m$,
proves \eqref{eq:central-infinite-component}.

To obtain the finite component, subtract
\eqref{eq:central-infinite-component} from the full deleted sum
\eqref{eq:singular-central}.  Indeed,
\begin{align*}
 \sum_{j=2m+1}^{\infty}\frac{(-1)^{j-1}q^j}{[j]_q}
 &=\sum_{r=m+1}^{\infty}\frac{q^{2r-1}}{[2r-1]_q}
   -\sum_{r=m+1}^{\infty}\frac{q^{2r}}{[2r]_q}\\
 &=\sum_{r=1}^{\infty}\frac{q^{2r-1}}{[2r-1]_q}
   -\sum_{r=1}^{\infty}\frac{q^{2r+2m}}{[2r+2m]_q}
   -\sum_{r=1}^{m}\frac{q^{2r-1}}{[2r-1]_q}.
\end{align*}
The first two sums constitute the infinite component, leaving exactly
\eqref{eq:central-finite-component}.
\end{proof}

\section{The classical finite part}

Define the alternating harmonic number
\[
 \overline H_N=\sum_{r=1}^{N}\frac{(-1)^{r-1}}r.
\]

\begin{theorem}[Classical deleted-singularity identity]\label{thm:classical-singular-central}
For every $m\ge0$,
\[
 \sum_{\substack{k\ge0\\k\ne m}}
 \binom{2k}{k}\frac1{4^k(k-m)}
 =\binom{2m}{m}
 \frac{\log2-\overline H_{2m}}{2^{2m-1}}.
 \tag{18.16}\label{eq:classical-singular-central}
\]
\end{theorem}

\begin{proof}
Put $c_k=(1/2)_k/k!=\binom{2k}{k}/4^k$.  For $a$ away from the nonpositive integers, the series
\[
 T(a)=\sum_{k=0}^{\infty}\frac{c_k}{k+a}
\]
converges locally uniformly and defines a meromorphic function.  For $\Re a>0$, termwise integration of
$(1-x)^{-1/2}=\sum_{k\ge0}c_kx^k$ gives
\[
 T(a)=\int_0^1x^{a-1}(1-x)^{-1/2}\,\dd x
 =B(a,1/2)
 =\frac{\Gamma(a)\Gamma(1/2)}{\Gamma(a+1/2)}.
 \tag{18.17}\label{eq:T-beta}
\]
Analytic continuation makes \eqref{eq:T-beta} valid wherever both sides are defined.  The desired deleted sum is the finite part at $a=-m$:
\[
 \lim_{a\to-m}\left(T(a)-\frac{c_m}{a+m}\right).
 \tag{18.18}\label{eq:finite-part}
\]
Write $a=-m+\eps$.  The standard expansions of the gamma function at a negative integer and at a regular point are
\begin{align*}
 \Gamma(-m+\eps)
 &=\frac{(-1)^m}{m!\,\eps}
   \left(1+(H_m-\gamma)\eps+O(\eps^2)\right),\\
 \Gamma(1/2-m+\eps)
 &=\Gamma(1/2-m)
   \left(1+\psi(1/2-m)\eps+O(\eps^2)\right).
\end{align*}
Moreover,
\[
 \frac{(-1)^m\Gamma(1/2)}{m!\Gamma(1/2-m)}=c_m.
\]
Thus the finite part in \eqref{eq:finite-part} is
\[
 c_m\bigl(H_m-\gamma-\psi(1/2-m)\bigr).
 \tag{18.19}\label{eq:finite-part-digamma}
\]
Using $\psi(1/2)=-\gamma-2\log2$ and the recurrence
$\psi(z+1)=\psi(z)+1/z$, we find
\[
 \psi(1/2-m)
 =-\gamma-2\log2+2\sum_{r=1}^{m}\frac1{2r-1}.
\]
Since
\[
 \overline H_{2m}
 =\sum_{r=1}^{m}\frac1{2r-1}-\frac12H_m,
\]
the bracket in \eqref{eq:finite-part-digamma} is
$2(\log2-\overline H_{2m})$.  Hence the finite part equals
\[
 2\frac{\binom{2m}{m}}{4^m}
 (\log2-\overline H_{2m}),
\]
which is \eqref{eq:classical-singular-central}.
\end{proof}

\begin{proposition}[Limit of the Lambert tail]\label{prop:lambert-tail-limit}
For fixed $m\ge0$,
\[
 \lim_{q\to1^-}
 \sum_{r=2m+1}^{\infty}\frac{(-1)^{r-1}q^r}{[r]_q}
 =\log2-\overline H_{2m}.
 \tag{18.20}\label{eq:lambert-tail-limit}
\]
\end{proposition}

\begin{proof}
For $0<q<1$, set $a_r(q)=q^r/[r]_q$.  These numbers are positive and decrease with $r$, since
\[
 \frac{a_{r+1}(q)}{a_r(q)}=q\frac{[r]_q}{[r+1]_q}<1.
\]
Moreover,
\[
 a_r(q)=\frac{q^r}{1+q+\cdots+q^{r-1}}
 \le\frac1r,
\]
because every summand in the denominator is at least $q^{r-1}$.  The alternating-series remainder after $N$ terms therefore has absolute value at most $a_{N+1}(q)\le1/(N+1)$, uniformly in $q$.  For each fixed $r$, $a_r(q)\to1/r$ as $q\to1^-$.  First pass to the limit in a finite partial sum and then let $N\to\infty$.  The result is the tail of the alternating harmonic series,
\[
 \sum_{r=2m+1}^{\infty}\frac{(-1)^{r-1}}r
 =\log2-\overline H_{2m}.
\]
\end{proof}

\chapter{Asymptotic regimes and Bernoulli expansions}\label{chap:asymptotics}

A $q$-analogue should explain both the deformation near $q=1$ and behavior far from it.  Gaussian coefficients have especially clean asymptotic descriptions because their logarithms are finite sums of a single elementary function.

\section{A complete expansion as \texorpdfstring{$q\to1$}{q to 1}}

Let the Bernoulli numbers be defined by
\[
 \frac{x}{e^x-1}=\sum_{r=0}^{\infty}B_r\frac{x^r}{r!}.
\]

\begin{lemma}[The basic logarithmic expansion]\label{lem:basic-bernoulli-log}
As $x\to0$,
\[
 \log\frac{1-e^{-x}}x
 =-\frac x2+
 \sum_{r=1}^{\infty}
 \frac{B_{2r}}{2r(2r)!}x^{2r}.
 \tag{19.1}\label{eq:basic-bernoulli-log}
\]
The series converges for $\abs x<2\pi$.
\end{lemma}

\begin{proof}
Differentiate the left side:
\[
 \frac{\dd}{\dd x}\log\frac{1-e^{-x}}x
 =\frac1{e^x-1}-\frac1x.
\]
By the Bernoulli generating function,
\[
 \frac1{e^x-1}-\frac1x
 =-\frac12+\sum_{r=1}^{\infty}
 B_{2r}\frac{x^{2r-1}}{(2r)!}.
\]
Integrate term by term and use the value $0$ at $x=0$.  The nearest nonzero singularities are at $x=\pm2\pi i$, giving radius $2\pi$.
\end{proof}

\begin{theorem}[Bernoulli expansion of a Gaussian coefficient]\label{thm:gaussian-bernoulli}
Let $0\le k\le n$, put $D=k(n-k)$, and set $q=e^{-t}$.  As $t\to0$,
\begin{align}
 \log\frac{\Qbinom nk{e^{-t}}}{\binom nk}
 =-&\frac D2t
 +\sum_{r=1}^{\infty}
 \frac{B_{2r}t^{2r}}{2r(2r)!}
 \sum_{j=1}^{k}
 \left((n-k+j)^{2r}-j^{2r}\right).
 \tag{19.2}\label{eq:gaussian-bernoulli}
\end{align}
For fixed $n,k$, the displayed series converges whenever $\abs t<2\pi/n$.
\end{theorem}

\begin{proof}
From the product formula,
\begin{align*}
 \Qbinom nk{e^{-t}}
 &=\prod_{j=1}^{k}
 \frac{1-e^{-(n-k+j)t}}{1-e^{-jt}}\\
 &=\binom nk
 \prod_{j=1}^{k}
 \frac{(1-e^{-(n-k+j)t})/((n-k+j)t)}
      {(1-e^{-jt})/(jt)}.
\end{align*}
Take logarithms and apply \Cref{lem:basic-bernoulli-log} to every factor.  The linear term is
\[
 -\frac t2\sum_{j=1}^{k}(n-k)=-\frac12k(n-k)t=-\frac D2t,
\]
and the even terms give the sum in \eqref{eq:gaussian-bernoulli}.  Every argument of the basic series has modulus below $2\pi$ when $\abs t<2\pi/n$, so the finite sum converges there.
\end{proof}

\begin{corollary}[First terms]\label{cor:gaussian-first-asymptotic}
With $D=k(n-k)$,
\begin{align}
 \Qbinom nk{e^{-t}}
 &=\binom nk
 \exp\left(
 -\frac D2t+\frac{D(n+1)}{24}t^2+O(t^4)
 \right),
 \tag{19.3}\label{eq:gaussian-first-asymptotic}\\
 &=\binom nk\left(
 1-\frac D2t+
 \left(\frac{D^2}{8}+\frac{D(n+1)}{24}\right)t^2
 +O(t^3)
 \right).
 \tag{19.4}\label{eq:gaussian-first-expanded}
\end{align}
\end{corollary}

\begin{proof}
For $r=1$ in \eqref{eq:gaussian-bernoulli},
\begin{align*}
 \sum_{j=1}^{k}\left((n-k+j)^2-j^2\right)
 &=k(n-k)^2+(n-k)k(k+1)\\
 &=k(n-k)(n+1)=D(n+1).
\end{align*}
Since $B_2/(2\cdot2!)=1/24$, the logarithmic formula follows.  Expanding the exponential gives the second formula.
\end{proof}

\begin{remark}[Reciprocity and centering]
The factor $e^{-Dt/2}$ is forced by reciprocity:
\[
 \Qbinom nk{e^{-t}}
 =e^{-Dt}\Qbinom nk{e^t}.
\]
Therefore the centered quantity $e^{Dt/2}\Qbinom nk{e^{-t}}$ is an even analytic function of $t$.  This explains, before any calculation, why all odd powers beyond the linear term vanish from the logarithm.
\end{remark}

\section{Fixed \texorpdfstring{$q$}{q} and growing parameters}

\begin{theorem}[Fixed-column limit]\label{thm:fixed-column-limit}
If $0<\abs q<1$ and $k$ is fixed, then
\[
 \lim_{n\to\infty}\qbinom nk=\frac1{(q;q)_k}.
 \tag{19.5}\label{eq:fixed-column-limit}
\]
More precisely,
\[
 \qbinom nk
 =\frac1{(q;q)_k}\left(1+O(\abs q^{n-k+1})\right).
 \tag{19.6}\label{eq:fixed-column-error}
\]
\end{theorem}

\begin{proof}
The product formula gives
\[
 \qbinom nk=\frac{(q^{n-k+1};q)_k}{(q;q)_k}.
\]
The numerator tends to $1$.  Since the number of factors is fixed,
\[
 \prod_{j=0}^{k-1}(1-q^{n-k+1+j})
 =1+O(\abs q^{n-k+1}),
\]
which proves both assertions.
\end{proof}

\begin{theorem}[Central fixed-$q$ limit]\label{thm:central-fixed-q}
For $0<\abs q<1$,
\[
 \lim_{m\to\infty}\qbinom{2m}{m}
 =\frac1{(q;q)_\infty}.
 \tag{19.7}\label{eq:central-fixed-q}
\]
\end{theorem}

\begin{proof}
By the factorial form,
\[
 \qbinom{2m}{m}
 =\frac{(q;q)_{2m}}{(q;q)_m^2}.
\]
Every finite shifted factorial tends to $(q;q)_\infty$ as its index tends to infinity.  The quotient therefore tends to
$(q;q)_\infty/(q;q)_\infty^2$.
\end{proof}

\begin{proposition}[Uniform elementary bound]\label{prop:gaussian-bound}
For $0<q<1$ and $0\le k\le n$,
\[
 1\le\qbinom nk\le\frac1{(q;q)_\infty}.
 \tag{19.8}\label{eq:gaussian-bound}
\]
\end{proposition}

\begin{proof}
The lower bound follows from the partition interpretation: the constant term is $1$ and every coefficient is nonnegative.  The same interpretation says that $\qbinom nk$ is the generating function of partitions fitting in a $k$ by $(n-k)$ rectangle.  Dropping the rectangle restriction embeds this set into all partitions, whose generating function is $(q;q)_\infty^{-1}$.  Evaluating at a positive $q<1$ gives the upper bound.
\end{proof}

\section{The regime \texorpdfstring{$q>1$}{q greater than 1}}

\begin{corollary}[Transfer through reciprocity]\label{cor:qgreaterone}
Let $q>1$ and $D=k(n-k)$.  Then
\[
 \qbinom nk=q^D\Qbinom nk{q^{-1}}.
 \tag{19.9}\label{eq:qgreaterone-reciprocity}
\]
Consequently, for fixed $k$,
\[
 \qbinom nk
 =\frac{q^{k(n-k)}}{(q^{-1};q^{-1})_k}
 \left(1+O(q^{-n+k-1})\right),
 \tag{19.10}\label{eq:qgreaterone-fixed-k}
\]
and
\[
 \qbinom{2m}{m}
 \sim\frac{q^{m^2}}{(q^{-1};q^{-1})_\infty}.
 \tag{19.11}\label{eq:qgreaterone-central}
\]
\end{corollary}

\begin{proof}
The reciprocity identity was proved in \Cref{thm:qbinom-structure}.  Apply \Cref{thm:fixed-column-limit,thm:central-fixed-q} with $q^{-1}$ in place of $q$.
\end{proof}

\section{Extreme specializations}

\begin{proposition}\label{prop:extreme-specializations}
For $0\le k\le n$,
\[
 \qbinom nk\bigg|_{q=0}=1,
 \qquad
 \lim_{q\to1}\qbinom nk=\binom nk,
\]
and the leading term as $q\to\infty$ is $q^{k(n-k)}$ with coefficient $1$.
\end{proposition}

\begin{proof}
The partition generating polynomial has constant term $1$, proving the first assertion.  The second follows by replacing every ratio $(1-q^a)/(1-q^b)$ in the product formula by its limit $a/b$.  Palindromicity and degree $k(n-k)$ show that the leading coefficient equals the constant coefficient, namely $1$.
\end{proof}

\chapter{Computation, certification, and formal proof architecture}\label{chap:formalization}

The identities in this book are well suited to exact computation because their natural objects are polynomials, rational functions, and formal power series.  They are also well suited to proof assistants, provided algebraic identities are separated cleanly from analytic convergence.  Lau, Lee, and Ono's 2026 formalization project~\cite{LauLeeOno2026} demonstrates this architecture for $q$-Pochhammer symbols, Bailey pairs, Jacobi's triple product, and the Rogers--Ramanujan identities.

\section{Dynamic programming for Gaussian polynomials}

A polynomial is represented here by its coefficient vector.  Multiplication by $q^h$ is therefore a shift by $h$ places.

\begin{algorithm}[Pascal dynamic program]\label{alg:gaussian-dp}
Given $n\ge0$ and $0\le k\le n$:
\begin{enumerate}
\item Replace $k$ by $\min(k,n-k)$, using symmetry.
\item Initialize $G_0(q)=1$ and $G_s(q)=0$ for $1\le s\le k$.
\item For $r=1,2,\ldots,n$, update in descending order
\[
 G_s(q)\leftarrow G_s(q)+q^{r-s}G_{s-1}(q)
 \qquad
 \bigl(s=\min(r,k),\ldots,1\bigr).
\]
\item Return $G_k(q)$.
\end{enumerate}
\end{algorithm}

\begin{theorem}[Correctness and complexity]\label{thm:gaussian-dp-correct}
\Cref{alg:gaussian-dp} returns $\qbinom nk$.  If $D=k(n-k)$ after the symmetry reduction, its coefficient-level arithmetic cost is $O(nkD)$ and its storage cost is $O(kD)$.
\end{theorem}

\begin{proof}
Let $G_s^{(r)}$ denote the value after row $r$.  We prove by induction on $r$ that
\[
 G_s^{(r)}=\Qbinom rsq
 \qquad(0\le s\le\min(r,k)).
\]
The assertion is true at $r=0$.  Descending updates ensure that both $G_s$ and $G_{s-1}$ on the right are still values from row $r-1$.  Hence the update is
\[
 G_s^{(r)}
 =\Qbinom{r-1}{s}q+q^{r-s}\Qbinom{r-1}{s-1}q
 =\Qbinom rsq
\]
by the second $q$-Pascal recurrence.  At $r=n$ this gives the desired output.

After replacing $k$ by $\min(k,n-k)$, we have $k\le n/2$.  Every intermediate polynomial $\Qbinom rsq$ with $s\le k$ has degree $s(r-s)\le k(n-k)=D$.  There are at most $nk$ shift-and-add updates, each touching at most $D+1$ coefficients.  Storing $k+1$ such vectors uses $O(kD)$ coefficients.
\end{proof}

\begin{remark}[Numerical specialization]
When only a numerical value is required, one should not first expand the polynomial.  For positive $q$ near $1$, the stable formula is
\[
 \log\qbinom nk
 =\sum_{j=1}^{k}
 \left[\log\bigl(-\operatorname{expm1}((n-k+j)\log q)\bigr)
      -\log\bigl(-\operatorname{expm1}(j\log q)\bigr)\right],
\]
where \texttt{expm1} avoids cancellation in $e^x-1$.  For exact symbolic work, the Pascal algorithm or the cyclotomic factorization is preferable.
\end{remark}

\section{A factorization algorithm}

\begin{algorithm}[Cyclotomic factor list]\label{alg:cyclotomic-list}
For each $1\le d\le n$, compute
\[
 e_d=\floor{n/d}-\floor{k/d}-\floor{(n-k)/d}.
\]
Return the list of $d$ for which $e_d=1$.
\end{algorithm}

\begin{theorem}[Correctness of the factor list]\label{thm:cyclotomic-list-correct}
The output of \Cref{alg:cyclotomic-list} is exactly the set of cyclotomic factors of $\qbinom nk$, each with multiplicity one.
\end{theorem}

\begin{proof}
This is \Cref{thm:cyclotomic-pochhammer}: the exponent of $\Phi_d(q)$ is $e_d$, and the floor calculation in that theorem proves $e_d\in\{0,1\}$.
\end{proof}

\begin{remark}
The factor-list algorithm takes only $O(n)$ integer floor operations.  It is therefore far cheaper than expanding $\qbinom nk$ when the purpose is root-of-unity evaluation, divisibility testing, or congruence work.  Polynomial expansion may be postponed until the selected cyclotomic factors are actually multiplied.
\end{remark}

\section{Telescoping certificates for finite sums}

A term $T(n,k)$ is called \emph{$q$-hypergeometric} if the ratios
\[
 \frac{T(n,k+1)}{T(n,k)},
 \qquad
 \frac{T(n+1,k)}{T(n,k)}
\]
are rational functions of $q^n$ and $q^k$ (and of any fixed parameters).  Shifted-factorial terms have this property because adjacent factors cancel.

\begin{theorem}[Finite telescoping certificate]\label{thm:telescoping-certificate}
Let
\[
 S_n=\sum_{k=L(n)}^{U(n)}T(n,k).
\]
Suppose there are functions $A_0(n),\ldots,A_r(n)$ and a certificate $G(n,k)$ such that
\[
 \sum_{j=0}^{r}A_j(n)T(n+j,k)
 =G(n,k+1)-G(n,k)
 \tag{20.1}\label{eq:certificate}
\]
for all relevant $k$, after extending $T$ by zero outside its natural finite range.  If $G(n,k)$ vanishes at both ends of the resulting finite summation interval, then
\[
 \sum_{j=0}^{r}A_j(n)S_{n+j}=0.
 \tag{20.2}\label{eq:certified-recurrence}
\]
\end{theorem}

\begin{proof}
Sum \eqref{eq:certificate} over an interval containing the support of every $T(n+j,k)$.  The left side becomes
$\sum_jA_j(n)S_{n+j}$.  The right side telescopes to the terminal value of $G$ minus its initial value, which is $0$ by hypothesis.
\end{proof}

\begin{corollary}[Identity certification]\label{cor:identity-certification}
If two finite sums satisfy the same recurrence of order $r$ and agree for $r$ consecutive initial values at every nonsingular parameter choice, then they are equal wherever the recurrence determines subsequent values.
\end{corollary}

\begin{proof}
Subtract the two sequences.  The difference satisfies the homogeneous recurrence and has $r$ consecutive zero initial values.  Induction through the recurrence gives zero thereafter.
\end{proof}

\begin{remark}[Creative telescoping]
The $q$-Zeilberger method~\cite{PetkovsekWilfZeilberger1996} searches algorithmically for a rational multiple $G(n,k)=R(q^n,q^k)T(n,k)$ satisfying \eqref{eq:certificate}.  Once found, the certificate is independently checkable by rational-function arithmetic.  The search may be sophisticated, but verification is elementary: divide by $T(n,k)$ and clear denominators.
\end{remark}

\section{Denominator clearing and specialization}

\begin{lemma}[Polynomial identity principle]\label{lem:polynomial-identity-principle}
Let $R(q),S(q)\in\Q(q)$.  If $R(q)=S(q)$ at infinitely many complex values of $q$ where both sides are defined, then $R=S$ in $\Q(q)$.
\end{lemma}

\begin{proof}
Write $R-S=A/B$ with $A,B\in\Q[q]$ and $B\ne0$.  The hypothesis says that $A$ has infinitely many zeros.  A nonzero polynomial has only finitely many zeros, so $A=0$.
\end{proof}

\begin{corollary}[Safe generic specialization]\label{cor:safe-specialization}
A rational $q$-identity proved under generic nonvanishing assumptions remains valid after any specialization at which the denominator-cleared identity is defined.  Apparent excluded values may therefore be recovered by cancellation or by taking limits.
\end{corollary}

\begin{proof}
Clear all denominators to obtain a polynomial identity.  Polynomial evaluation is valid at every specialization.  If the original rational expressions have removable singularities, divide out common factors or take their uniquely determined limits.
\end{proof}

This principle is particularly important at roots of unity.  One should first rewrite a Gaussian coefficient as its polynomial, not substitute a root into the factorial quotient and accept an indeterminate $0/0$.

\section{Integer-exponent translation and automated identity discovery}\label{sec:identity-discovery}

A large class of Gaussian identities can be generated by starting from a terminating basic hypergeometric summation and replacing every shifted factorial whose parameter is an integral power of $q$ by factorial blocks.  This old method is now amenable to systematic symbolic search; the identity finder of Zhong and Zhao~\cite{ZhongZhao2026} organizes the necessary side conditions and backtracking choices.  The proof certificate, however, remains elementary substitution.

\begin{lemma}[Two translation rules]\label{lem:integer-qpochhammer-translation}
For nonnegative integers $a,n,m,k$ with $0\le k\le m$,
\begin{align}
 \frac{(q^{a+1};q)_n}{(q;q)_n}
 &=\Qbinom{a+n}{n}{q},
 \label{eq:positive-integer-translation}\\
 \frac{(q^{-m};q)_k}{(q;q)_k}
 &=(-1)^kq^{-mk+\binom k2}\Qbinom{m}{k}{q}.
 \label{eq:negative-integer-translation}
\end{align}
More generally,
\begin{equation}\label{eq:factorial-block-translation}
 (q^{a+1};q)_n=\frac{(q;q)_{a+n}}{(q;q)_a}.
\end{equation}
\end{lemma}

\begin{proof}
The factors in \eqref{eq:factorial-block-translation} are precisely
$1-q^{a+1},\ldots,1-q^{a+n}$.  Division by $(q;q)_n$ gives
\eqref{eq:positive-integer-translation}.  For the negative rule, factor
$-q^{-m+j}$ from the $j$th factor:
\begin{align*}
 (q^{-m};q)_k
 &=\prod_{j=0}^{k-1}(1-q^{-m+j})\\
 &=(-1)^kq^{-mk+\binom k2}
   \prod_{j=0}^{k-1}(1-q^{m-j})\\
 &=(-1)^kq^{-mk+\binom k2}
   \frac{(q;q)_m}{(q;q)_{m-k}}.
\end{align*}
Divide by $(q;q)_k$ to obtain \eqref{eq:negative-integer-translation}.
\end{proof}

\begin{theorem}[Certified integral-parameter translation]\label{thm:certified-integral-translation}
Consider a terminating identity in which every summand and the closed form are products and quotients of monomials in $q$ and shifted factorials of the forms
\[
 (q^{a+1};q)_n,
 \qquad
 (q^{-m};q)_k,
 \qquad
 (q;q)_r,
\]
with integer-valued parameters satisfying side conditions that keep all displayed finite indices nonnegative and all original denominators nonzero.  Replacing the factors by \eqref{eq:positive-integer-translation}--\eqref{eq:factorial-block-translation}, cancelling common factorial blocks, and collecting monomial powers of $q$ produces an algebraically equivalent identity in Gaussian coefficients and residual $q$-factorials.  If all residual factorial multiplicities cancel, the output is a pure $q$-binomial identity.
\end{theorem}

\begin{proof}
Each replacement is an equality by \Cref{lem:integer-qpochhammer-translation}.  Products of equal rational functions are equal, cancellation of a common nonzero factor preserves equality under the stated generic assumptions, and collecting monomials is ordinary exponent arithmetic.  Hence the transformed summand agrees term by term with the original summand, and the transformed closed form agrees with the original closed form.  Summing the finitely many terms preserves equality.  Finally, clearing denominators turns the result into a polynomial identity, so removable exceptional specializations are recovered by \Cref{cor:safe-specialization}.
\end{proof}

\begin{remark}[Constraint generation]
The nontrivial search problem is not verification but representation: the same factorial block can participate in several Gaussian coefficients, and the integer parameters must satisfy inequalities such as $0\le k\le m$.  An automated finder records the multiplicity of each $(q;q)_r$, searches for compatible triples
\[
 \frac{(q;q)_u}{(q;q)_v(q;q)_{u-v}}
 =\Qbinom uvq,
\]
and carries the corresponding linear inequalities as side conditions.  Once a candidate decomposition is printed, \Cref{thm:certified-integral-translation} reduces its proof to exact algebra.
\end{remark}

\begin{corollary}[Mechanically checkable certificates]\label{cor:identity-finder-certificate}
A computer-generated integral-parameter translation can be certified without trusting the search program: expand every Gaussian coefficient into factorial quotients, apply base reversal to every negative parameter, clear denominators, and compare the resulting Laurent monomials and finite products.
\end{corollary}

\begin{proof}
The prescribed steps reverse the equalities in
\Cref{lem:integer-qpochhammer-translation}.  After denominator clearing, equality is decidable by polynomial arithmetic.  Therefore the search code may be treated as a conjecture generator while the emitted certificate is checked independently.
\end{proof}

\begin{historicalnote}
Recent formal libraries reinforce the same separation of concerns.  Eberl's Isabelle development~\cite{Eberl2024QAnalogues} formalizes Gaussian polynomials, finite and infinite $q$-Pochhammer symbols, Euler expansions, the $q$-binomial theorem, and $q$-Vandermonde.  The Lean developments cited earlier and the Fabius-function repository~\cite{ReshetnikovProveIt2026} similarly distinguish denominator-free polynomial recurrences from analytic infinite products.  This architecture is exactly what makes automated discovery compatible with proof-level reliability.
\end{historicalnote}

\section{Formal series versus analytic series}

\begin{theorem}[Formal-to-analytic transfer]\label{thm:formal-analytic-transfer}
Suppose
\[
 \sum_{n\ge0}a_nz^n=\sum_{n\ge0}b_nz^n
\]
coefficientwise in $\C[[z]]$.  If both series converge for $\abs z<R$, then they define the same analytic function on that disk.  Conversely, if two power series converge in a disk and define the same function on a subset with an accumulation point in the disk, then their coefficients are equal.
\end{theorem}

\begin{proof}
The forward implication is immediate because the two partial sums agree term by term.  For the converse, the identity theorem makes the analytic functions equal throughout the disk; differentiating $n$ times at $0$ gives $n!a_n=n!b_n$.
\end{proof}

\begin{principle}[A scalable proof architecture]\label{prin:proof-architecture}
For a mixed algebraic-analytic $q$-identity, the following order minimizes hidden assumptions:
\begin{enumerate}
\item define finite shifted factorials and prove their algebraic laws;
\item prove terminating identities after denominator clearing;
\item define formal infinite products by coefficient stabilization when possible;
\item separately establish analytic convergence for $\abs q<1$;
\item transfer the formal coefficient identity to analytic functions;
\item only then take limits, differentiate parameters, or specialize at roots of unity.
\end{enumerate}
\end{principle}

\begin{proof}
Each step is justified by results already proved: finite product algebra in \Cref{chap:pochhammer}; polynomial identity and specialization by \Cref{lem:polynomial-identity-principle,cor:safe-specialization}; formal-to-analytic transfer by \Cref{thm:formal-analytic-transfer}; and the convergence theorems in \Cref{chap:infinite-qbinomial,chap:basic-hypergeometric,chap:bilateral}.  The ordering prevents an analytic operation from being used before its convergence hypothesis has been supplied.
\end{proof}

\begin{historicalnote}
Formal verification forces distinctions that handwritten mathematics often suppresses: a finite product indexed by a natural number is not the same object as a bilateral product; a polynomial Gaussian coefficient is not definitionally the same as a rational factorial quotient; and a formal power-series identity does not automatically contain a proof of analytic convergence.  The Lean development of Lau, Lee, and Ono~\cite{LauLeeOno2026} turns these distinctions into reusable interfaces rather than obstacles.
\end{historicalnote}

\part{Geometric filters, dyadic structure, and the Fabius bridge}

\chapter{Geometric interpolation and exact \texorpdfstring{$q$}{q}-filters}\label{chap:geometric-filters}
\index{geometric interpolation}
\index{Lagrange interpolation!geometric nodes}
\index{Richardson extrapolation}
\index{Gaussian coefficient!reversed row}

The geometric grid
\[
  1,q,q^2,\ldots,q^n
\]
is the natural interpolation grid for $q$-calculus.  The Newton basis on this grid was developed in \Cref{chap:qnewton}.  The complementary Lagrange viewpoint is equally useful: its weights are signed Gaussian coefficients, and every moment beyond the interpolation range is again a Gaussian coefficient.  This gives exact coefficient filters, closed error formulas for geometric Richardson extrapolation, and the algebraic mechanism behind several constructions in the Fabius-function literature.

\section{A reversed Gaussian row}

For $0\le k\le n$ define
\begin{equation}\label{eq:reversed-gaussian-row}
 N_{n,k}(q)
 :=(-1)^{n-k}q^{\binom{n-k+1}{2}}\qbinom nk,
\end{equation}
and put $N_{n,k}(q)=0$ outside this range.  Although the definition is written with a quotient notation, $N_{n,k}(q)$ is a polynomial in $q$ with integral coefficients up to the displayed monomial and sign.

\begin{theorem}[Reversed finite $q$-binomial theorem]\label{thm:reversed-qbinomial}
For every $n\ge0$,
\begin{equation}\label{eq:reversed-product}
 \sum_{k=0}^{n}N_{n,k}(q)z^k
 =\prod_{j=1}^{n}(z-q^j).
\end{equation}
Consequently, for every integer $d$ with $0\le d\le n$,
\begin{equation}\label{eq:reversed-moments}
 \sum_{k=0}^{n}N_{n,k}(q)q^{kd}
 =
 \begin{cases}
   (q;q)_n,&d=0,\\
   0,&1\le d\le n.
 \end{cases}
\end{equation}
\end{theorem}

\begin{proof}
Write the product on the right as
\[
 \prod_{j=1}^{n}(z-q^j)
 =z^n\prod_{j=1}^{n}\left(1-\frac{q^j}{z}\right)
 =z^n(q/z;q)_n.
\]
The finite $q$-binomial theorem, \Cref{thm:finite-qbinomial}, gives
\[
 z^n(q/z;q)_n
 =\sum_{r=0}^{n}(-1)^r q^{\binom r2+r}\qbinom nr z^{n-r}.
\]
Put $k=n-r$.  Since $\qbinom nr=\qbinom nk$ by symmetry and
$\binom r2+r=\binom{r+1}{2}=\binom{n-k+1}{2}$, this is exactly
\eqref{eq:reversed-product}.

Now substitute $z=q^d$.  At $d=0$ the product is
$\prod_{j=1}^{n}(1-q^j)=(q;q)_n$.  If $1\le d\le n$, the factor with
$j=d$ vanishes.  This proves \eqref{eq:reversed-moments}.
\end{proof}

\begin{remark}[Denominator-free meaning]\label{rem:denominator-free-reversed}
The identity \eqref{eq:reversed-product} belongs to $\Z[q,z]$ and therefore remains valid after mapping $q$ into any commutative ring.  No division by $1-q^j$ is needed.  This is the robust way to use Gaussian coefficients when $q$ is a root of unity or when the coefficient ring has zero divisors.
\end{remark}

\section{Lagrange weights at zero}

Assume henceforth that $q$ lies in a field $K$ and that
\[
 q^i\ne q^j\qquad(0\le i<j\le n).
\]
Equivalently, the nodes $1,q,\ldots,q^n$ are distinct.  Define
\begin{equation}\label{eq:geometric-lagrange-weight}
 \lambda_{n,k}^{(q)}
 :=\frac{N_{n,k}(q)}{(q;q)_n}
 =\frac{(-1)^{n-k}q^{\binom{n-k+1}{2}}}{(q;q)_n}\qbinom nk.
\end{equation}

\begin{theorem}[Evaluation at zero on a geometric grid]\label{thm:geometric-evaluation-zero}
For every polynomial $p\in K[X]$ of degree at most $n$,
\begin{equation}\label{eq:geometric-eval-zero}
 p(0)=\sum_{k=0}^{n}\lambda_{n,k}^{(q)}p(q^k).
\end{equation}
Thus $\lambda_{n,k}^{(q)}$ are precisely the Lagrange interpolation weights for evaluation at $0$ on the nodes $1,q,\ldots,q^n$.
\end{theorem}

\begin{proof}
It is enough to verify \eqref{eq:geometric-eval-zero} on the basis
$1,X,\ldots,X^n$.  For $p(X)=1$, \eqref{eq:reversed-moments} divided by
$(q;q)_n$ gives $1$.  For $p(X)=X^d$ with $1\le d\le n$, the same identity gives $0=p(0)$.  Linearity proves the formula for every polynomial of degree at most $n$.

The Lagrange interpolation functional at $0$ is the unique linear combination of the point evaluations $p\mapsto p(q^k)$ that agrees with $p\mapsto p(0)$ on all polynomials of degree at most $n$: uniqueness follows from invertibility of the Vandermonde matrix of the distinct nodes.  Hence the displayed coefficients are the Lagrange weights.
\end{proof}

\begin{corollary}[Explicit product form]\label{cor:geometric-lagrange-product}
For $0\le k\le n$,
\begin{equation}\label{eq:geometric-lagrange-product}
 \lambda_{n,k}^{(q)}
 =\prod_{\substack{0\le j\le n\\j\ne k}}
   \frac{-q^j}{q^k-q^j}.
\end{equation}
\end{corollary}

\begin{proof}
The right side is the usual Lagrange basis polynomial associated with the node $q^k$, evaluated at $0$.  By uniqueness in \Cref{thm:geometric-evaluation-zero}, it equals \eqref{eq:geometric-lagrange-weight}.  One can also simplify the product directly by splitting the factors with $j<k$ and $j>k$.
\end{proof}

\begin{example}
For $n=2$ the weights are
\[
 \lambda_{2,0}^{(q)}=\frac{q^3}{(1-q)(1-q^2)},\qquad
 \lambda_{2,1}^{(q)}=-\frac{q(1+q)}{(1-q)(1-q^2)},\qquad
 \lambda_{2,2}^{(q)}=\frac1{(1-q)(1-q^2)}.
\]
They satisfy
\[
 \sum_{k=0}^{2}\lambda_{2,k}^{(q)}=1,
 \qquad
 \sum_{k=0}^{2}\lambda_{2,k}^{(q)}q^k=0,
 \qquad
 \sum_{k=0}^{2}\lambda_{2,k}^{(q)}q^{2k}=0.
\]
\end{example}

\section{All residual moments}

The vanishing moments in \eqref{eq:reversed-moments} are only the beginning.  Every later moment has a closed form.  The key fact is a universal remainder formula involving complete homogeneous symmetric functions.

\begin{theorem}[Universal Lagrange residual formula]\label{thm:universal-lagrange-residual}
Let $v_1,\ldots,v_N$ be distinct elements of a field, let
\[
 \ell_i(x)=\prod_{\substack{1\le j\le N\\j\ne i}}
 \frac{x-v_j}{v_i-v_j},
\]
and let $h_r$ denote the complete homogeneous symmetric polynomial, with $h_0=1$.  For every $r\ge0$,
\begin{equation}\label{eq:universal-residual}
 \sum_{i=1}^{N}\ell_i(x)v_i^{N+r}
 =x^{N+r}-\prod_{i=1}^{N}(x-v_i)\,
 h_r(x,v_1,\ldots,v_N).
\end{equation}
For exponents $0\le d<N$, the left side with $N+r$ replaced by $d$ is simply $x^d$.
\end{theorem}

\begin{proof}
The assertion for $d<N$ is ordinary Lagrange interpolation.  For $r\ge0$, put
\[
 P(T)=\prod_{i=1}^{N}(T-v_i),
 \qquad
 H_r(T)=T^{N+r}-P(T)h_r(T,v_1,\ldots,v_N).
\]
We first show that $\deg H_r<N$.  Write $e_j=e_j(v_1,\ldots,v_N)$ and
$h_j=h_j(v_1,\ldots,v_N)$.  Then
\[
 P(T)=\sum_{j=0}^{N}(-1)^je_jT^{N-j},
 \qquad
 h_r(T,v_1,\ldots,v_N)=\sum_{s=0}^{r}T^s h_{r-s}.
\]
For $0\le \ell\le r$, the coefficient of $T^{N+r-\ell}$ in their product is
\[
 \sum_{j=0}^{\ell}(-1)^je_jh_{\ell-j}.
\]
By elementary--complete orthogonality, \Cref{cor:e-h-orthogonality}, this is $1$ when $\ell=0$ and $0$ for $1\le\ell\le r$.  Hence the terms of degrees $N,N+1,\ldots,N+r$ in $P(T)h_r(T,v_1,\ldots,v_N)$ agree with those of $T^{N+r}$, and therefore $\deg H_r<N$.

At every node $v_i$ we have $P(v_i)=0$, so
$H_r(v_i)=v_i^{N+r}$.  Since $H_r$ has degree less than $N$, Lagrange interpolation gives
\[
 H_r(x)=\sum_{i=1}^{N}\ell_i(x)v_i^{N+r}.
\]
Substituting the definition of $H_r$ proves \eqref{eq:universal-residual}.
\end{proof}

\begin{theorem}[Complete geometric moment formula]\label{thm:complete-geometric-moments}
Under the distinct-node hypothesis, for every integer $d\ge0$,
\begin{equation}\label{eq:complete-geometric-moments}
 \sum_{k=0}^{n}\lambda_{n,k}^{(q)}q^{kd}
 =
 \begin{cases}
  1,&d=0,\\
  0,&1\le d\le n,\\
  (-1)^nq^{\binom{n+1}{2}}\Qbinom{d-1}{n}{q},&d\ge n+1.
 \end{cases}
\end{equation}
Equivalently, with the convention that a Gaussian coefficient is zero outside its natural range,
\begin{equation}\label{eq:complete-geometric-moments-compact}
 \sum_{k=0}^{n}\lambda_{n,k}^{(q)}q^{kd}
 =(-1)^nq^{\binom{n+1}{2}}\Qbinom{d-1}{n}{q}
 \qquad(d>0).
\end{equation}
\end{theorem}

\begin{proof}
The cases $0\le d\le n$ were proved in \Cref{thm:geometric-evaluation-zero}.  Let $d=n+1+r$ with $r\ge0$.  Apply \Cref{thm:universal-lagrange-residual} with target $x=0$ and nodes
$v_k=q^k$ for $0\le k\le n$.  There are $N=n+1$ nodes, and the weights $\ell_k(0)$ are $\lambda_{n,k}^{(q)}$.  Since
\[
 \prod_{k=0}^{n}(0-q^k)=(-1)^{n+1}q^{\binom{n+1}{2}},
\]
we obtain
\[
 \sum_{k=0}^{n}\lambda_{n,k}^{(q)}q^{k(n+1+r)}
 =(-1)^nq^{\binom{n+1}{2}}h_r(1,q,\ldots,q^n).
\]
The geometric specialization in \Cref{thm:weighted-specializations} gives
\[
 h_r(1,q,\ldots,q^n)=\Qbinom{n+r}{n}{q}.
\]
Because $n+r=d-1$, this is the desired formula.
\end{proof}

\begin{corollary}[Scaled geometric nodes]\label{cor:scaled-geometric-moments}
Let $c\ne0$ and use the same weights $\lambda_{n,k}^{(q)}$ on the nodes $cq^k$.  Then for $d>0$,
\[
 \sum_{k=0}^{n}\lambda_{n,k}^{(q)}(cq^k)^d
 =c^d(-1)^nq^{\binom{n+1}{2}}\Qbinom{d-1}{n}{q}.
\]
In particular, these weights evaluate every polynomial of degree at most $n$ at $0$ independently of the nonzero scale $c$.
\end{corollary}

\begin{proof}
Factor $c^d$ from the moment sum and apply \Cref{thm:complete-geometric-moments}.  Exactness on polynomials follows by linearity.
\end{proof}

\section{Formal coefficient filters}

For a formal power series $f(X)=\sum_{m\ge0}a_mX^m$, define
\begin{equation}\label{eq:geometric-filter-definition}
 (\mathcal G_{n,q}f)(X)
 :=\sum_{k=0}^{n}\lambda_{n,k}^{(q)}f(q^kX).
\end{equation}
The sum is finite, so this is well defined over any coefficient field in which the weights exist.

\begin{theorem}[Diagonal action and exact tail]\label{thm:geometric-filter-diagonal}
For every formal power series $f(X)=\sum_{m\ge0}a_mX^m$,
\begin{align}
 [X^m]\mathcal G_{n,q}f
 &=\left(\sum_{k=0}^{n}\lambda_{n,k}^{(q)}q^{km}\right)a_m,
 \label{eq:filter-diagonal}\\
 \mathcal G_{n,q}f(X)
 &=a_0+(-1)^nq^{\binom{n+1}{2}}
   \sum_{r\ge0}\Qbinom{n+r}{n}{q}a_{n+1+r}X^{n+1+r}.
 \label{eq:filter-exact-tail}
\end{align}
Thus the filter deletes the coefficients of degrees $1,\ldots,n$ and multiplies every later coefficient by an explicit Gaussian polynomial.
\end{theorem}

\begin{proof}
Substitution gives
\[
 \mathcal G_{n,q}f(X)
 =\sum_{k=0}^{n}\lambda_{n,k}^{(q)}
   \sum_{m\ge0}a_mq^{km}X^m.
\]
The outer sum is finite, so coefficient extraction may be interchanged with it, proving \eqref{eq:filter-diagonal}.  Apply \Cref{thm:complete-geometric-moments}: the multiplier is $1$ at $m=0$, is $0$ for $1\le m\le n$, and at $m=n+1+r$ is
$(-1)^nq^{\binom{n+1}{2}}\Qbinom{n+r}{n}{q}$.  This gives \eqref{eq:filter-exact-tail}.
\end{proof}

\begin{corollary}[Geometric Richardson extrapolation]\label{cor:geometric-richardson}
Let $f$ be analytic at $0$.  The values
\[
 f(z),\ f(qz),\ \ldots,\ f(q^nz)
\]
produce the extrapolation formula
\[
 \sum_{k=0}^{n}\lambda_{n,k}^{(q)}f(q^kz)
 =f(0)+O(z^{n+1})
 \qquad(z\to0),
\]
and the complete asymptotic error is the convergent specialization of
\eqref{eq:filter-exact-tail}.
\end{corollary}

\begin{proof}
Apply \Cref{thm:geometric-filter-diagonal} to the Taylor series of $f$.  The first nonconstant surviving term has degree $n+1$.
\end{proof}

\begin{theorem}[A quantitative analytic bound]\label{thm:geometric-filter-bound}
Assume $0<q<1$ and
$f(z)=\sum_{m\ge0}a_mz^m$ with
$\abs{a_m}\le M R^{-m}$.  If $\abs z<R$, then
\begin{equation}\label{eq:geometric-filter-bound}
 \left|\sum_{k=0}^{n}\lambda_{n,k}^{(q)}f(q^kz)-f(0)\right|
 \le
 \frac{M q^{\binom{n+1}{2}}}{(q;q)_n}
 \frac{(\abs z/R)^{n+1}}{1-\abs z/R}.
\end{equation}
\end{theorem}

\begin{proof}
For $r\ge0$,
\[
 \Qbinom{n+r}{n}{q}
 =\frac{(q^{r+1};q)_n}{(q;q)_n}
 \le\frac1{(q;q)_n},
\]
because every factor in $(q^{r+1};q)_n$ lies in $(0,1)$.  Take absolute values in \eqref{eq:filter-exact-tail}, use the coefficient bound, and sum the resulting geometric series.
\end{proof}

\begin{remark}
The exact multiplier in \eqref{eq:filter-exact-tail}, rather than merely the order estimate $O(z^{n+1})$, is the principal advantage of a geometric grid.  The complete tail is controlled by one Gaussian row.  This phenomenon is used repeatedly in geometric extrapolation, in formal power-series filters, and in the dyadic identities of the next chapter.
\end{remark}

\chapter{Half-base arithmetic and Prouhet--Thue--Morse cancellation}\label{chap:halfbase-thuemorse}
\index{$q=1/2$ specialization}
\index{Mersenne product}
\index{Prouhet problem}
\index{Thue--Morse sequence}

The specialization $q=1/2$ turns factors $1-q^j$ into normalized Mersenne numbers.  It is especially important in exact formulas for the Fabius and Rvachev functions.  At this base, the finite $q$-binomial theorem produces a compact $n+1$-point Prouhet functional, while the classical Thue--Morse product produces a $2^n$-point functional.  Both annihilate all polynomials of degree below $n$, but their weights and node geometries are very different.

\section{Mersenne normalization}

Define
\begin{equation}\label{eq:mersenne-product}
 \mathfrak m_0:=1,
 \qquad
 \mathfrak m_n:=\prod_{j=1}^{n}(2^j-1)
 \quad(n\ge1).
\end{equation}

\begin{theorem}[Half-base factorial and Gaussian formulas]\label{thm:halfbase-normalization}
For $n\ge0$ and $0\le k\le n$,
\begin{align}
 (1/2;1/2)_n
 &=2^{-\binom{n+1}{2}}\mathfrak m_n,
 \label{eq:halfbase-factorial}\\
 \Qbinom nk{1/2}
 &=2^{-k(n-k)}
   \frac{\mathfrak m_n}{\mathfrak m_k\mathfrak m_{n-k}}.
 \label{eq:halfbase-gaussian}
\end{align}
\end{theorem}

\begin{proof}
Each factor satisfies
\[
 1-2^{-j}=2^{-j}(2^j-1).
\]
Multiplication for $1\le j\le n$ gives \eqref{eq:halfbase-factorial}, because
$1+2+\cdots+n=\binom{n+1}{2}$.  Substitute this expression into
\[
 \Qbinom nk{1/2}
 =\frac{(1/2;1/2)_n}{(1/2;1/2)_k(1/2;1/2)_{n-k}}.
\]
The power of $2$ in the quotient is
\[
 -\binom{n+1}{2}+\binom{k+1}{2}+\binom{n-k+1}{2}
 =-k(n-k),
\]
which proves \eqref{eq:halfbase-gaussian}.
\end{proof}

\begin{theorem}[Dyadic evaluation of the finite product]\label{thm:dyadic-product-evaluation}
For integers $m,n\ge0$,
\begin{equation}\label{eq:dyadic-product-evaluation}
 (2^m;1/2)_n
 =
 \begin{cases}
  0,&m<n,\\[3pt]
  (-1)^n\dfrac{\mathfrak m_m}{\mathfrak m_{m-n}},&m\ge n.
 \end{cases}
\end{equation}
Equivalently,
\begin{equation}\label{eq:halfbase-sum-evaluation}
 \sum_{k=0}^{n}(-1)^k2^{-\binom k2}
 \Qbinom nk{1/2}2^{mk}
 =
 \begin{cases}
  0,&m<n,\\[3pt]
  (-1)^n\dfrac{\mathfrak m_m}{\mathfrak m_{m-n}},&m\ge n.
 \end{cases}
\end{equation}
\end{theorem}

\begin{proof}
By definition,
\[
 (2^m;1/2)_n=\prod_{j=0}^{n-1}(1-2^{m-j}).
\]
If $m<n$, the factor with $j=m$ is zero.  If $m\ge n$, every factor is negative and
\[
 \prod_{j=0}^{n-1}(1-2^{m-j})
 =(-1)^n\prod_{r=m-n+1}^{m}(2^r-1)
 =(-1)^n\frac{\mathfrak m_m}{\mathfrak m_{m-n}}.
\]
The finite $q$-binomial theorem at $q=1/2$ and $z=2^m$ turns the left side into the sum in \eqref{eq:halfbase-sum-evaluation}.
\end{proof}

\begin{corollary}[Root locus]\label{cor:halfbase-root-locus}
The polynomial
\[
 \sum_{k=0}^{n}(-1)^k2^{-\binom k2}\Qbinom nk{1/2}z^k
\]
has the simple roots $1,2,4,\ldots,2^{n-1}$ and no others.
\end{corollary}

\begin{proof}
It is $(z;1/2)_n=\prod_{j=0}^{n-1}(1-z/2^j)$, a degree-$n$ polynomial with exactly the displayed roots.
\end{proof}

\section{A compressed geometric Prouhet functional}

Define nodes and weights
\begin{equation}\label{eq:halfbase-prouhet-data}
 \xi_k:=-2^k,
 \qquad
 w_{n,k}:=(-1)^k2^{-\binom k2}\Qbinom nk{1/2}
 \qquad(0\le k\le n).
\end{equation}

\begin{theorem}[Geometric Prouhet identity]\label{thm:geometric-prouhet}
For every polynomial $p$ of degree at most $n$,
\begin{equation}\label{eq:geometric-prouhet}
 \sum_{k=0}^{n}w_{n,k}p(\xi_k)
 =\mathfrak m_n[X^n]p(X).
\end{equation}
In particular, the functional on the left annihilates every polynomial of degree less than $n$.
\end{theorem}

\begin{proof}
It suffices to take $p(X)=X^d$ for $0\le d\le n$.  With $q=1/2$,
\begin{align*}
 \sum_{k=0}^{n}w_{n,k}\xi_k^d
 &=(-1)^d\sum_{k=0}^{n}(-1)^kq^{\binom k2}
   \Qbinom nkq q^{-kd}\\
 &=(-1)^d(q^{-d};q)_n
\end{align*}
by the finite $q$-binomial theorem.  If $d<n$, the factor with index $j=d$ in
$(q^{-d};q)_n$ is zero.  At $d=n$, base reversal gives
\[
 (q^{-n};q)_n=(-1)^nq^{-\binom{n+1}{2}}(q;q)_n.
\]
Therefore
\[
 (-1)^n(q^{-n};q)_n
 =2^{\binom{n+1}{2}}(1/2;1/2)_n
 =\mathfrak m_n
\]
by \eqref{eq:halfbase-factorial}.  Linearity now proves \eqref{eq:geometric-prouhet}.
\end{proof}

\begin{corollary}[Affine form]\label{cor:geometric-prouhet-affine}
For constants $a,b$ and every polynomial $p$ of degree at most $n$,
\[
 \sum_{k=0}^{n}w_{n,k}p(a+b\xi_k)
 =b^n\mathfrak m_n[X^n]p(X).
\]
\end{corollary}

\begin{proof}
Apply \Cref{thm:geometric-prouhet} to the polynomial $X\mapsto p(a+bX)$.  Its leading coefficient is $b^n[X^n]p(X)$.
\end{proof}

\begin{remark}
The identity uses only $n+1$ geometrically spaced nodes.  The price for this compression is that the weights are nonuniform rational numbers involving the half-base Gaussian row.  The Thue--Morse construction below uses $2^n$ equally weighted signs on an arithmetic grid instead.
\end{remark}

\section{The Thue--Morse product}

For a nonnegative integer $r$, let $s_2(r)$ be the sum of the binary digits of $r$ and put
\[
 \tau_r:=(-1)^{s_2(r)}.
\]
This is the signed Thue--Morse sequence.

\begin{theorem}[Finite Thue--Morse product]\label{thm:thue-morse-product}
For every $n\ge0$,
\begin{equation}\label{eq:thue-morse-product}
 \sum_{r=0}^{2^n-1}\tau_r z^r
 =\prod_{j=0}^{n-1}(1-z^{2^j}).
\end{equation}
\end{theorem}

\begin{proof}
Every $r$ with $0\le r<2^n$ has a unique binary expansion
$r=\sum_{j=0}^{n-1}\varepsilon_j2^j$ with $\varepsilon_j\in\{0,1\}$.  Expanding the product on the right, choose either $1$ or $-z^{2^j}$ from the $j$th factor.  The choice vector $(\varepsilon_0,\ldots,\varepsilon_{n-1})$ contributes
\[
 (-1)^{\varepsilon_0+\cdots+\varepsilon_{n-1}}
 z^{\varepsilon_0+2\varepsilon_1+\cdots+2^{n-1}\varepsilon_{n-1}}
 =\tau_rz^r.
\]
Summing over all binary choices gives the identity.
\end{proof}

\begin{theorem}[Prouhet cancellation and first surviving moment]\label{thm:thue-morse-moments}
For $n\ge1$,
\begin{equation}\label{eq:thue-morse-moments}
 \sum_{r=0}^{2^n-1}\tau_r r^d
 =
 \begin{cases}
  0,&0\le d<n,\\
  (-1)^n2^{\binom n2}n!,&d=n.
 \end{cases}
\end{equation}
More generally, for every polynomial $p$ of degree at most $n$,
\begin{equation}\label{eq:thue-morse-polynomial-functional}
 \sum_{r=0}^{2^n-1}\tau_r p(a+br)
 =(-1)^n2^{\binom n2}n!\,b^n[X^n]p(X).
\end{equation}
\end{theorem}

\begin{proof}
Let
$P_n(z)=\sum_{r<2^n}\tau_rz^r$.  By \Cref{thm:thue-morse-product}, $P_n$ has a zero of order exactly $n$ at $z=1$, because every factor $1-z^{2^j}$ has a simple zero there.  The Euler operator $D=z\frac{d}{dz}$ satisfies
\[
 D^dP_n(1)=\sum_{r<2^n}\tau_rr^d.
\]
For $d<n$, this is zero because $D^d$ is a linear combination of ordinary derivatives of orders at most $d$, all of which vanish at $1$.

For $d=n$, the coefficient of the ordinary $n$th derivative in $D^n$ is $z^n$, while all lower derivatives still vanish at $1$.  Hence
\[
 D^nP_n(1)=P_n^{(n)}(1).
\]
The leading coefficient of the local expansion at $1$ is the product of the first derivatives of the $n$ factors:
\[
 P_n^{(n)}(1)
 =n!\prod_{j=0}^{n-1}\left.\frac{d}{dz}(1-z^{2^j})\right|_{z=1}
 =n!\prod_{j=0}^{n-1}(-2^j)
 =(-1)^n2^{\binom n2}n!.
\]
This proves \eqref{eq:thue-morse-moments}.  Apply it termwise to the polynomial $X\mapsto p(a+bX)$; all terms below degree $n$ vanish, and its leading coefficient is $b^n[X^n]p(X)$.  This gives \eqref{eq:thue-morse-polynomial-functional}.
\end{proof}

\begin{corollary}[Equal power sums]\label{cor:thue-morse-prouhet-partition}
Partition $\{0,1,\ldots,2^n-1\}$ according to the parity of $s_2(r)$.  The two classes have equal sums of $d$th powers for every $0\le d<n$.
\end{corollary}

\begin{proof}
Move the negative-sign terms in \eqref{eq:thue-morse-moments} to the other side.
\end{proof}

\section{Two cancellation mechanisms}

\begin{principle}[Geometric versus binary Prouhet filters]\label{prin:two-prouhet-filters}
There are two complementary exact functionals of order $n$:
\begin{align*}
 \sum_{k=0}^{n}(-1)^k2^{-\binom k2}\Qbinom nk{1/2}p(-2^k)
 &=\mathfrak m_n[X^n]p,\\
 \sum_{r=0}^{2^n-1}(-1)^{s_2(r)}p(r)
 &=(-1)^n2^{\binom n2}n![X^n]p,
\end{align*}
valid for $\deg p\le n$.  The first uses a Gaussian row on a geometric grid; the second uses uniform signs generated by binary digits on an arithmetic grid.
\end{principle}

\begin{proof}
These are \Cref{thm:geometric-prouhet,thm:thue-morse-moments}.
\end{proof}

The common product mechanism is now visible.  In the geometric case the annihilating polynomial is $(q^{-d};q)_n$ evaluated at integer $d$; in the binary case it is the product $\prod_{j<n}(1-z^{2^j})$ and its order of vanishing at $z=1$.  Both encode a hierarchy of moment cancellations by forcing a product to acquire prescribed zeros.

\chapter{Geometric uniform sums, \texorpdfstring{$q$}{q}-cumulants, and the Fabius--Rvachev functions}\label{chap:geometric-uniform-fabius}
\index{Fabius function}
\index{Rvachev up function}
\index{geometric uniform sum}
\index{cumulant}
\index{Bell polynomial}
\index{sinc product}

A second bridge from $q$-products to the Fabius theory is probabilistic.  Replace discrete binary digits by independent uniform coordinates and weight them geometrically.  The resulting random series has a smooth compactly supported density, a Fourier transform given by an infinite product of sinc functions, and cumulants containing the factors $1-q^{2m}$.  At $q=1/2$ its distribution function is the Fabius function and its centered density is the Rvachev up function~\cite{Fabius1966,Rvachev1990,AriasDeReyna1982,ReshetnikovProveIt2026}.

\section{The geometric uniform random series}

Let $0<q<1$, let $U_0,U_1,\ldots$ be independent random variables uniformly distributed on $[0,1]$, and define
\begin{equation}\label{eq:geometric-uniform-series}
 X_q:=(1-q)\sum_{j=0}^{\infty}q^jU_j.
\end{equation}

\begin{theorem}[Convergence, support, symmetry, and self-similarity]\label{thm:geometric-uniform-basic}
The series \eqref{eq:geometric-uniform-series} converges absolutely for every outcome and satisfies $0\le X_q\le1$.  Its topological support is $[0,1]$, and
\begin{equation}\label{eq:geometric-uniform-mean-symmetry}
 \mathbb E X_q=\frac12,
 \qquad
 X_q\overset{d}=1-X_q.
\end{equation}
If $X_q'$ is an independent copy of $X_q$, then
\begin{equation}\label{eq:geometric-uniform-fixed-point}
 X_q\overset{d}=(1-q)U_0+qX_q'.
\end{equation}
Consequently its distribution function $F_q(x)=\mathbb P(X_q\le x)$ satisfies
\begin{equation}\label{eq:geometric-uniform-cdf-equation}
 F_q(x)=\int_0^1
 F_q\!\left(\frac{x-(1-q)u}{q}\right)\,du,
\end{equation}
where $F_q(y)=0$ for $y\le0$ and $F_q(y)=1$ for $y\ge1$.
\end{theorem}

\begin{proof}
Because $0\le U_j\le1$,
\[
 0\le (1-q)\sum_{j=0}^{N-1}q^jU_j
 \le(1-q)\sum_{j=0}^{N-1}q^j\le1.
\]
The tail is bounded by
\[
 (1-q)\sum_{j=N}^{\infty}q^j=q^N,
\]
so the series converges absolutely and takes values in $[0,1]$.

Every $x\in[0,1]$ belongs to the support: if all of the first $N$ variables lie sufficiently close to $x$, then their weighted partial sum lies close to $x(1-q^N)$, and the remaining tail is at most $q^N$; the corresponding finite cylinder has positive probability.  No point outside $[0,1]$ can belong to the support.

Linearity of expectation and the geometric sum give
\[
 \mathbb EX_q=(1-q)\sum_{j\ge0}q^j\mathbb EU_j
 =\frac{1-q}{2}\frac1{1-q}=\frac12.
\]
Since $1-U_j$ are again independent uniform variables,
\[
 1-X_q=(1-q)\sum_{j\ge0}q^j(1-U_j)
\]
has the same distribution as $X_q$.  Splitting off the first term and reindexing the tail proves \eqref{eq:geometric-uniform-fixed-point}.  Conditioning on $U_0=u$ then gives
\[
 \mathbb P((1-q)u+qX_q'\le x)
 =F_q\!\left(\frac{x-(1-q)u}{q}\right),
\]
and integration over $u\in[0,1]$ proves \eqref{eq:geometric-uniform-cdf-equation}.
\end{proof}

\section{Fourier product and cumulants}

Use the entire function
\[
 \operatorname{sinc}z:=
 \begin{cases}
  \dfrac{\sin z}{z},&z\ne0,\\
  1,&z=0.
 \end{cases}
\]

\begin{theorem}[Characteristic function]\label{thm:geometric-uniform-characteristic}
For real $t$,
\begin{equation}\label{eq:geometric-uniform-characteristic}
 \varphi_q(t):=\mathbb E\,e^{itX_q}
 =e^{it/2}\prod_{j=0}^{\infty}
 \operatorname{sinc}\!\left(\frac{(1-q)q^jt}{2}\right).
\end{equation}
The product converges uniformly on compact subsets of the complex $t$-plane and therefore defines an entire function.
\end{theorem}

\begin{proof}
For a uniform variable $U$ on $[0,1]$ and any complex $s$,
\[
 \mathbb Ee^{sU}=\int_0^1e^{su}\,du
 =e^{s/2}\frac{\sinh(s/2)}{s/2}.
\]
At $s=ita$ this becomes
$e^{ita/2}\operatorname{sinc}(at/2)$.  Independence gives the finite product for the partial sum.  The phases combine because
$\sum_{j\ge0}(1-q)q^j=1$.

For compact $t$-sets, the arguments of the sinc factors are $O(q^j)$.  Since
$\operatorname{sinc}z=1+O(z^2)$ near $0$, the series of deviations from $1$ is bounded by a constant multiple of $\sum q^{2j}$, uniformly on the compact set.  The standard uniform convergence criterion for infinite products proves locally uniform convergence to an entire function.  Passing to the limit in the finite characteristic functions gives \eqref{eq:geometric-uniform-characteristic}.
\end{proof}

Let $Y_q=X_q-1/2$.  We use Bernoulli numbers in the convention
\begin{equation}\label{eq:bernoulli-convention-fabius}
 \frac{z}{e^z-1}=\sum_{m=0}^{\infty}B_m\frac{z^m}{m!}.
\end{equation}
The cumulants $\kappa_m(Y_q)$ are defined by
\[
 \log\mathbb Ee^{tY_q}
 =\sum_{m\ge1}\kappa_m(Y_q)\frac{t^m}{m!}
\]
near $t=0$.

\begin{theorem}[$q$-cumulants of the geometric uniform sum]\label{thm:q-cumulants-uniform}
All odd cumulants of $Y_q$ vanish, and for $m\ge1$,
\begin{equation}\label{eq:q-cumulants-uniform}
 \kappa_{2m}(Y_q)
 =\frac{B_{2m}}{2m}\,
   \frac{(1-q)^{2m}}{1-q^{2m}}.
\end{equation}
Equivalently,
\begin{equation}\label{eq:q-cumulant-log-mgf}
 \log\mathbb Ee^{tY_q}
 =\sum_{m=1}^{\infty}
   \frac{B_{2m}}{2m(2m)!}
   \frac{(1-q)^{2m}}{1-q^{2m}}t^{2m}
\end{equation}
for $t$ in a neighborhood of $0$.
\end{theorem}

\begin{proof}
For $U$ uniform on $[0,1]$,
\[
 \mathbb Ee^{t(U-1/2)}
 =\frac{\sinh(t/2)}{t/2}.
\]
Differentiating the logarithm gives
\[
 \frac{d}{dt}\log\frac{\sinh(t/2)}{t/2}
 =\frac12\coth(t/2)-\frac1t.
\]
The Bernoulli generating function implies the classical Laurent expansion
\[
 \frac12\coth(t/2)-\frac1t
 =\sum_{m=1}^{\infty}B_{2m}\frac{t^{2m-1}}{(2m)!}.
\]
Integrating from $0$ to $t$ yields
\begin{equation}\label{eq:centered-uniform-logmgf}
 \log\frac{\sinh(t/2)}{t/2}
 =\sum_{m=1}^{\infty}
   \frac{B_{2m}}{2m(2m)!}t^{2m}.
\end{equation}
Now
\[
 Y_q=(1-q)\sum_{j\ge0}q^j(U_j-1/2).
\]
Cumulant generating functions add under independence.  Substitute
$t(1-q)q^j$ into \eqref{eq:centered-uniform-logmgf}, sum over $j$, and use
$\sum_{j\ge0}q^{2mj}=1/(1-q^{2m})$.  This proves
\eqref{eq:q-cumulant-log-mgf}; comparison of coefficients gives
\eqref{eq:q-cumulants-uniform}.  Only even powers occur, so all odd cumulants vanish.
\end{proof}

\begin{definition}[Complete exponential Bell polynomials]
The polynomials $\mathcal B_n(x_1,\ldots,x_n)$ are defined by
\begin{equation}\label{eq:complete-bell-definition}
 \exp\!\left(\sum_{r\ge1}x_r\frac{t^r}{r!}\right)
 =\sum_{n\ge0}\mathcal B_n(x_1,\ldots,x_n)\frac{t^n}{n!}.
\end{equation}
\end{definition}

\begin{corollary}[All moments]\label{cor:geometric-uniform-moments}
Let $\kappa_r=\kappa_r(Y_q)$, with the values in
\Cref{thm:q-cumulants-uniform}.  Then
\begin{equation}\label{eq:centered-moments-bell}
 \mathbb E(Y_q^n)=\mathcal B_n(\kappa_1,\ldots,\kappa_n).
\end{equation}
The moments can also be generated recursively from $\mu_0=1$ by
\begin{equation}\label{eq:moment-cumulant-recurrence}
 \mu_n=\sum_{r=1}^{n}\binom{n-1}{r-1}\kappa_r\mu_{n-r},
 \qquad \mu_n:=\mathbb E(Y_q^n).
\end{equation}
Finally,
\begin{equation}\label{eq:raw-moments-from-centered}
 \mathbb E(X_q^n)=\sum_{j=0}^{n}\binom nj2^{j-n}\mu_j.
\end{equation}
\end{corollary}

\begin{proof}
The moment generating function is the exponential of the cumulant generating function.  Comparing it with the defining identity
\eqref{eq:complete-bell-definition} proves \eqref{eq:centered-moments-bell}.  If
$M(t)=\sum\mu_nt^n/n!$ and $K(t)=\sum\kappa_rt^r/r!$, then
$M=e^K$ and $M'=K'M$.  Comparing the coefficient of $t^{n-1}$ gives
\eqref{eq:moment-cumulant-recurrence}.  The last formula is the binomial expansion of $(Y_q+1/2)^n$.
\end{proof}

\begin{corollary}[The half-base cumulants]\label{cor:halfbase-cumulants}
At $q=1/2$,
\begin{equation}\label{eq:halfbase-cumulants}
 \kappa_{2m}(Y_{1/2})
 =\frac{B_{2m}}{2m(2^{2m}-1)},
 \qquad
 \kappa_{2m+1}(Y_{1/2})=0.
\end{equation}
Hence all moments of $X_{1/2}$ are rational and are computable to arbitrary order by \eqref{eq:moment-cumulant-recurrence}.
\end{corollary}

\begin{proof}
In \eqref{eq:q-cumulants-uniform},
\[
 \frac{(1-1/2)^{2m}}{1-2^{-2m}}
 =\frac1{2^{2m}-1}.
\]
The Bernoulli numbers are rational, so the recurrence uses only rational operations.
\end{proof}

\section{Superpolynomial Fourier decay and smoothness}

\begin{theorem}[Log-Gaussian decay in frequency]\label{thm:geometric-uniform-fourier-decay}
Put
\[
 A=\frac{1-q}{2},
 \qquad
 b=\log(1/q)>0.
\]
For sufficiently large $\abs t$, with $L=\log(A\abs t)$,
\begin{equation}\label{eq:log-gaussian-fourier-bound}
 \abs{\varphi_q(t)}
 \le \exp\!\left(-\frac{L^2}{2b}+\frac L2\right).
\end{equation}
In particular, for every $M\ge0$,
\[
 \abs t^M\varphi_q(t)\in L^1(\R).
\]
\end{theorem}

\begin{proof}
The phase $e^{it/2}$ has absolute value $1$, and
$\abs{\operatorname{sinc}u}\le\min(1,1/\abs u)$.  Choose
\[
 N=\floor{L/b}+1.
\]
For $0\le j<N$ we have $A\abs t q^j\ge1$.  Hence
\begin{align*}
 \abs{\varphi_q(t)}
 &\le\prod_{j=0}^{N-1}\frac1{A\abs t q^j}\\
 &=\exp\!\left(-NL+\frac{bN(N-1)}2\right).
\end{align*}
Because $L\ge b(N-1)$, the exponent is at most
$-bN(N-1)/2$.  For $L$ sufficiently large,
$N\ge L/b$ and $N-1\ge L/b-1$, so
\[
 -\frac{bN(N-1)}2
 \le-\frac{L^2}{2b}+\frac L2.
\]
This proves \eqref{eq:log-gaussian-fourier-bound}.

After the change of variables $L=\log(A t)$, the right side is a Gaussian in $L$ multiplied by $e^{L/2}$.  It therefore dominates every negative power of $t$.  More explicitly, for each $M$ the exponent
$-L^2/(2b)+(M+3/2)L$ tends to $-\infty$ quadratically, which proves integrability of $\abs t^M\varphi_q(t)$.
\end{proof}

\begin{theorem}[Smooth compactly supported density]\label{thm:geometric-uniform-smooth-density}
For every $0<q<1$, the distribution of $X_q$ has a $C^\infty$ density $f_q$ supported on $[0,1]$.  For every $m\ge0$,
\begin{equation}\label{eq:geometric-uniform-density-fourier}
 f_q^{(m)}(x)
 =\frac1{2\pi}\int_{\R}(-it)^m e^{-itx}\varphi_q(t)\,dt.
\end{equation}
All derivatives of $f_q$ vanish at the endpoints $0$ and $1$.
\end{theorem}

\begin{proof}
By \Cref{thm:geometric-uniform-fourier-decay}, $t^m\varphi_q(t)$ is integrable for every $m$.  Fourier inversion for a finite measure with integrable characteristic function gives the density
\[
 f_q(x)=\frac1{2\pi}\int_{\R}e^{-itx}\varphi_q(t)\,dt.
\]
For completeness, one may obtain this directly by convolving the distribution with a centered Gaussian of variance $\varepsilon$: its density is the inverse transform of
$\varphi_q(t)e^{-\varepsilon t^2/2}$.  Dominated convergence as
$\varepsilon\downarrow0$ yields the displayed $f_q$ uniformly.  Against every compactly supported continuous test function, the Gaussian convolutions converge both to integration against $f_q$ and weakly to the original measure; hence that measure has density $f_q$.

Because $t^m\varphi_q(t)$ is integrable, differentiation under the integral is justified by dominated convergence and gives \eqref{eq:geometric-uniform-density-fourier}.  Thus $f_q\in C^\infty(\R)$.  The random variable is supported on $[0,1]$, so the density is zero on the two open complementary half-lines.  Smoothness then forces every derivative to have value $0$ at both boundary points.
\end{proof}

\section{Finite spline approximants}

Put
\begin{equation}\label{eq:finite-geometric-uniform-sum}
 X_{q,N}:=(1-q)\sum_{j=0}^{N-1}q^jU_j,
 \qquad
 a_j:=(1-q)q^j.
\end{equation}
For $x\in\R$ and an integer $r\ge0$, write
$x_+^r=(\max(x,0))^r$; when $r=0$, interpret $x_+^0$ as the indicator of $x>0$ away from the knots.

The following nonidentically scaled uniform-sum formula is a box-spline form of the generalized Irwin--Hall distribution; compare Bradley and Gupta~\cite{BradleyGupta2004}.

\begin{theorem}[Generalized uniform-sum spline]\label{thm:finite-uniform-spline}
For $N\ge1$, the density of $X_{q,N}$ is
\begin{equation}\label{eq:finite-uniform-spline}
 f_{q,N}(x)
 =\frac1{(N-1)!\prod_{j=0}^{N-1}a_j}
 \sum_{A\subseteq\{0,\ldots,N-1\}}
 (-1)^{\abs A}
 \left(x-\sum_{j\in A}a_j\right)_+^{N-1}.
\end{equation}
It is a compactly supported piecewise polynomial of degree at most $N-1$, and its Fourier transform is the corresponding finite product in
\eqref{eq:geometric-uniform-characteristic}.
\end{theorem}

\begin{proof}
The density of $a_jU_j$ is $a_j^{-1}\one_{[0,a_j]}$.  Thus $f_{q,N}$ is the convolution of these $N$ box densities.  For $s>0$, its Laplace transform is
\[
 \prod_{j=0}^{N-1}\frac{1-e^{-a_js}}{a_js}.
\]
On the other hand,
\[
 \int_0^\infty e^{-sx}\frac{(x-c)_+^{N-1}}{(N-1)!}\,dx
 =\frac{e^{-cs}}{s^N}.
\]
Taking the Laplace transform of the right side of
\eqref{eq:finite-uniform-spline} therefore gives
\[
 \frac1{s^N\prod_j a_j}
 \sum_A(-1)^{\abs A}e^{-s\sum_{j\in A}a_j}
 =\prod_{j=0}^{N-1}\frac{1-e^{-a_js}}{a_js},
\]
because the subset sum factors as $\prod_j(1-e^{-a_js})$.  Uniqueness of the Laplace transform proves the density formula.  Piecewise polynomiality is immediate from the truncated-power representation.  The Fourier product follows either by taking Fourier transforms of the box convolutions or by the finite version of the proof of \Cref{thm:geometric-uniform-characteristic}.
\end{proof}

\begin{theorem}[Dyadic Thue--Morse spline]\label{thm:dyadic-thue-morse-spline}
At $q=1/2$,
\begin{equation}\label{eq:dyadic-thue-morse-spline}
 f_{1/2,N}(x)
 =\frac{2^{\binom{N+1}{2}}}{(N-1)!}
 \sum_{r=0}^{2^N-1}(-1)^{s_2(r)}
 \left(x-\frac r{2^N}\right)_+^{N-1}.
\end{equation}
Thus, after translating by the mean $(1-2^{-N})/2$, the inverse Fourier transform of the $N$th partial sinc product is a Thue--Morse signed spline on the dyadic grid.
\end{theorem}

\begin{proof}
Here $a_j=2^{-j-1}$, so
\[
 \prod_{j=0}^{N-1}a_j=2^{-\binom{N+1}{2}}.
\]
Every subset $A$ corresponds to a binary word $(\varepsilon_0,\ldots,\varepsilon_{N-1})$ and
\[
 \sum_{j\in A}2^{-j-1}
 =\sum_{j=0}^{N-1}\varepsilon_j2^{-j-1}.
\]
Reversing the order of the bits identifies this number with $r/2^N$ for a unique
$0\le r<2^N$, while $\abs A$ and $s_2(r)$ have the same parity.  Substitute these facts into \eqref{eq:finite-uniform-spline}.
\end{proof}

\begin{corollary}[A coupling convergence bound]\label{cor:finite-uniform-convergence}
Almost surely,
\begin{equation}\label{eq:finite-uniform-tail-bound}
 0\le X_q-X_{q,N}\le q^N.
\end{equation}
Consequently,
\begin{equation}\label{eq:cdf-coupling-bound}
 F_q(x-q^N)\le F_{q,N}(x)\le F_q(x+q^N).
\end{equation}
In particular,
\begin{equation}\label{eq:cdf-uniform-rate}
 \sup_x\abs{F_{q,N}(x)-F_q(x)}
 \le \norm{f_q}_\infty q^N.
\end{equation}
\end{corollary}

\begin{proof}
The omitted tail is nonnegative and bounded by
$(1-q)\sum_{j\ge N}q^j=q^N$, proving
\eqref{eq:finite-uniform-tail-bound}.  The inclusions
\[
 \{X_q\le x-q^N\}\subseteq\{X_{q,N}\le x\}
 \subseteq\{X_q\le x+q^N\}
\]
give \eqref{eq:cdf-coupling-bound}.  Since $F_q$ has bounded density,
$\abs{F_q(x+h)-F_q(x)}\le\norm{f_q}_\infty\abs h$; apply this on both sides.
\end{proof}

\section{The Fabius and Rvachev specializations}

Set $q=1/2$, write $X=X_{1/2}$, and let
\[
 F(x):=\mathbb P(X\le x),
\]
extended by $F(x)=0$ for $x\le0$ and $F(x)=1$ for $x\ge1$.

\begin{theorem}[The probabilistic Fabius equation]\label{thm:probabilistic-fabius}
The function $F$ is $C^\infty$, satisfies
\begin{align}
 F(1-x)&=1-F(x),\label{eq:fabius-symmetry}\\
 F(x)&=\int_{2x-1}^{2x}F(u)\,du,\label{eq:fabius-integral-equation}\\
 F'(x)&=2F(2x)-2F(2x-1).\label{eq:fabius-global-derivative}
\end{align}
In particular, for $0\le x\le1/2$,
\begin{equation}\label{eq:fabius-local-derivative}
 F'(x)=2F(2x).
\end{equation}
Together with $F(0)=0$, $F(1)=1$, and symmetry, this is the standard Fabius function.
\end{theorem}

\begin{proof}
Smoothness follows from \Cref{thm:geometric-uniform-smooth-density}.  Since $X$ has no atoms and $X\overset d=1-X$,
\[
 F(1-x)=\mathbb P(1-X\le1-x)
 =\mathbb P(X\ge x)=1-F(x),
\]
which proves \eqref{eq:fabius-symmetry}.

At $q=1/2$, the fixed-point equation is
$X\overset d=(U+X')/2$.  Conditioning on $U=u$ gives
\[
 F(x)=\int_0^1F(2x-u)\,du
 =\int_{2x-1}^{2x}F(v)\,dv,
\]
proving \eqref{eq:fabius-integral-equation}.  Differentiate using the fundamental theorem of calculus to obtain
\eqref{eq:fabius-global-derivative}.  On $0\le x\le1/2$ the term $F(2x-1)$ is zero, giving \eqref{eq:fabius-local-derivative}.
\end{proof}

Let $f=F'$ and define the centered density
\begin{equation}\label{eq:up-definition}
 \operatorname{up}(y):=
 \begin{cases}
  \dfrac12 f\!\left(\dfrac{y+1}{2}\right),&-1\le y\le1,\\[6pt]
  0,&\abs y>1.
 \end{cases}
\end{equation}
This is the Rvachev up function in the probability normalization
$\int_{\R}\operatorname{up}(y)\,dy=1$.

\begin{theorem}[Infinite sinc product and the up equation]\label{thm:up-sinc-product}
With the Fourier convention
$\widehat g(t)=\int_{\R}e^{ity}g(y)\,dy$,
\begin{equation}\label{eq:up-sinc-product}
 \widehat{\operatorname{up}}(t)
 =\prod_{k=1}^{\infty}\operatorname{sinc}\!\left(\frac{t}{2^k}\right).
\end{equation}
Moreover,
\begin{equation}\label{eq:up-functional-differential}
 \operatorname{up}'(y)
 =2\operatorname{up}(2y+1)-2\operatorname{up}(2y-1)
 \qquad(y\in\R).
\end{equation}
\end{theorem}

\begin{proof}
The random variable with density $\operatorname{up}$ is
$Y=2X-1$.  Writing $V_j=2U_j-1$, we have independent uniform variables on $[-1,1]$ and
\[
 Y=\sum_{j=0}^{\infty}2^{-j-1}V_j.
\]
The characteristic function of $V_j$ is $\operatorname{sinc}t$.  Independence therefore gives
\[
 \mathbb Ee^{itY}
 =\prod_{j=0}^{\infty}\operatorname{sinc}(t/2^{j+1}),
\]
which is \eqref{eq:up-sinc-product}.

From \eqref{eq:up-definition},
$f(x)=2\operatorname{up}(2x-1)$.  Differentiate
\eqref{eq:fabius-global-derivative} to obtain
\[
 f'(x)=4f(2x)-4f(2x-1).
\]
Set $x=(y+1)/2$.  Because
$\operatorname{up}'(y)=f'((y+1)/2)/4$, substitution of
$f(t)=2\operatorname{up}(2t-1)$ gives
\[
 \operatorname{up}'(y)
 =2\operatorname{up}(2y+1)-2\operatorname{up}(2y-1).
\]
The identity remains valid outside the support because all boundary derivatives vanish.
\end{proof}

\begin{remark}[Where the $q$-objects enter]
The parameter $q$ controls four structures at once:
\begin{enumerate}
\item the geometric coordinate weights $(1-q)q^j$;
\item the factors $1-q^{2m}$ in the cumulants \eqref{eq:q-cumulants-uniform};
\item the geometric interpolation nodes and Gaussian residual moments of \Cref{chap:geometric-filters};
\item at $q=1/2$, the Mersenne products, dyadic Thue--Morse splines, and the Fabius--Rvachev functional equations.
\end{enumerate}
Thus the appearance of $q$-Pochhammer and Gaussian coefficients in exact Fabius formulas is not accidental: both arise from the same geometric scaling law.
\end{remark}

\appendix

\chapter{Formula atlas}\label{app:atlas}

This appendix is a navigation table, not a substitute for hypotheses and proofs.  Each formula points back to the theorem where its exact domain and derivation are given.

\index{$q$-Pochhammer symbol}
\index{Gaussian coefficient}
\index{$q$-binomial coefficient|see{Gaussian coefficient}}
\index{basic hypergeometric series}
\index{cyclotomic polynomial}
\index{Jacobi triple product}
\index{Rogers--Ramanujan identities}
\index{Bailey pair}
\index{Jackson integral}
\index{$q$-gamma function}
\index{Lambert series}
\index{cyclic sieving}
\index{Newton interpolation}

\footnotesize
\begin{longtable}{>{\raggedright\arraybackslash}p{0.23\textwidth}>{\raggedright\arraybackslash}p{0.55\textwidth}>{\raggedright\arraybackslash}p{0.14\textwidth}}
\toprule
Name & Formula & Location\\
\midrule
\endfirsthead
\toprule
Name & Formula & Location\\
\midrule
\endhead
Concatenation & $(a;q)_{m+n}=(a;q)_m(aq^m;q)_n$ & \Cref{prop:concat}\\[1pt]
Dissection & $(a;q)_{rn}=\prod_{j=0}^{r-1}(aq^j;q^r)_n$ & \Cref{prop:dissection}\\[1pt]
Dissection with remainder & $(a;q)_{rn+u}=\prod_{s<u}(aq^s;q^r)_{n+1}\prod_{s\ge u}(aq^s;q^r)_n$ & \Cref{cor:dissection-remainder}\\[1pt]
Base reversal & $(a;q)_n=(-a)^nq^{\binom n2}(a^{-1}q^{1-n};q)_n$ & \Cref{prop:base-reversal}\\[1pt]
Negative index & $(a;q)_{-n}=1/(aq^{-n};q)_n$ & \Cref{prop:negative-index}\\[1pt]
Infinite logarithm & $\log(a;q)_\infty=-\sum_{r\ge1}a^r/(r(1-q^r))$ & \Cref{thm:lambert-log}\\[1pt]
Gaussian product & $\qbinom nk=(q;q)_n/((q;q)_k(q;q)_{n-k})$ & \Cref{prop:qbinom-products}\\[1pt]
Pascal I & $\qbinom nk=q^k\qbinom{n-1}k+\qbinom{n-1}{k-1}$ & \Cref{thm:qpascal}\\[1pt]
Pascal II & $\qbinom nk=\qbinom{n-1}k+q^{n-k}\qbinom{n-1}{k-1}$ & \Cref{thm:qpascal}\\[1pt]
Reciprocity & $\qbinom nk=q^{k(n-k)}\Qbinom nk{q^{-1}}$ & \Cref{thm:qbinom-structure}\\[1pt]
Finite theorem & $(z;q)_n=\sum_k(-1)^kq^{\binom k2}\qbinom nkz^k$ & \Cref{thm:finite-qbinomial}\\[1pt]
Reciprocal theorem & $(z;q)_{n+1}^{-1}=\sum_{k\ge0}\qbinom{n+k}kz^k$ & \Cref{thm:reciprocal-finite}\\[1pt]
$q$-Vandermonde & $\qbinom{m+n}r=\sum_k q^{(m-k)(r-k)}\qbinom mk\qbinom n{r-k}$ & \Cref{thm:q-vandermonde}\\[1pt]
$q$-binomial inversion & $b_n=\sum_k\qbinom nka_k\iff a_n=\sum_k(-1)^{n-k}q^{\binom{n-k}2}\qbinom nkb_k$ & \Cref{thm:qbinomial-inversion}\\[1pt]
Infinite theorem & $\sum_{n\ge0}(a;q)_nz^n/(q;q)_n=(az;q)_\infty/(z;q)_\infty$ & \Cref{thm:infinite-qbinomial}\\[1pt]
Euler I & $(z;q)_\infty=\sum_{n\ge0}(-1)^nq^{\binom n2}z^n/(q;q)_n$ & \Cref{cor:euler-infinite}\\[1pt]
Euler II & $1/(z;q)_\infty=\sum_{n\ge0}z^n/(q;q)_n$, $\abs z<1$ & \Cref{cor:euler-infinite}\\[1pt]
$q$-Gauss & ${}_2\phi_1(a,b;c;q,c/(ab))=(c/a,c/b;q)_\infty/(c,c/(ab);q)_\infty$ & \Cref{thm:q-gauss}\\[1pt]
Heine & ${}_2\phi_1(a,b;c;q,z)=\dfrac{(b,az;q)_\infty}{(c,z;q)_\infty}\,{}_2\phi_1(c/b,z;az;q,b)$ & \Cref{thm:heine}\\[1pt]
$q$-Pfaff--Saalsch\"utz & Terminating balanced ${}_3\phi_2$ product evaluation & \Cref{thm:q-pfaff}\\[1pt]
$q$-Chu--Vandermonde & Two terminating ${}_2\phi_1$ evaluations & \Cref{cor:q-chu}\\[1pt]
Ramanujan ${}_1\psi_1$ & ${}_1\psi_1(a;b;q,z)=(q,b/a,az,q/(az);q)_\infty/(b,q/a,z,b/(az);q)_\infty$ & \Cref{thm:1psi1}\\[1pt]
Jacobi triple product & $\sum_{n\in\Z}(-1)^nq^{\binom n2}z^n=(q,z,q/z;q)_\infty$ & \Cref{thm:jtp}\\[1pt]
Pentagonal theorem & $(q;q)_\infty=\sum_{n\in\Z}(-1)^nq^{n(3n-1)/2}$ & \Cref{thm:pentagonal}\\[1pt]
$q$-derivative & $D_qf(x)=(f(x)-f(qx))/((1-q)x)$ & \Cref{prop:qderivative-rules}\\[1pt]
Jackson integral & $\int_0^af(x)\,\dd_qx=(1-q)a\sum_{n\ge0}q^nf(aq^n)$ & \Cref{thm:q-fundamental}\\[1pt]
$q$-gamma & $\Gamma_q(x)=(1-q)^{1-x}(q;q)_\infty/(q^x;q)_\infty$ & \Cref{thm:qgamma-basic}\\[1pt]
Cyclotomic factorization & $\qbinom nk=\prod_d\Phi_d(q)^{\floor{n/d}-\floor{k/d}-\floor{(n-k)/d}}$ & \Cref{thm:cyclotomic-pochhammer}\\[1pt]
$q$-Lucas & $\Qbinom{ad+b}{rd+s}q\equiv\binom ar\Qbinom bsq\pmod{\Phi_d(q)}$ & \Cref{thm:q-lucas}\\[1pt]
$q$-Babbage & $\Qbinom{an}{bn}q\equiv\Qbinom ab{q^{n^2}}\pmod{\Phi_n(q)^2}$ & \Cref{thm:q-babbage}\\[1pt]
Negative upper index & $\Qbinom{-m}kq=(-1)^kq^{-mk-\binom k2}\Qbinom{m+k-1}kq$ & \Cref{prop:negative-gaussian}\\[1pt]
Geometric Newton basis & $P(x)=\sum_kc_k(-1)^kq^{\binom k2}(x;q^{-1})_k$ & \Cref{thm:geometric-newton}\\[1pt]
Central reduction & $\Qbinom{2k}k{q^2}/(-q;q)_{2k}=(q;q^2)_k/(q^2;q^2)_k$ & \Cref{lem:central-reduction}\\[1pt]
Bernoulli expansion & All-orders expansion of $\log(\Qbinom nk{e^{-t}}/\binom nk)$ & \Cref{thm:gaussian-bernoulli}\\
\bottomrule
\end{longtable}
\normalsize

\chapter{Limit and specialization dictionary}\label{app:limits}

\section{Elementary limits}

For fixed nonnegative integers, the following transitions are direct consequences of the product formulas:
\begin{align*}
 \lim_{q\to1}[n]_q&=n,
 &\lim_{q\to1}(1-q)^{-n}(q^a;q)_n&=(a)_n,\\
 \lim_{q\to1}\qbinom nk&=\binom nk,
 &\lim_{n\to\infty}\qbinom nk&=\frac1{(q;q)_k}\quad(\abs q<1),\\
 \lim_{q\to1^-}\Gamma_q(x)&=\Gamma(x),
 &\lim_{q\to0}\qbinom nk&=1\quad(0\le k\le n).
\end{align*}
Each has been proved in the main text; the relevant references are \Cref{prop:classical-poch-limit,cor:qbinom-classical,thm:fixed-column-limit,thm:qgamma-classical-limit,prop:extreme-specializations}.

\section{Classical shadows of the main summations}

Under the standard scaling $a=q^A$, $b=q^B$, and $q\to1^-$, terminating basic hypergeometric sums tend to ordinary hypergeometric sums.  The mechanism is
\[
 \frac{(q^A;q)_n}{(1-q)^n}\longrightarrow(A)_n,
 \qquad
 \frac{(q;q)_n}{(1-q)^n}\longrightarrow n!.
\]
Thus:
\begin{enumerate}
\item the finite $q$-binomial theorem tends to Newton's binomial theorem;
\item $q$-Chu--Vandermonde tends to Chu--Vandermonde;
\item $q$-Pfaff--Saalsch\"utz tends to Pfaff--Saalsch\"utz;
\item the Jackson $q$-beta integral tends to Euler's beta integral;
\item \Cref{thm:regular-central-sum} tends to \Cref{cor:regular-central-classical};
\item \Cref{thm:singular-central} tends, at the level of its Lambert tail and prefactor, to \Cref{thm:classical-singular-central}.
\end{enumerate}

\section{Roots of unity}

At a root of unity, factorial quotients should be replaced by polynomial formulas before evaluation.  The basic dictionary is:
\begin{align*}
 \Phi_d(q)\mid\qbinom nk
 &\Longleftrightarrow k\bmod d>n\bmod d,\\
 \Qbinom{ad+b}{rd+s}q
 &\equiv\binom ar\Qbinom bsq\pmod{\Phi_d(q)},\\
 \qbinom nk\bigg|_{q=-1}
 &\text{ is given by \Cref{cor:qminusone}.}
\end{align*}
These statements turn apparently singular substitutions into finite enumerations and are the algebraic basis of the subset cyclic-sieving theorem.

\chapter{Proof-dependency guide}\label{app:dependencies}

The monograph can be read linearly, but several shorter routes are useful.

\section*{Route to Gaussian combinatorics}
\[
 \text{finite products}
 \longrightarrow q\text{-Pascal}
 \longrightarrow \text{finite }q\text{-binomial theorem}
 \longrightarrow \text{partitions, paths, subspaces}.
\]
Read \Cref{chap:pochhammer,chap:gaussian,chap:finite-qbinomial,chap:combinatorics,chap:finite-geometry}.

\section*{Route to product--sum identities}
\[
\begin{gathered}
 \text{infinite }q\text{-binomial theorem}
 \longrightarrow \text{Heine transformation}
 \longrightarrow q\text{-Gauss},\\
 \text{negative indices and bilateral convergence}
 \longrightarrow {}_1\psi_1,\\
 \text{Euler expansions}
 \longrightarrow \text{Jacobi triple product}
 \longrightarrow \text{Bailey pairs}.
\end{gathered}
\]
Read \Cref{chap:infinite-qbinomial,chap:basic-hypergeometric,chap:bilateral,chap:jtp,chap:bailey}.

\section*{Route to arithmetic}
\[
 1-q^n=\prod_{d\mid n}\Phi_d(q)
 \longrightarrow \text{carry criterion}
 \longrightarrow q\text{-Lucas}
 \longrightarrow \text{cyclic sieving and }q\text{-Babbage}.
\]
Read \Cref{chap:cyclotomic}.

\section*{Route to analytic deformation}
\[
\begin{gathered}
 \log(a;q)_\infty
 \longrightarrow \text{Lambert series and dilogarithms}
 \longrightarrow q\text{-gamma}\\
 \longrightarrow \text{central sums and Bernoulli asymptotics}.
\end{gathered}
\]
Read \Cref{chap:pochhammer,chap:qgamma,chap:central,chap:asymptotics}.

\backmatter

\begin{thebibliography}{99}

\bibitem{Andrews1976Partitions}
G.~E. Andrews,
\emph{The Theory of Partitions},
Addison--Wesley, 1976; Cambridge University Press reprint, 1998.

\bibitem{Andrews1986QSeries}
G.~E. Andrews,
\emph{$q$-Series: Their Development and Application in Analysis, Number Theory, Combinatorics, Physics, and Computer Algebra},
CBMS Regional Conference Series in Mathematics 66, American Mathematical Society, 1986.

\bibitem{Andrews1999Congruences}
G.~E. Andrews,
$q$-analogs of the binomial coefficient congruences of Babbage, Wolstenholme and Glaisher,
\emph{Discrete Mathematics} \textbf{204} (1999), 15--25.

\bibitem{AndrewsAskeyRoy1999}
G.~E. Andrews, R.~Askey, and R.~Roy,
\emph{Special Functions},
Encyclopedia of Mathematics and its Applications 71, Cambridge University Press, 1999.

\bibitem{Bailey1935}
W.~N. Bailey,
\emph{Generalized Hypergeometric Series},
Cambridge Tracts in Mathematics and Mathematical Physics 32, Cambridge University Press, 1935.

\bibitem{Bressoud1999}
D.~M. Bressoud,
\emph{Proofs and Confirmations: The Story of the Alternating Sign Matrix Conjecture},
MAA Spectrum, Mathematical Association of America, 1999.

\bibitem{ChernDilcherJiu2026}
S.~Chern, K.~Dilcher, and L.~Jiu,
New central $q$-binomial identities,
preprint, arXiv:2608.05498, 2026.

\bibitem{DLMF17}
NIST Digital Library of Mathematical Functions,
Chapter 17, ``$q$-Hypergeometric and Related Functions,''
F.~W.~J. Olver et al., editors, \url{https://dlmf.nist.gov/17}, accessed August 2026.

\bibitem{Fine1988}
N.~J. Fine,
\emph{Basic Hypergeometric Series and Applications},
Mathematical Surveys and Monographs 27, American Mathematical Society, 1988.

\bibitem{Gasper1995}
G.~Gasper,
Lecture notes for an introductory minicourse on $q$-series,
preprint, arXiv:math/9509223, 1995.

\bibitem{GasperRahman2004}
G.~Gasper and M.~Rahman,
\emph{Basic Hypergeometric Series}, second edition,
Encyclopedia of Mathematics and its Applications 96, Cambridge University Press, 2004.

\bibitem{GoldmanRota1970}
J.~Goldman and G.-C.~Rota,
On the foundations of combinatorial theory IV: finite vector spaces and Eulerian generating functions,
\emph{Studies in Applied Mathematics} \textbf{49} (1970), 239--258.

\bibitem{Heine1878}
E.~Heine,
\emph{Handbuch der Kugelfunctionen}, second edition,
G.~Reimer, Berlin, 1878.

\bibitem{Ismail2005}
M.~E.~H. Ismail,
\emph{Classical and Quantum Orthogonal Polynomials in One Variable},
Encyclopedia of Mathematics and its Applications 98, Cambridge University Press, 2005.

\bibitem{Jackson1908}
F.~H. Jackson,
On $q$-functions and a certain difference operator,
\emph{Transactions of the Royal Society of Edinburgh} \textbf{46} (1908), 253--281.

\bibitem{KacCheung2002}
V.~Kac and P.~Cheung,
\emph{Quantum Calculus},
Universitext, Springer, 2002.

\bibitem{KoekoekLeskySwarttouw2010}
R.~Koekoek, P.~A. Lesky, and R.~F. Swarttouw,
\emph{Hypergeometric Orthogonal Polynomials and Their $q$-Analogues},
Springer Monographs in Mathematics, Springer, 2010.

\bibitem{Konvalina1998}
J.~Konvalina,
Generalized binomial coefficients and the subset--subspace problem,
\emph{Advances in Applied Mathematics} \textbf{21} (1998), 228--240.

\bibitem{Kupershmidt2000}
B.~A. Kupershmidt,
$q$-Newton binomial: from Euler to Gauss,
\emph{Journal of Nonlinear Mathematical Physics} \textbf{7} (2000), 244--262.

\bibitem{LauLeeOno2026}
K.~Lau, S.~Lee, and K.~Ono,
Formalized $q$-series: the Rogers--Ramanujan identities and beyond,
preprint, arXiv:2607.01544, 2026.

\bibitem{Macdonald1995}
I.~G. Macdonald,
\emph{Symmetric Functions and Hall Polynomials}, second edition,
Oxford Mathematical Monographs, Oxford University Press, 1995.

\bibitem{O'Hara1990}
K.~M. O'Hara,
Unimodality of Gaussian coefficients: a constructive proof,
\emph{Journal of Combinatorial Theory, Series A} \textbf{53} (1990), 29--52.

\bibitem{PetkovsekWilfZeilberger1996}
M.~Petkov\v{s}ek, H.~S. Wilf, and D.~Zeilberger,
\emph{$A=B$},
A K Peters, Wellesley, MA, 1996.

\bibitem{ReinerStantonWhite2004}
V.~Reiner, D.~Stanton, and D.~White,
The cyclic sieving phenomenon,
\emph{Journal of Combinatorial Theory, Series A} \textbf{108} (2004), 17--50.

\bibitem{Rogers1894}
L.~J. Rogers,
Second memoir on the expansion of certain infinite products,
\emph{Proceedings of the London Mathematical Society} \textbf{25} (1894), 318--343.

\bibitem{Rota1964}
G.-C.~Rota,
On the foundations of combinatorial theory I: theory of M\"obius functions,
\emph{Zeitschrift f\"ur Wahrscheinlichkeitstheorie und Verwandte Gebiete} \textbf{2} (1964), 340--368.

\bibitem{Sagan1992}
B.~E. Sagan,
Congruence properties of $q$-analogs,
\emph{Advances in Mathematics} \textbf{95} (1992), 127--143.

\bibitem{Slater1966}
L.~J. Slater,
\emph{Generalized Hypergeometric Functions},
Cambridge University Press, 1966.

\bibitem{Stanley2012}
R.~P. Stanley,
\emph{Enumerative Combinatorics}, Volume 1, second edition,
Cambridge Studies in Advanced Mathematics 49, Cambridge University Press, 2012.

\bibitem{Stanley1999}
R.~P. Stanley,
\emph{Enumerative Combinatorics}, Volume 2,
Cambridge Studies in Advanced Mathematics 62, Cambridge University Press, 1999.

\bibitem{Straub2011Ljunggren}
A.~Straub,
A $q$-analog of Ljunggren's binomial congruence,
\emph{DMTCS Proceedings}, FPSAC 2011, 897--902; arXiv:1103.3258.

\bibitem{Zeilberger1989}
D.~Zeilberger,
Kathy O'Hara's constructive proof of the unimodality of the Gaussian polynomials,
\emph{American Mathematical Monthly} \textbf{96} (1989), 590--602.

\bibitem{Zudilin2019}
W.~Zudilin,
Congruences for $q$-binomial coefficients,
\emph{Annals of Combinatorics} \textbf{23} (2019), 1123--1135.

\bibitem{AriasDeReyna1982}
J.~Arias de Reyna,
An infinitely differentiable function with compact support: definition and properties,
\emph{Revista de la Real Academia de Ciencias Exactas, Fisicas y Naturales de Madrid} \textbf{76} (1982), 21--38;
English translation, arXiv:1702.05442.

\bibitem{BradleyGupta2004}
D.~M. Bradley and R.~C. Gupta,
On the distribution of the sum of $n$ non-identically distributed uniform random variables,
preprint, arXiv:math/0411298, 2004.

\bibitem{Eberl2024QAnalogues}
M.~Eberl,
Combinatorial $q$-analogues,
\emph{Archive of Formal Proofs}, 2024, current document revision February 2026,
\url{https://isa-afp.org/entries/Combinatorial_Q_Analogues.html}.

\bibitem{Fabius1966}
J.~Fabius,
A probabilistic example of a nowhere analytic $C^\infty$-function,
\emph{Zeitschrift f\"ur Wahrscheinlichkeitstheorie und Verwandte Gebiete} \textbf{5} (1966), 173--174.

\bibitem{ReshetnikovProveIt2026}
V.~Reshetnikov,
\emph{ProveIt: Fabius Function development and documentation},
GitHub repository, 2026,
\url{https://github.com/VladimirReshetnikov/ProveIt}.

\bibitem{Rvachev1990}
V.~A. Rvachev,
Compactly supported solutions of functional-differential equations and their applications,
\emph{Russian Mathematical Surveys} \textbf{45} (1990), no.~1, 87--120.

\bibitem{ZhongZhao2026}
H.~Zhong and L.~Zhao,
$q$-binomial identities finder,
\emph{Advances in Applied Mathematics} \textbf{172} (2026), article 102965;
arXiv:2408.00186.

\end{thebibliography}

\printindex

\end{document}
